% Standalone source. All manuscript text, references, and printed listings are embedded.
\documentclass[reqno]{amsart}
\usepackage[T1]{fontenc}
\usepackage{lmodern}
\usepackage{amsmath,amssymb,mathtools,mathrsfs}
\usepackage{booktabs,longtable,array}
\usepackage{enumitem}
\usepackage{xcolor}
\usepackage{microtype}
\usepackage[unicode,pdfusetitle,hidelinks]{hyperref}
\usepackage{listings}
\setlength{\emergencystretch}{2em}
\setcounter{tocdepth}{1}
\allowdisplaybreaks[2]
\newtheorem{theorem}{Theorem}[section]
\newtheorem{proposition}[theorem]{Proposition}
\newtheorem{lemma}[theorem]{Lemma}
\newtheorem{corollary}[theorem]{Corollary}
\theoremstyle{definition}
\newtheorem{definition}[theorem]{Definition}
\theoremstyle{remark}
\newtheorem{remark}[theorem]{Remark}
\newtheorem{example}[theorem]{Example}
\newcommand{\C}{\mathbb C}
\newcommand{\R}{\mathbb R}
\newcommand{\N}{\mathbb N}
\newcommand{\D}{\mathbb D}
\newcommand{\HH}{\mathbb H}
\DeclareMathOperator{\Tr}{Tr}
\DeclareMathOperator{\Id}{Id}
\DeclareMathOperator{\supp}{supp}
\DeclareMathOperator{\rank}{rank}
\DeclareMathOperator{\diag}{diag}
\DeclareMathOperator{\osc}{osc}
\DeclareMathOperator{\dist}{dist}


\title[Lee--Yang quantum dynamics and channels]{Lee--Yang preserving quantum Markov semigroups and channels}
\author{Dongsheng Wei}
\address{Independent Researcher}
\email{dongshengwei2025@icloud.com}
\date{}
\subjclass[2020]{Primary 81P48, 47D06; Secondary 30C15, 81P45, 22E46}
\keywords{Lee--Yang polynomial, quantum Markov semigroup, quantum channel, stability preserver, finite spin, generator classification}
\hypersetup{pdftitle={Lee-Yang preserving quantum Markov semigroups and channels},pdfauthor={Dongsheng Wei}}
\begin{document}
\begin{abstract}
We classify continuous quantum Markov semigroups that preserve the
Lee--Yang property of every density matrix in the class. On two or more
qubits, ordinary preservation forces tensor-product dynamics, with each
local generator characterized by two explicit semidefinite conditions;
it is then equivalent to preservation with arbitrary inert reference
systems. For arbitrary finite spins with at least one spin greater than
one half, preservation of all pure Lee--Yang inputs already excludes
intersite terms and forces isotropic diffusion with axial drift at every
higher-spin site. We also classify invertible static qubit maps with
stable full Choi tensors, give the complete one-qubit full-Choi
classification including singular maps,
and obtain a root parametrization for independently covariant maps
between finite spins. Finally, correlated inputs yield an exact
first-order defect and dimension-independent quantitative rigidity for
diagonal Hamiltonians, whereas product tests do not provide uniform
robustness. The proofs combine physical positivity with polynomial
root collisions and coherent tensor contraction.
\end{abstract}
\maketitle
\tableofcontents
\clearpage
\section{Introduction}\label{intro:section}
The Lee--Yang property associates to a quantum density matrix a polynomial
that is zero-free when its independent row and column variables lie in the
unit disk. It extends the use of stable polynomial tensors from equilibrium
operators to quantum channels. A channel can preserve this property for a
particular state, for every state in the class on its given system, or even
in the presence of an arbitrary inert reference system. These requirements
lead to different static problems. Our main result is that continuous
physical dynamics impose a much more rigid structure.

On two or more qubits, ordinary preservation of all Lee--Yang density
matrices forces a completely positive, trace-preserving (CPTP) semigroup
to be a tensor product of one-qubit
semigroups. Each local generator is specified by an axial Hamiltonian and
two explicit positive semidefinite matrix conditions
(Theorem~\ref{dyn:main}). No locality, unitality, or restriction on the
initial GKLS noise operators is assumed. A finite family of equatorial
product states already detects every intersite term. Further two-site
inputs prevent compensation between local defects and recover the exact
local cone. Ordinary and complete preservation consequently coincide for
these semigroups. On one qubit, the ordinary cone is larger, and we give
its separate characterization in Proposition~\ref{dyn:one-qubit}.

For arbitrary finite spins, the classification becomes stricter. If at
least one site has spin at least one, preservation of all pure Lee--Yang
inputs for sufficiently small times already forces a sum of local
generators (Theorem~\ref{dyn:arbitrary-spin}). Every site of spin at least
one carries isotropic spin diffusion and an axial Hamiltonian; qubit
sites retain the preceding semidefinite cone. The proof first reduces
arbitrary noise operators to linear spin operators by repeated-root
contacts of physical pure states. Positivity makes the coefficients of
forbidden high-order terms sums of squares, so they cannot cancel.
This is where polynomial root geometry enters the physical classification.

The static problem retains features that a flow from the identity cannot
exhibit. We classify all invertible Hermiticity-preserving (HP),
trace-preserving (TP) qubit maps with stable full
Choi tensors as local maps followed by a permutation of the physical
wires (Theorem~\ref{stat:invertible}). We then give the complete
one-qubit full-Choi classification, including singular maps, and structural results
for singular multisite channels and stable isometries. The separate
covariant problem at arbitrary input and output spins has an exact root
parametrization (Theorem~\ref{spin:covariant}). Under independent local
$SU(2)$ covariance and trace preservation, stable coherent Choi kernels
force tensor factorization and complete positivity, even when positivity
was not imposed initially.

The last section measures a limitation of product tests. For diagonal
Hamiltonians, correlated two-block inputs recover an exact first-order
trace-norm defect. This defect controls distance to the local diagonal
commutator space with a dimension-independent constant
(Theorem~\ref{rob:diagonal}). In contrast, product-input tests alone admit
families with vanishing defects while the distance to the local space stays
bounded below. The comparison constant comes from the Kalton--Roberts
approximate-additivity theorem; the polynomial contribution is the exact
root-contact identity connecting that theorem to physical distance.

\subsection*{Relation to earlier work}
Our use of disk-stable tensors and contraction follows the classical
Asano and stability-preserver theories
\cite{AsanoJapanese70,BB09,BB10II}. Wong, Bravyi, Gosset, and Liu
\cite{WBG26} formulate Lee--Yang quantum tensors and provide useful
preserving channels. Here the main question is necessity: which
continuous physical dynamics preserve the entire class, when neither
locality nor a prescribed family of jumps is given? The finite-dimensional
GKLS representation \cite{GKS76,Lindblad76} supplies the general
physical generator; the zero-free constraints supply the additional
classification.

Isotropic spin diffusion and its multipole spectrum are established
objects \cite{DenisMartin22,FilippovKuzhamuratova19}. Our higher-spin
result derives that form from Lee--Yang preservation without assuming
rotation covariance. Likewise, the representation theory of covariant
channels is inherited from \cite{AlNuwairan14,ARS24}; the root restrictions
and forced factorization in Theorem~\ref{spin:covariant} are additional
consequences of zero-freeness. The short application to rotation
convolution semigroups is stated with the Gaussian recognition results
of Voit \cite{Voit03} in view, rather than as a separate characterization
of Brownian motion.

The finite-box polynomial generator classification and its two precise
local tools are imported from the companion manuscript
\cite{WeiGenerators2026}. Section~\ref{gen:section} states them with
all degree conventions: the real and complex generator theorems,
the order-two box decomposition, and the lower-half-plane root-cluster
lemma. In particular, the necessity proof at higher spin uses the latter
two statements separately. The physical and static proofs retained here
are given in full, with the relevant imported statement cited at each use.

\subsection*{Organization}
Section~\ref{pre:section} fixes the polynomial and coherent normalizations,
and Section~\ref{gen:section} records the imported results.
Sections~\ref{dyn:section} and~\ref{dynspin:subsection} prove the
qubit and finite-spin dynamical classifications.
Sections~\ref{stat:section} and~\ref{loc:section} treat static qubit maps.
Section~\ref{spin:section} gives the independently covariant finite-spin
classification, and Section~\ref{rob:section} proves the quantitative
rigidity results.


\section{Conventions and zero-free kernels}\label{pre:section}

\subsection{Regions, degrees, and preservation}
Write $\HH=\{z\in\C:\Im z>0\}$ and $\D_r=\{z\in\C:|z|<r\}$,
with $\D=\D_1$.
\begin{definition}[Stability and preservation]\label{pre:stability}
A nonzero polynomial is \emph{$\Omega$-stable} if it
has no zero when all its variables lie in $\Omega$.
For a product of different regions the corresponding meaning is used.
A polynomial is \emph{real stable} if its coefficients are real and it is
$\HH$-stable. In one variable this is equivalent to real-rootedness.
An operator \emph{preserves} one of these classes when it sends every
member to another member or to zero.
\end{definition}

Let $\N=\{0,1,2,\ldots\}$. For $\kappa\in\N^n$ put
\[
 V_\kappa=\{f\in\C[z_1,\ldots,z_n]:\deg_{z_i}f\leq\kappa_i\},
 \qquad
 \binom{\kappa}{\alpha}=\prod_i\binom{\kappa_i}{\alpha_i}.
\]
Variables of capacity zero can be deleted. A degree bound is not an
assertion that the polynomial has that degree. All finite-box statements
include degree drops and constants. For a real operator, its action on
complex polynomials always means its complex-linear extension.
Every operator on $V_\kappa$ has a unique canonical expression
\begin{equation}\label{pre:normal-order}
 A=\sum_{\alpha\leq\kappa}P_\alpha(z)\partial^\alpha,
\end{equation}
as an operator on $V_\kappa$. Uniqueness follows by applying the expression
successively to monomials, starting with $1$: the equation for $z^\alpha$
determines $\alpha!P_\alpha$ after all lower terms are known. The phrase
\emph{differential order} refers to this expression.

A finite-dimensional generator $A$ preserves a class when $e^{tA}$
preserves it for all $t\geq0$. The polynomial and quantum operator
classifications below take place in finite-dimensional spaces.

\begin{definition}[Complete stability preservation]\label{pre:complete}
An operator is \emph{completely stable} if its tensor product with the
identity in arbitrarily many extra polynomial variables preserves the
same complex stable class, for every finite choice of the extra degree
bounds.
\end{definition}
This is a statement about polynomial
stability, distinct from complete positivity of a quantum channel.


\subsection{Polarization and the disk kernel}
The normalized polarization of $z_i^a$ into $\kappa_i$ binary slots is
\[
 z_i^a\longmapsto
 \binom{\kappa_i}{a}^{-1}
 e_a(z_{i1},\ldots,z_{i\kappa_i}).
\]
It is symmetric in the slots and specializes to $z_i^a$ on the diagonal.
Polarization preserves stability in the convex circular regions used
here, including open disks and half-planes; diagonalization preserves
it directly. We use the Grace--Walsh--Szeg\H{o} coincidence
theorem in this form; see \cite{BB09,BB10II,BB09Circular}.
The binomial normalization is part of the definition, not an optional
choice of basis.\label{pre:polarization}

For $T:V_\kappa\to V_\lambda$ define its disk kernel by
\begin{equation}\label{pre:kernel}
 K_T(z,w)=T_z\prod_i(1+z_iw_i)^{\kappa_i}
 =\sum_{\alpha\leq\kappa}\binom{\kappa}{\alpha}
       (Tz^\alpha)(z)w^\alpha.
\end{equation}
The notation $z_i$ in the first expression denotes the input variables
on which $T$ acts; its result is a polynomial in the output variables.
The kernel has output degree bound $\lambda$ and input degree bound
$\kappa$. In particular $K_{\Id}(z,w)=\prod_i(1+z_iw_i)^{\kappa_i}$.

For two polynomials of degree at most $\kappa$ in paired variables define
their bilinear contraction by
\begin{equation}\label{pre:contraction}
 [f(u),g(v)]_\kappa
   =\sum_{\alpha\leq\kappa}
      \binom{\kappa}{\alpha}^{-1} f_\alpha g_\alpha,
 \qquad
 f(u)=\sum_\alpha f_\alpha u^\alpha,\quad
 g(v)=\sum_\alpha g_\alpha v^\alpha .
\end{equation}
There is no complex conjugation in this formula. Matrix composition is
exactly contraction of the corresponding kernels:
\[
 K_{ST}(z,w)=[K_S(z,u),K_T(v,w)]_\kappa,
\]
where $\kappa$ is the intermediate capacity.

\begin{lemma}[Disk contraction]\label{pre:grace}
Suppose $f(z,u)$ and $g(v,w)$ are stable in products of disks and
$\deg_{u_i}f,\deg_{v_i}g\leq\kappa_i$. If their paired radii
$r_i,s_i$ satisfy $r_is_i\geq1$, then their contraction
\eqref{pre:contraction} is stable in the remaining disks or is identically
zero. If all $r_is_i>1$, contraction at any point in the remaining
disks is nonzero.
\end{lemma}
\begin{proof}
First consider a multiaffine polynomial
$F(u,v)=a+bu+cv+duv$ that is nonzero on $\D_r\times\D_s$.
Then $a\ne0$ and $F(rt,st)$ is nonzero for $|t|<1$.
If $d\ne0$, the product of its two roots has modulus at least one,
so $|d|rs\le|a|$; if $d=0$, this inequality is immediate.
Thus $a+d\ne0$ when $rs>1$.
The same argument applies after fixing any uncontracted variables
in their respective disks. Consequently the operation
$a+bu+cv+duv\mapsto a+d$ preserves stability under strict
separation of the paired radii. It may be iterated on a general
multiaffine polynomial, without a factorization assumption.

Polarize the paired variables of $f$ and $g$ into $\kappa_i$
slots for coordinate $i$, and apply this operation to each pair
of corresponding slots in their product.
For a fixed multi-index $\alpha$, the matching slot subsets have
cardinality $\binom{\kappa}{\alpha}$, while the two normalized
polarizations contribute the factor
$\binom{\kappa}{\alpha}^{-2}$.
The result is therefore exactly \eqref{pre:contraction}.
This proves strict nonvanishing, including degree deficiencies.
It is the weighted form of circular-domain contraction;
see \cite{AsanoJapanese70,BB09,BB10II}.
For equality of paired radii, replace every paired variable in each input
polynomial by $(1-\varepsilon)$ times that variable, apply the strict
statement, and let $\varepsilon\downarrow0$. Hurwitz's theorem gives
the stable-or-zero conclusion. The argument is valid after fixing
any values of the unpaired variables in their respective disks.
\end{proof}

The finite-box symbol theorem \cite{BB09,BB10II} identifies stability
preservers of rank at least two by the appropriate stable kernel;
the low-rank alternatives are retained whenever they arise.
In particular an invertible disk-stability preserver has a stable
kernel \eqref{pre:kernel}. 

The Cayley isomorphism is
\begin{equation}\label{pre:cayley}
 (\mathcal C_\kappa f)(x)
   =\prod_i(x_i+\mathrm i)^{\kappa_i}
      f\left(\frac{x_1-\mathrm i}{x_1+\mathrm i},\ldots,
              \frac{x_n-\mathrm i}{x_n+\mathrm i}\right).
\end{equation}
It is a linear isomorphism between the degree boxes, and $f$ is
$\D$-stable if and only if $\mathcal C_\kappa f$ is $\HH$-stable.
Thus the generator classifications in half-plane coordinates and
the physical classifications in disk coordinates concern conjugate
finite-dimensional semigroups.

\subsection{Coherent coordinates and physical tensors}
Equip $V_\kappa$ with the inner product for which
$\binom{\kappa}{\alpha}^{1/2}z^\alpha$ is an orthonormal basis.
Equivalently, identify the standard basis $|\alpha\rangle$ of
$\bigotimes_i\C^{\kappa_i+1}$ with
$\binom{\kappa}{\alpha}^{1/2}z^\alpha$. The polynomial of a vector
$\psi$ is
\[
 f_\psi(z)=\sum_{\alpha\leq\kappa}
          \binom{\kappa}{\alpha}^{1/2}\psi_\alpha z^\alpha.
\]
An operator $M$ on this Hilbert space has coherent tensor
\begin{equation}\label{pre:coherent}
 F_M(z,w)=\sum_{\alpha,\beta\leq\kappa}
 \binom{\kappa}{\alpha}^{1/2}\binom{\kappa}{\beta}^{1/2}
 M_{\alpha\beta}z^\alpha w^\beta .
\end{equation}
This is precisely \eqref{pre:kernel} for its action in coherent
coordinates. Bra variables are independent complex variables;
they are not conjugated when testing zero-freeness.
For a positive semidefinite density matrix $\rho$, the assertion
that $F_\rho$ is $\D$-stable is called its \emph{Lee--Yang property}.

For a linear map $\Phi$ between matrix algebras the full coherent
Choi tensor is defined by replacing $M_{\alpha\beta}$ in
\eqref{pre:coherent} with the coefficients
$\Phi(|\gamma\rangle\langle\delta|)_{\alpha\beta}$ and adding
input monomials $u^\gamma v^\delta$ with the corresponding square-root
binomial weights. Its stability is called \emph{full-Choi
Lee--Yang preservation}. For qubits all capacities are one, so
these weights are one. They cannot be omitted for higher spin.
Ordinary physical preservation tests only positive semidefinite
Lee--Yang inputs; it is a weaker condition for isolated maps.

We use the Pauli matrices
\[
 X=\begin{pmatrix}0&1\\1&0\end{pmatrix},\qquad
 Y=\begin{pmatrix}0&-\mathrm i\\\mathrm i&0\end{pmatrix},\qquad
 Z=\begin{pmatrix}1&0\\0&-1\end{pmatrix}.
\]
Thus $E_i=z_i\partial_i$ represents $|1\rangle\langle1|_i$
and $Z_i=I-2E_i$ on a qubit.
For general capacity $E_i=z_i\partial_i$ has eigenvalues
$0,\ldots,\kappa_i$. 


\section{Imported polynomial generator results}\label{gen:section}

We state the exact polynomial results needed in the physical proofs.
All four statements in this section are imported from \cite{WeiGenerators2026};
the physical classifications and their proofs are developed below.

Let \(\kappa=(\kappa_1,\ldots,\kappa_d)\) have finite positive integer
coordinates, and write
\[
 V_\kappa^{\mathbb F}
 =\mathbb F[z_1,\ldots,z_d]_{\deg_{z_i}\leq\kappa_i},
 \qquad \mathbb F\in\{\R,\C\}.
\]
Coordinates of capacity zero are omitted.  If no coordinate remains, the
space consists of constants and every scalar generator has the requisite
preservation property.  Stability and complete preservation have the meanings
of Definitions~\ref{pre:stability} and~\ref{pre:complete}; in this section the
domain of stability is \(\HH^d\).
A real operator acts on \(V_\kappa^\C\) by complexification.

For a polynomial \(q\) in one variable of degree at most two, define
\begin{equation}
 J_{i,\kappa_i}(q)
 =q(z_i)\partial_i-\frac{\kappa_i}{2}q'(z_i).
 \label{gen:eq-J}
\end{equation}
These are endomorphisms of the box.  The three choices \(q=1,z_i,z_i^2\)
give
\[
 J_i^-=\partial_i,\qquad
 J_i^0=z_i\partial_i-\frac{\kappa_i}{2},\qquad
 J_i^+=z_i^2\partial_i-\kappa_i z_i.
\]
Their span is the sitewise \(\mathfrak{sl}_2\) algebra.

Every endomorphism has a unique canonical normal representation
\begin{equation}
 A=\sum_{\alpha\leq\kappa}P_\alpha(z)\partial^\alpha.
 \label{gen:eq-canonical}
\end{equation}
Here the coefficient polynomials are not assumed to belong to the original
box.  The order of \(A\) means the largest \(|\alpha|\) for which its
canonical coefficient \(P_\alpha\) is nonzero.  The uniqueness of this
representation follows by triangular recursion on the monomial basis
and is important at small capacities: in particular
\(\partial_i^2\) is absent when \(\kappa_i=1\).

\begin{theorem}[Real generators; \cite{WeiGenerators2026}]
\label{gen:main-real}
Let \(A\in\operatorname{End}_{\R}(V_\kappa^\R)\).  The following conditions
are equivalent.
\begin{enumerate}
\item For every \(t\geq0\), \(e^{tA}\) preserves real stability.
\item For every \(t\geq0\), its complexification preserves complex stability,
also after adjoining arbitrary inert polynomial variables.
\item The canonical representation of \(A\) has order at most two, and its
second-order coefficients satisfy
\begin{align}
 P_{2e_i}(s)&\leq0
 &&\text{for all }s\in\R,\quad \kappa_i\geq2,
 \label{gen:eq-real-unary}\\
 P_{e_i+e_j}(s,t)&\leq0
 &&\text{for all }(s,t)\in\R^2,\quad i<j.
 \label{gen:eq-real-pair}
\end{align}
\end{enumerate}
In the third condition, box preservation forces \(P_{2e_i}\) to depend only
on \(z_i\), with degree at most four, and \(P_{e_i+e_j}\) to depend only on
\(z_i,z_j\), with bidegree at most \((2,2)\).  There are no further
inequalities on the zeroth- and first-order coefficients beyond the
identities already forced by \(A(V_\kappa^\R)\subseteq V_\kappa^\R\).
\end{theorem}

The coefficient \(P_{e_i+e_j}\) in this statement is the full coefficient
of \(\partial_i\partial_j\), with \(i<j\).  Thus a convention writing the
second-order part as \(\sum_{i,j}Q_{ij}\partial_i\partial_j\), with \(Q\)
symmetric, has \(P_{e_i+e_j}=2Q_{ij}\).
We use the full-coefficient convention throughout this section.

\begin{theorem}[Complex generators; \cite{WeiGenerators2026}]
\label{gen:main-complex}
Let \(A\in\operatorname{End}_{\C}(V_\kappa^\C)\).  The following conditions
are equivalent.
\begin{enumerate}
\item For every \(t\geq0\), \(e^{tA}\) preserves complex stability.
\item For every \(t\geq0\), \(e^{tA}\) preserves complex stability after
adjoining arbitrary inert polynomial variables.
\item There are a real generator \(R\) satisfying
Theorem~\ref{gen:main-real}, a real scalar \(c\), and real polynomials
\(q_i\) of degree at most two, nonnegative on \(\R\), such that
\begin{equation}
 A=R+i\left(c\Id+\sum_{i=1}^dJ_{i,\kappa_i}(q_i)\right).
 \label{gen:eq-complex-decomposition}
\end{equation}
\end{enumerate}
Equivalently, the canonical representation has order at most two, all its
second-order coefficients are real and satisfy
\eqref{gen:eq-real-unary}--\eqref{gen:eq-real-pair}, and every
\(\Im P_{e_i}\) is nonnegative on \(\R^d\).  In this formulation box
preservation forces
\(\Im P_{e_i}(z)=q_i(z_i)\) with \(\deg q_i\leq2\), and forces the
corresponding imaginary zeroth-order term in
\eqref{gen:eq-complex-decomposition}.
\end{theorem}

The two theorems identify ordinary and complete polynomial preservation
for semigroups starting at the identity. Their proofs are in the companion
paper \cite{WeiGenerators2026}; they are used here as imported results.

\begin{proposition}[Box decomposition; \cite{WeiGenerators2026}]
\label{gen:box-decomposition}
Every box endomorphism of canonical order at most two is a linear
combination of \(\Id\), the operators \(J_i^-,J_i^0,J_i^+\), and their
products of length two, including products at the same site.
Every endomorphism of order at most one has a unique expression
\begin{equation}
 c\Id+\sum_{i=1}^dJ_{i,\kappa_i}(q_i),
 \qquad c\in\mathbb F,\quad q_i\in\mathbb F[z_i]_{\leq2}.
 \label{gen:eq-first-order}
\end{equation}
Consequently the second-order coefficients have the variable dependence
and degree bounds stated in Theorem~\ref{gen:main-real}.  The vector space
of endomorphisms of order at most two has dimension
\begin{equation}
 1+3d+5\#\{i:\kappa_i\geq2\}+9\binom d2.
 \label{gen:eq-order-two-dimension}
\end{equation}
\end{proposition}

\begin{lemma}[Lower-half-plane root clusters; \cite{WeiGenerators2026}]
\label{gen:cluster-moments}
Let \(p_t(s)\), \(t\geq0\), be a curve of complex polynomials of uniformly
bounded degree whose coefficients are right differentiable at zero.
Suppose
\[
 p_0(s)=s^m g(s),\qquad m\geq1,\qquad g(0)\neq0,
\]
and that \(p_t\) has no zero in \(\HH\) for all sufficiently small \(t\).
Let \(C_t\) be the monic factor consisting of the \(m\) roots converging
to zero.  Its coefficients are right differentiable, and
\begin{equation}
 \left.\frac{d}{dt}C_t(s)\right|_{t=0}
 =c_1s^{m-1}+c_2s^{m-2},
 \qquad
 \Im c_1\geq0,\quad c_2\in\R,\quad c_2\leq0.
 \label{gen:eq-cluster-variation}
\end{equation}
When \(m=1\), the \(c_2\) term is absent.  If every \(p_t\) is real
rooted, then \(c_1\) is real as well.
\end{lemma}


\section{Physical quantum Markov dynamics}\label{dyn:section}

We classify continuous physical evolutions that preserve the Lee--Yang
property of every density matrix in the class.  The physical tensor
decomposition and the computational basis are fixed throughout this section.

For a matrix $M$ on $n$ qubits, put
\begin{equation}\label{dyn:polynomial}
 F_M(z,w)=\sum_{a,b\in\{0,1\}^n}M_{ab}z^a w^b.
\end{equation}
A density matrix is called \emph{Lee--Yang} if this polynomial does not vanish
on $\D^{2n}$.  Thus the row and column variables are independent.  For a pure
state, $F_{|\psi\rangle\langle\psi|}(z,w)
=f_\psi(z)\overline{f_\psi(\bar w)}$, where
$f_\psi(z)=\sum_a\psi_a z^a$.  A channel is \emph{ordinarily preserving} if it
preserves these density matrices on its given system, and is \emph{completely
preserving} if it does so after adjoining any number of inert qubits.  Complete
positivity of a channel and complete preservation in this sense are distinct
conditions.

Write $\sigma_x=X$, $\sigma_y=Y$, and $\sigma_z=Z$, with
$Z=\diag(1,-1)$.  The Pauli matrices in all generator formulas below are
unnormalized.  We use the standard finite-dimensional GKLS representation
\cite{GKS76,Lindblad76}.  If $\mathcal P_n$ denotes the nonidentity Hermitian
Pauli strings, every generator of a continuous CPTP semigroup has a
representation
\begin{equation}\label{dyn:gkls}
 \mathcal L(\rho)=-i[H,\rho]
 +\sum_{P,Q\in\mathcal P_n}K_{PQ}
 \left(P\rho Q-\tfrac12\{QP,\rho\}\right),
 \qquad H=H^*,\quad K\succeq0.
\end{equation}
Identity components in the noise have been absorbed into the Hamiltonian.

\begin{theorem}[Ordinary and complete preservation]\label{dyn:main}
Let $n\ge2$, and let $\Phi_t=e^{t\mathcal L}$ be a continuous CPTP semigroup
on $n$ qubits.  The following conditions are equivalent.
\begin{enumerate}
\item $\Phi_t$ ordinarily preserves Lee--Yang density matrices for every $t\ge0$.
\item The same preservation holds for all $t$ in some interval $[0,\varepsilon)$.
\item $\Phi_t$ completely preserves Lee--Yang density matrices for every $t\ge0$.
\item The full Choi tensor of $\Phi_t$ is Lee--Yang for every $t\ge0$.
\item The generator is a sum of one-qubit generators,
\begin{equation}\label{dyn:local-form}
 \mathcal L=\sum_{j=1}^n\mathcal L_j\otimes\Id_{\widehat j},\qquad
 \mathcal L_j(\rho)=-i[h_jZ,\rho]
 +\sum_{\alpha,\beta=x,y,z}(K_j)_{\alpha\beta}
 \left(\sigma_\alpha\rho\sigma_\beta-
 \tfrac12\{\sigma_\beta\sigma_\alpha,\rho\}\right),
\end{equation}
where $h_j\in\R$ and, for a real symmetric $2\times2$ matrix $B_j$, a vector
$d_j\in\R^2$, and $c_j\in\R$,
\begin{equation}\label{dyn:two-lmis}
 K_j=\begin{pmatrix}B_j&i d_j\\-i d_j^{\mathsf T}&c_j\end{pmatrix}\succeq0,
 \qquad
 \begin{pmatrix}
 c_jI_2-\operatorname{adj}B_j&d_j\\
 d_j^{\mathsf T}&\Tr B_j
 \end{pmatrix}\succeq0.
\end{equation}
\end{enumerate}
In this case $\Phi_t=\bigotimes_j e^{t\mathcal L_j}$.  The equivalence between
ordinary and complete preservation is false on one qubit.
\end{theorem}

We prove the theorem in four steps.  A fixed finite family of product inputs
first excludes every interaction.  Two-site pure states then fix the reality
of the remaining local coefficients.  Finally, positive definite separable
states recover the individual local inequalities even though both sites
evolve.  The resulting inequalities are reduced to
\eqref{dyn:two-lmis} without dividing by a possibly vanishing trace.

\subsection{A finite family detecting every interaction}

\begin{lemma}[A one-sided root collision]\label{dyn:collision}
Let $g_t$ be a differentiable curve in a finite-dimensional polynomial space,
with $g_0\not\equiv0$.  Suppose $0$ is a root of $g_0$ of multiplicity $m$,
and $g_t$ is upper-half-plane stable along a sequence $t\downarrow0$ of
positive times.  If $m\ge3$, then $\dot g_0(0)=0$.  If $m=2$ and
$g_0(x)=c x^2+O(x^3)$, then $\dot g_0(0)/c$ is real and nonpositive.
\end{lemma}

\begin{proof}
Choose a small circle about zero containing no other root of $g_0$.
The nearby cluster has a monic factor
$b_t(x)=x^m+a_1(t)x^{m-1}+\cdots+a_m(t)$ with $b_0=x^m$.
Its coefficients are differentiable: the contour integrals of
$x^k g_t'(x)/g_t(x)$ give the cluster power sums, and Newton identities give
its elementary symmetric functions.  This construction also applies when
the ambient degree changes.  The complementary factor is nonzero at zero.

At each of the specified stable times, write the cluster roots as
$\lambda_j=x_j+i y_j$, where $y_j\le0$.  Differentiability of the first
power sum gives $\sum_j|y_j|=O(t)$.  The second power sum is also $O(t)$,
so its real part gives $\sum_jx_j^2=O(t)$, since
$\sum_j y_j^2=O(t^2)$.  Thus every $\lambda_j=O(\sqrt t)$ and
$\sum_j\lambda_j^k=o(t)$ for $k\ge3$.  Newton identities imply
\[
 a_k(t)=o(t)\quad(k\ge3),\qquad
 a_2(t)=-\tfrac12\sum_jx_j^2+o(t).
\]
For the second equality, the imaginary part of the second power sum is
$O(t^{3/2})$ and $a_1(t)^2=O(t^2)$.  Consequently $a_2'(0)$ is real and
nonpositive, while $a_k'(0)=0$ for $k\ge3$.  Multiplication by the
complementary factor proves the assertions.  All derivatives are fixed by
the differentiable curve, so estimates along the given sequence suffice.
\end{proof}

\begin{theorem}[Finite equatorial tests]\label{dyn:trine}
Let $\Phi_t=e^{t\mathcal L}$ be a continuous CPTP semigroup on $n$ qubits.
Set $\mathcal E=\{0,\sqrt3,-\sqrt3\}$ and
\[
 \psi_a=\frac{(1+ia)|0\rangle+(1-ia)|1\rangle}
 {\sqrt{2(1+a^2)}},\qquad
 \rho_{\boldsymbol a}=\bigotimes_{j=1}^n
 |\psi_{a_j}\rangle\langle\psi_{a_j}|,\quad
 \boldsymbol a\in\mathcal E^n.
\]
Suppose that, for each of these $3^n$ densities, there is a sequence of
positive times tending to zero along which its image is Lee--Yang.
Then $\mathcal L$ is a sum of one-qubit GKLS generators.
The sequences may depend on the input.
\end{theorem}

\begin{proof}
The coefficient Cayley transform is
\begin{equation}\label{dyn:cayley}
 C=\frac1{\sqrt2}\begin{pmatrix}i&-i\\1&1\end{pmatrix},\qquad
 f(z)\longmapsto\frac{s+i}{\sqrt2}
 f\left(\frac{s-i}{s+i}\right).
\end{equation}
It takes disk stability to upper-half-plane stability.  Applied to all row
and column legs of a density tensor, it is a coefficient transformation,
not a unitary conjugation of that density matrix.  Direct multiplication
gives $C\psi_a=(-a,1)^{\mathsf T}/\sqrt{1+a^2}$ and
$C\overline{\psi_a}=(a,1)^{\mathsf T}/\sqrt{1+a^2}$.  The transformed
input is therefore a nonzero scalar multiple of
\begin{equation}\label{dyn:paired-product}
 p_{\boldsymbol a}(z,w)=\prod_j(z_j-a_j)(w_j+a_j).
\end{equation}

For a binary coefficient matrix $M=(M_{ab})$, its polynomial action has
the first-order form
\[
 M_{00}+M_{10}s+
 [M_{01}+(M_{11}-M_{00})s-M_{10}s^2]\partial_s.
\]
The identities $CXC^*=-Z$, $CYC^*=X$, $CZC^*=-Y$ consequently give
the derivative coefficients
\begin{equation}\label{dyn:pauli-vectors}
 v(s)=(2s,1-s^2,i(1+s^2))
\end{equation}
for left multiplication.  Right multiplication uses transposed Paulis;
its derivative vector at a general column coordinate $t$ is
$\widetilde v(t)=(2t,-(1-t^2),i(1+t^2))$.  In particular
$\widetilde v(-a)=-\overline{v(a)}$.

Use \eqref{dyn:gkls}.  A zero diagonal entry of a positive semidefinite
matrix has a zero row and column.  Suppose that $k\ge2$ is the greatest
weight of a Pauli string with $K_{PP}>0$, and choose such a string with
support $S$, $|S|=k$.  All noise rows of larger weight vanish.  Put both
coordinates of each site in $S$ at their roots in
\eqref{dyn:paired-product}, keeping every other coordinate away from its
root, and restrict to a line with every coordinate slope positive.
The initial polynomial has a root of multiplicity $2k\ge4$.
Lemma~\ref{dyn:collision} makes its first variation zero at this point.

Hamiltonian and anticommutator terms act only on row legs or only on column
legs, so leave selected roots untouched.  A sandwich $P\rho Q$ can survive
only if both supports contain $S$.  Maximality makes both supports exactly
$S$.  Dividing out the nonzero complementary product and the input scalar
therefore gives
\begin{equation}\label{dyn:psd-contact}
 0=(-1)^k\sum_{\supp P=\supp Q=S}
 K_{PQ}V_P(\boldsymbol a)\overline{V_Q(\boldsymbol a)},
 \qquad V_P(\boldsymbol a)=\prod_{j\in S}v_{P_j}(a_j).
\end{equation}
The sum, without its sign, is a positive semidefinite quadratic form.
Thus the exact-support block $K_{SS}$ kills
$\overline{V(\boldsymbol a)}$ for every $\boldsymbol a\in\mathcal E^S$.
The three components in \eqref{dyn:pauli-vectors} form a basis of the
quadratic polynomials.  Evaluation at the three distinct points of
$\mathcal E$ is invertible.  The resulting tensor products span
$\C^{3^k}$, forcing $K_{SS}=0$, a contradiction.  Every surviving noise
index consequently has weight one.  This still permits cross-site noise
blocks; we eliminate them below.

The anticommutator operator $\sum K_{PQ}QP$ now has support at most two.
If a Hamiltonian term has maximal support $S$ with $|S|=k\ge3$, put just
the row coordinates of $S$ at their roots.  All column coordinates and
remaining row coordinates stay away from their roots.  The $k$-fold
contact again has zero first variation.  No noise term, anticommutator,
or column Hamiltonian removes all these row roots.  The surviving row
Hamiltonian equals a nonzero common factor times
\[
 -i\sum_{\supp P=S}h_P\prod_{j\in S}v_{P_j}(a_j).
\]
The same finite tensor evaluation basis forces every displayed $h_P$ to
vanish.  Descending in the maximal support leaves only Hamiltonian terms
of weight at most two.  The transformed superoperator now acts on at most
two individual row or column legs at a time.

Fix distinct physical sites $i,j$.  A row--row double contact with coordinates
$s,t\in\mathcal E$ uses input parameters $a_i=s,a_j=t$.  A row--column
contact uses $a_i=s,a_j=-t$, which is available because $\mathcal E=-\mathcal E$.
Keep all other coordinates off their roots.  Only the corresponding mixed
principal coefficient survives.  Dividing by the initial quadratic
coefficient divides that coefficient by the positive product of the two
selected slopes.  Lemma~\ref{dyn:collision} therefore makes both mixed
coefficients real and nonpositive on this nine-point grid.

Let $K^{ij}$ be the $3\times3$ cross-site noise block, and write the
Hamiltonian pair as $\sum_{\alpha,\beta}h_{\alpha\beta}
\sigma_{i,\alpha}\sigma_{j,\beta}$ with $h$ real.  The row--column
coefficient matrix is $K^{ij}$.  Since distinct-site Paulis commute,
the two conjugate noise blocks in the anticommutator combine to give
the row--row coefficient matrix
\begin{equation}\label{dyn:cross-matrices}
 -\Re K^{ij}-ih.
\end{equation}
Here the real part is taken entrywise; there is no transpose in this formula.

Put
\[
 u(s)=\frac{(2s,1-s^2)}{1+s^2},\qquad R=\diag(1,-1).
\]
The normalized row and column derivative vectors are $(u(s),i)$ and
$(Ru(t),i)$.  Uniform averaging over $\mathcal E$ gives
\begin{equation}\label{dyn:trine-moments}
 \mathbb E u=0,\qquad \mathbb E uu^{\mathsf T}=\tfrac12I_2.
\end{equation}
Initially write $K^{ij}=\left(\begin{smallmatrix}A&B\\D&e\end{smallmatrix}\right)$
with complex blocks.  The imaginary part of its normalized row--column
coefficient is the imaginary part of
$u^{\mathsf T}ARv+i u^{\mathsf T}B+iDRv-e$.
Its nine values vanish.  Averaging these values, and then their products
with $u$, $Rv$, and $u(Rv)^{\mathsf T}$, proves
$\Im e=0$, $\Re B=0$, $\Re D=0$, and $\Im A=0$.
Hence $B=ib$, $D=id^{\mathsf T}$ with $A,b,d,e$ real, and
\begin{equation}\label{dyn:lr-inequality}
 u^{\mathsf T}ARv-b^{\mathsf T}u-d^{\mathsf T}Rv-e\le0.
\end{equation}
The imaginary part obtained from \eqref{dyn:cross-matrices} is
$-u^{\mathsf T}h_{TT}v+h_{zz}$, where $T=\{x,y\}$.
The same moments give $h_{TT}=0=h_{zz}$.  Write
$k=h_{Tz}$ and $\ell^{\mathsf T}=h_{zT}$.  The real row--row inequality is
\begin{equation}\label{dyn:ll-inequality}
 e-u^{\mathsf T}Av+k^{\mathsf T}u+\ell^{\mathsf T}v\le0.
\end{equation}
Both inequalities have already been divided by the positive factor
$(1+s^2)(1+t^2)$.  Their unweighted averages are respectively $-e$ and $e$,
so $e=0$.  Each of the nine nonpositive values in each inequality has
zero average, and thus each value is zero.  Applying
\eqref{dyn:trine-moments} once more gives $A=b=d=k=\ell=0$.
Consequently $K^{ij}=0$ and the pair Hamiltonian vanishes.

Every remaining block $K^{ii}$ is positive semidefinite.  The remaining
terms of \eqref{dyn:gkls} are therefore genuine one-qubit GKLS generators.
They commute on distinct tensor factors.  The proof compared only
derivatives at zero, so the different stable time sequences present no
compatibility issue.
\end{proof}

\begin{corollary}\label{dyn:immediate-failure}
If a CPTP generator is not a sum of one-qubit generators, one of the fixed
inputs in Theorem~\ref{dyn:trine} has a non-Lee--Yang output at every
sufficiently small positive time.
\end{corollary}

\begin{proof}
Otherwise each input would have Lee--Yang output at positive times
arbitrarily close to zero, providing the sequences in that theorem.
\end{proof}

The theorem makes no assertion that its number of tests is minimal or that
testing the output polynomial is computationally efficient.  Nor can a
single isolated positive time replace the sequences: the interaction
$H=Z_1Z_2$ generates the identity channel at $t=\pi$.

\subsection{The local reality forced by a second physical site}

We first record the elementary one-qubit state geometry.

\begin{lemma}\label{dyn:hemisphere}
A qubit density matrix $\rho=(I+r\cdot\sigma)/2$ is Lee--Yang if and only
if $r_z\ge0$.
\end{lemma}

\begin{proof}
Write $\rho=\left(\begin{smallmatrix}a&b\\\bar b&d\end{smallmatrix}\right)$,
so $a+d=1$, $|b|^2\le ad$, and
$F_\rho(z,w)=g(w)+z h(w)$ with $g=a+bw$, $h=\bar b+dw$.
On $|w|=1$,
\begin{equation}\label{dyn:bilinear-boundary}
 |g(w)|^2-|h(w)|^2
 =(a-d)[a+d+2\Re(bw)].
\end{equation}
If $a>d$, then $a>|b|$, so $g$ is zero-free on the closed disk.
The right side is nonnegative.  The maximum principle applied to $h/g$
proves $|h/g|\le1$ on the disk and hence the desired zero-freeness.
If $a=d$, add a positive multiple of $|0\rangle\langle0|$ and take its
nonzero coefficient limit; Hurwitz's theorem gives the assertion.
If $a<d$, positivity gives $d>|b|$, and
\eqref{dyn:bilinear-boundary} is strictly negative on the circle.
Thus a nearby interior point has $|h|>|g|$ and $h\ne0$;
choosing $z=-g/h$ gives an interior zero.  This proves both directions.
\end{proof}

\begin{proposition}\label{dyn:local-reality}
Suppose $n\ge2$ and a CPTP semigroup ordinarily preserves Lee--Yang
densities on a neighborhood of time zero.  In the local representation
given by Theorem~\ref{dyn:trine}, each noise matrix and Hamiltonian have
the form
\begin{equation}\label{dyn:real-local-shape}
 K=\begin{pmatrix}B&i d\\-i d^{\mathsf T}&c\end{pmatrix},
 \quad B=B^{\mathsf T}\in\R^{2\times2},\quad d\in\R^2,
 \qquad H=hZ.
\end{equation}
Writing $T=\Tr B$ and $J=\left(\begin{smallmatrix}0&-1\\1&0\end{smallmatrix}\right)$,
the local Bloch equation is
\begin{equation}\label{dyn:bloch}
 \dot r_\perp=M r_\perp+k,\qquad \dot r_z=-2T r_z,
 \quad M=2[B-(T+c)I_2]+2hJ,
 \quad k=4Jd.
\end{equation}
In particular the finite-time map has the form
\begin{equation}\label{dyn:affine}
 (r_\perp,r_z)\longmapsto(D_t r_\perp+\tau_t,\lambda_t r_z),
 \quad D_t\in GL_2(\R),\quad \lambda_t=e^{-2Tt}>0.
\end{equation}
\end{proposition}

\begin{proof}
After locality is known, let $q_j(s,w)$ denote the coefficient of
$\partial_s\partial_w$ in the Cayley-transformed generator at site $j$.
For a general Hermitian noise matrix write
\[
 K=\begin{pmatrix}
 \alpha&\delta+i\eta&r+i\mu\\
 \delta-i\eta&\beta&t+i\nu\\
 r-i\mu&t-i\nu&\gamma
 \end{pmatrix}.
\]
Multiplying the row and column derivative vectors in
\eqref{dyn:pauli-vectors} gives
\begin{equation}\label{dyn:imaginary-q}
 \Im q(s,w)=2(s+w)[-\eta(1-sw)+r(1+sw)+t(w-s)].
\end{equation}
The Hamiltonian and anticommutator act on single legs and do not enter
this coefficient.

Fix two sites.  For real $a,b$ and $c_0>0$, the polynomial
$f(z_1,z_2)=(z_1-a)(z_2-b)-c_0$ is upper-half-plane stable: solving for
$z_2$ when $z_1\in\HH$ gives a point in the lower half-plane.
Its inverse Cayley transform defines a pure physical Lee--Yang input.
The transformed conjugate factor is
$g(w_1,w_2)=(w_1+a)(w_2+b)-c_0$, since
$\overline C=\diag(-1,1)C$ and there are two sites.
For arbitrary nonzero real $u,v$, choose the real contact
\[
 z_1=a+u,\quad z_2=b+c_0/u,\qquad
 w_1=-a+v,\quad w_2=-b+c_0/v.
\]
Here $f=g=0$, $\nabla f=(c_0/u,u)$ and
$\nabla g=(c_0/v,v)$.  Their directional derivatives along any strictly
positive direction are nonzero.  The initial product has a double root
with nonzero real leading coefficient.  Lower-order generator terms
vanish at the contact, and its first variation is
\begin{equation}\label{dyn:hyperbola-contact}
 h=\frac{c_0^2}{uv}q_1(a+u,-a+v)
   +uv\,q_2(b+c_0/u,-b+c_0/v).
\end{equation}
Lemma~\ref{dyn:collision} makes $h$ real.  Additional sites may be padded
by equatorial factors kept away from their roots; their generators
contribute zero at this contact.

Formula \eqref{dyn:imaginary-q} gives $\Im q_2(b,-b)=0$.
The coefficient of $c_0$ in the imaginary part of
\eqref{dyn:hyperbola-contact} is therefore
\[
 v\,\partial_s\Im q_2(b,-b)+u\,\partial_w\Im q_2(b,-b)=0.
\]
Independent variation of $u,v$ makes both derivatives vanish.
By \eqref{dyn:imaginary-q}, this says
$-\eta(1+b^2)+r(1-b^2)-2tb=0$ for every real $b$.
Its three coefficients give $\eta=r=t=0$.  Interchanging the sites proves
the same conclusion at the other site, and every site has a partner.
This proves the stated shape of $K$.

At an equatorial Bloch vector the dissipator with this shape has zero
longitudinal velocity.  A Hamiltonian $h_xX+h_yY+h_zZ$ contributes
$\dot r_z=2(h_xr_y-h_yr_x)$.  Product inputs and
Lemma~\ref{dyn:hemisphere} require this quantity to be nonnegative at
every equatorial point.  Replacing that point by its negative gives
$h_x=h_y=0$.  Finally, multiplication of the Pauli matrices in the
generator gives \eqref{dyn:bloch}; in particular
$\mathcal L(I)=-4d_yX+4d_xY$, fixing the displacement sign.
The linear differential equation gives $D_t=e^{tM}$ and
$\lambda_t=e^{-2Tt}$, as claimed.
\end{proof}

\subsection{Separable conditional blocks prevent compensation}

\begin{lemma}[Conditional-block cancellation]\label{dyn:cancellation}
Let a qubit channel $\Psi$ have the affine form
$(r_\perp,r_z)\mapsto(D r_\perp+\tau,\lambda r_z)$, with $\lambda>0$.
If $\alpha,\beta\in\C$ satisfy $|\beta|>|\alpha|$, the polynomial of
$\alpha\Psi(|0\rangle\langle0|)+\beta\Psi(|1\rangle\langle1|)$ has a zero
in $\D^2$.
\end{lemma}

\begin{proof}
The scalar $\alpha+\beta$ is nonzero.  Dividing the matrix by
$(\alpha+\beta)/2$ gives $I+\tau\cdot\sigma_\perp+qZ$, where
\[
 q=\lambda\frac{\alpha-\beta}{\alpha+\beta},\qquad
 \Re q=\lambda\frac{|\alpha|^2-|\beta|^2}{|\alpha+\beta|^2}<0.
\]
Its polynomial is
$(1+q)+(\tau_x+i\tau_y)z+(\tau_x-i\tau_y)w+(1-q)zw$.
Choose a phase $\theta$ making
$(\tau_x+i\tau_y)e^{i\theta}=|\tau|$ and substitute
$z=e^{i\theta}u$, $w=e^{-i\theta}u$; if $\tau=0$ any phase works.
The quadratic has nonzero leading coefficient $1-q$ and root product
$(1+q)/(1-q)$ of modulus less than one.  At least one root has modulus
less than one, giving the asserted pair of disk coordinates.
\end{proof}

The lemma applies to the actual evolving channel at a second site.
It does not require that this channel fix the reference state.  In
particular, if positive matrices $A,B$ give a Lee--Yang input proportional
to $A\otimes|0\rangle\langle0|+B\otimes|1\rangle\langle1|$, ordinary
preservation by $\Phi\otimes\Psi$ forces
\begin{equation}\label{dyn:dominance-necessary}
 |F_{\Phi(B)}(z,w)|\le |F_{\Phi(A)}(z,w)|\qquad (z,w\in\D).
\end{equation}
Indeed a strict reverse inequality supplies $\alpha,\beta$ for
Lemma~\ref{dyn:cancellation}.  This also covers $\alpha=0$.

\begin{proposition}[Recovery of each local inequality]\label{dyn:local-inequality}
Under the assumptions of Proposition~\ref{dyn:local-reality}, each local
matrix in \eqref{dyn:real-local-shape} satisfies
\begin{equation}\label{dyn:circle-criterion}
 r^{\mathsf T}Bv+d\cdot(v-r)\le c
 \qquad (r,v\in\R^2,\ |r|=|v|=1).
\end{equation}
Necessity follows from positive definite separable inputs classical on
one of two selected sites.
\end{proposition}

\begin{proof}
Fix an active site and a distinct reference site.  For $0<\kappa<1/2$
and $e\in\R^2$ with $|e|=1$, put
\[
 U=I+\tfrac12Z,\qquad
 V=\tfrac18I+\tfrac{\sqrt{\kappa(1-\kappa)}}4\sigma_e
        +\tfrac\kappa4Z,\qquad \sigma_e=e_xX+e_yY,
\]
and define
\begin{equation}\label{dyn:separable-family}
 A=\tfrac12(U+V),\quad B_0=\tfrac12(U-V),\qquad
 \rho_{\kappa,e}=\tfrac12
 [A\otimes|0\rangle\langle0|+B_0\otimes|1\rangle\langle1|].
\end{equation}
Both blocks are positive definite, because
$\lambda_{\min}(U)=1/2$ and
$\|V\|=1/8+\sqrt\kappa/4<1/2$.
Their traces sum to $2$, so this is a normalized positive definite
separable density matrix.

For a Hermitian matrix $Q=q_0I+q_\perp\cdot\sigma_\perp+q_zZ$, parameterize
the distinguished boundary by
$z=e^{i(\gamma+\delta)}$, $w=e^{i(\gamma-\delta)}$, and set
$x=\cos\gamma$, $y=\sin\gamma$, $n=(\cos\delta,-\sin\delta)$.
Expansion of \eqref{dyn:polynomial} gives
\begin{equation}\label{dyn:torus-evaluation}
 \tfrac12e^{-i\gamma}F_Q(z,w)=q_0x+q_\perp\cdot n-iq_zy.
\end{equation}
Write $u=1/2$, $v_0=1/8$,
$w_0=2v_0\sqrt{\kappa(1-\kappa)}$ and
$\zeta=\kappa v_0/u$, so $U=I+uZ$ and
$V=v_0I+w_0\sigma_e+\zeta Z$.  Formula
\eqref{dyn:torus-evaluation} yields
\begin{align}\label{dyn:input-dominance}
 \mathscr D_0(x,n)
 &:=\tfrac14(|F_A|^2-|F_{B_0}|^2)\notag\\
 &=v_0[(1-\kappa)x^2
 +2\sqrt{\kappa(1-\kappa)}x(e\cdot n)+\kappa]\notag\\
 &\ge v_0(\sqrt{1-\kappa}|x|-\sqrt\kappa)^2\ge0.
\end{align}
The matrix $A$ is positive definite with strictly positive $Z$ coefficient.
Consequently its diagonal entries obey $a>d>0$, and
\eqref{dyn:bilinear-boundary} is strictly positive on the circle.
The proof of Lemma~\ref{dyn:hemisphere} then gives zero-freeness of $F_A$
on the closed bidisk.  Applying the maximum principle separately in
$z,w$ to $F_{B_0}/F_A$ extends \eqref{dyn:input-dominance} to that bidisk.
Hence $F_A+z_2w_2F_{B_0}$ is nonzero on $\D^4$, proving that
\eqref{dyn:separable-family} is Lee--Yang.  Other sites may be padded by
maximally mixed states.

For the active map write $D_t,\tau_t,\lambda_t$ as in
\eqref{dyn:affine}.  At time zero their derivatives are $M,k,-2T$ in
\eqref{dyn:bloch}.  The reference map has $\lambda_t>0$ at every finite
time.  By \eqref{dyn:dominance-necessary}, ordinary preservation forces
the active output block difference to be nonnegative on the disk and,
by continuity, on its boundary.  Its exact normalized expression is
\begin{equation}\label{dyn:output-dominance}
 \mathscr D_t(x,n)=
 (x+\tau_t\cdot n)
 [v_0(x+\tau_t\cdot n)+w_0n^{\mathsf T}D_te]
 +\lambda_t^2u\zeta(1-x^2).
\end{equation}
Choose $n=-e$ and $x=\rho:=\sqrt{\kappa/(1-\kappa)}\in(0,1)$.
The input difference is zero.  Its derivative must therefore be
nonnegative.  Differentiating \eqref{dyn:output-dominance} gives
\begin{align}\label{dyn:derivative-dominance}
 \frac{\mathscr D'_0}{2v_0\kappa}
 &=\rho k\cdot n-n^{\mathsf T}Mn
      -2T\frac{1-2\kappa}{1-\kappa}\notag\\
 &=2(c-n^{\mathsf T}Bn)+2T\rho^2+\rho k\cdot n\ge0.
\end{align}
For explicit verification, the three derivative contributions before
division are $(k\cdot n)(2v_0\rho-w_0)$,
$-\rho w_0n^{\mathsf T}Mn$, and
$-4Tu\zeta(1-\rho^2)$.  The Hamiltonian term drops out because
$n^{\mathsf T}Jn=0$.

Set $n=Jm$, $|m|=1$.  Since $k=4Jd$ and
$(Jm)^{\mathsf T}B(Jm)=T-m^{\mathsf T}Bm$, we obtain
\begin{equation}\label{dyn:radial-inequality}
 c-T+m^{\mathsf T}Bm+T\rho^2+2\rho\,m\cdot d\ge0
 \qquad (|m|=1,\ 0\le\rho\le1).
\end{equation}
The endpoints follow by continuity.
For arbitrary unit $r,v$, set $(v-r)/2=\rho m$.
Their mean is perpendicular to $m$ and has squared length $1-\rho^2$;
when $\rho=0$, choose either perpendicular direction for $m$.
In two dimensions this gives
\[
 r^{\mathsf T}Bv+d\cdot(v-r)
 =T-m^{\mathsf T}Bm-T\rho^2+2\rho\,m\cdot d.
\]
Apply \eqref{dyn:radial-inequality} with $-m$ to prove
\eqref{dyn:circle-criterion}.
\end{proof}

\subsection{The exact local semidefinite cone}

\begin{lemma}[Reduction of the circle inequality]\label{dyn:lmi-reduction}
Suppose $B=B^{\mathsf T}\in\R^{2\times2}$ and
$\left(\begin{smallmatrix}B&i d\\-i d^{\mathsf T}&c\end{smallmatrix}\right)\succeq0$.
Then \eqref{dyn:circle-criterion} is equivalent to
\begin{equation}\label{dyn:second-lmi}
 \begin{pmatrix}cI_2-\operatorname{adj}B&d\\d^{\mathsf T}&\Tr B\end{pmatrix}
 \succeq0.
\end{equation}
If $T=\Tr B>0$, the equivalent scalar condition is
$c\ge\lambda_{\max}(\operatorname{adj}B+dd^{\mathsf T}/T)$.
If $T=0$, positivity forces $B=d=0$, and every $c\ge0$ is allowed.
\end{lemma}

\begin{proof}
Unitary congruence, followed if necessary by changing the sign of the last
coordinate, gives
$\left(\begin{smallmatrix}B&d\\d^{\mathsf T}&c\end{smallmatrix}\right)\succeq0$.
In particular $B\succeq0$ and $dd^{\mathsf T}\preceq cB$.
The assertion for $T=0$ follows immediately.  Suppose $T>0$.
Writing $(v-r)/2=\rho m$ as above, the expression to maximize is
\[
 T-m^{\mathsf T}Bm-T\rho^2+2\rho\,d\cdot m,
 \qquad |m|=1,\quad0\le\rho\le1.
\]
If $c<T$, then $|d|^2\le c\lambda_{\max}(B)<T^2$.
Choose the sign of $m$ so $d\cdot m\ge0$.
The unconstrained maximum in $\rho$ is attained at
$\rho=|d\cdot m|/T<1$.  The maximum is consequently
\[
 \lambda_{\max}\left(TI_2-B+\frac{dd^{\mathsf T}}T\right)
 =\lambda_{\max}\left(\operatorname{adj}B+\frac{dd^{\mathsf T}}T\right).
\]
This proves equivalence in this case.

If $c\ge T$, put $b_m=m^{\mathsf T}Bm$ and use
$|d\cdot m|\le\sqrt{cb_m}$.  The difference between $c$ and the expression
above is bounded below by
\[
 c-T+b_m+T\rho^2-2\sqrt{cb_m}\rho
 =(\sqrt{b_m}-\sqrt c\rho)^2+(c-T)(1-\rho^2)\ge0.
\]
The scalar matrix condition also holds automatically, because
\[
 dd^{\mathsf T}\preceq cB\preceq T[B+(c-T)I_2],
\]
the second inequality being equivalent to
$(c-T)(TI_2-B)\succeq0$.
Finally, the Schur complement with positive last entry $T$ proves
equivalence of the scalar condition and \eqref{dyn:second-lmi}.
The separate $c\ge T$ argument avoids an invalid unconstrained radial
maximization when $|d|>T$.
\end{proof}

\begin{lemma}[Choi tensors and complete preservation]\label{dyn:choi}
For a CPTP map $\Phi$ on $n$ qubits, the following are equivalent:
its full Choi tensor is Lee--Yang; it completely preserves physical
Lee--Yang states; and its extension by $n$ inert qubits takes the maximally
entangled state to a Lee--Yang state.
For any invertible complex-linear map, without a positivity or trace
assumption, its full Choi tensor is Lee--Yang if and only if its
coefficient action preserves all complex disk-stable polynomials in the
$2n$ row and column variables.
\end{lemma}

\begin{proof}
With $|\Omega\rangle=\sum_a|a\rangle|a\rangle$, the Choi entries are
$J_{ai,bj}=[\Phi(|i\rangle\langle j|)]_{ab}$.
The normalized maximally entangled input is Lee--Yang, since its polynomial
is proportional to $\prod_k(1+z_kz'_k)(1+w_kw'_k)$.
Its output is the normalized Choi matrix.
Conversely, for a joint input $\rho$, the output coefficients are
\[
 [ (\Phi\otimes\Id)(\rho)]_{ar,bs}
 =\sum_{i,j}J_{ai,bj}\rho_{ir,js}.
\]
This is bilinear tensor contraction with no conjugation of entries.
Lemma~\ref{pre:grace} gives a Lee--Yang polynomial or zero;
trace preservation excludes zero.
This proves the physical equivalences.

Let $M$ be the superoperator matrix in coefficient coordinates.  Its full
tensor differs from the Choi tensor only by a permutation of indices.
Contraction with an arbitrary stable coefficient vector gives stability
or zero, and invertibility excludes zero.  In the other direction, if
$M$ preserves all stable polynomials, apply it for each fixed
$w\in\D^{2n}$ to $\prod_j(1+z_jw_j)$.
The resulting joint polynomial is exactly the full tensor of $M$ and
is nonzero for every $(z,w)\in\D^{4n}$.
\end{proof}

\begin{proposition}[One-qubit complete generators]\label{dyn:local-complete}
A one-qubit CPTP semigroup is completely preserving if and only if its
generator has the form \eqref{dyn:real-local-shape} and satisfies the two
positive semidefinite conditions in \eqref{dyn:two-lmis}.
\end{proposition}

\begin{proof}
For necessity, extend the semigroup by one inert qubit.
Propositions~\ref{dyn:local-reality} and \ref{dyn:local-inequality} give the
real local shape and \eqref{dyn:circle-criterion} for the active site.
Lemma~\ref{dyn:lmi-reduction} gives the two matrix conditions.

For sufficiency, transform both coefficient legs by \eqref{dyn:cayley}.
Write $B=\left(\begin{smallmatrix}a&u\\u&b\end{smallmatrix}\right)$ and
$d=(x,y)^{\mathsf T}$.  The transformed generator is real.  This follows
directly from the Pauli identities: the mixed terms pair the purely
imaginary transverse--longitudinal noise entries with the purely imaginary
$Z$ derivative; the anticommutator is
$-(a+b+c)I+y(Z_1+Z_2)+x(X_1-X_2)$; and the longitudinal Hamiltonian is real
after the transform.  Its full coefficient of $\partial_s\partial_t$ is
\begin{align}\label{dyn:local-q}
 q(s,t)={}&4ast-b(1-s^2)(1-t^2)-c(1+s^2)(1+t^2)\notag\\
 &+2u[t(1-s^2)-s(1-t^2)]\notag\\
 &+2x[t(1+s^2)-s(1+t^2)]
 -y[(1-s^2)(1+t^2)+(1+s^2)(1-t^2)].
\end{align}
There is no extra factor of two in this convention.  With
\[
 r(s)=\frac{(2s,1-s^2)}{1+s^2},\qquad
 v(t)=\frac{(2t,t^2-1)}{1+t^2},
\]
we have
\begin{equation}\label{dyn:q-circle}
 \frac{q(s,t)}{(1+s^2)(1+t^2)}
 =r(s)^{\mathsf T}Bv(t)+d\cdot(v(t)-r(s))-c\le0.
\end{equation}
The real two-variable specialization of the finite-box generator theorem
(Theorem~\ref{gen:main-complex}) says that a real endomorphism of the binary
polynomial box generates upper-half-plane stability preservers precisely
when this mixed coefficient is nonpositive on $\R^2$.
It applies here because the transformed operator is an endomorphism of
that box.  Thus its exponential preserves complex stability.
Undoing the Cayley similarity and applying Lemma~\ref{dyn:choi}, with
invertibility of the exponential, proves complete preservation.
\end{proof}

\begin{proof}[Proof of Theorem~\ref{dyn:main}]
Assume preservation on a neighborhood of zero.  Theorem~\ref{dyn:trine}
gives locality.  Propositions~\ref{dyn:local-reality},
\ref{dyn:local-inequality}, and Lemma~\ref{dyn:lmi-reduction} give
\eqref{dyn:local-form}--\eqref{dyn:two-lmis}.  Conversely those conditions
give complete preservation of every local semigroup by
Proposition~\ref{dyn:local-complete}.  Their tensor product completely
preserves by successive application to the individual sites.
Lemma~\ref{dyn:choi} equates this with the full Choi condition.
Complete preservation implies ordinary preservation, which implies its
small-time version.  The one-qubit exception is proved next.
\end{proof}

\subsection{The Hermiticity-preserving branch without positivity}

For the full Choi condition, locality does not require complete positivity.
The local semidefinite reduction does require it, and must be replaced by
the circle inequality when positivity is absent.

\begin{theorem}[Hermiticity- and trace-preserving semigroups]\label{dyn:hp}
Let $T_t=e^{t\mathcal L}$ be a continuous semigroup of complex-linear,
Hermiticity-preserving, trace-preserving maps on $n$ qubits.  Its full
Choi tensor is Lee--Yang for every $t\ge0$ if and only if
$\mathcal L=\sum_j\mathcal L_j\otimes\Id_{\widehat j}$, where each
$\mathcal L_j$ has the algebraic form \eqref{dyn:local-form}, with
\[
 H_j=h_jZ,\qquad
 K_j=\begin{pmatrix}B_j&i d_j\\-i d_j^{\mathsf T}&c_j\end{pmatrix},
 \quad B_j=B_j^{\mathsf T}\in\R^{2\times2},\quad
 d_j\in\R^2,\quad c_j,h_j\in\R,
\]
and satisfies
\begin{equation}\label{dyn:hp-circle}
 c_j\ge\max_{r,v\in S^1}
 [r^{\mathsf T}B_jv+d_j\cdot(v-r)].
\end{equation}
Here no positive semidefiniteness of $K_j$ is imposed.  Requiring the
Choi property only on a neighborhood of zero is equivalent.  If complete
positivity is also imposed, the local conditions reduce exactly to
\eqref{dyn:two-lmis}.
\end{theorem}

\begin{proof}
We first justify the algebraic representation used in the proof.
A Hermiticity-preserving, trace-annihilating generator has a representation
\eqref{dyn:gkls} with $H=H^*$ and $K=K^*$, without the condition
$K\succeq0$.  Indeed, its Choi matrix is Hermitian.  In the orthogonal
vectorized Pauli basis, its compression to the orthogonal complement
of $|I\rangle\!\rangle$ uniquely specifies the Hermitian coefficients
$K_{PQ}$ of the sandwich terms.  Subtract the corresponding dissipator,
which is Hermiticity-preserving and trace-annihilating.  The remaining
Hermitian Choi matrix has zero compression to that complement, hence
is of the form
$|B\rangle\!\rangle\langle\!\langle I|
 +|I\rangle\!\rangle\langle\!\langle B|$ for some matrix $B$.
It represents $\rho\mapsto B\rho+\rho B^*$.
Trace annihilation gives $B+B^*=0$, so $B=-iH$ with $H$ Hermitian.
This proves the claimed representation.

The last assertion of Lemma~\ref{dyn:choi} and the Cayley transform show
that the full Choi hypothesis places the coefficient generator in
Theorem~\ref{gen:main-complex} on $2n$ binary legs.
It consequently has differential order at most two.
This implies tensor support at most two on those individual legs:
expand in the local endomorphism basis
$I,\partial_s,s\partial_s-1/2,s^2\partial_s-s$.
At a nonzero support of maximal size, its full derivative coefficient
is a linear combination of independent products of $1,s,s^2$.
If the support size exceeds two, the order bound forces all its
coefficients to vanish.  Descending eliminates every larger support.
Local Cayley similarity preserves the identity and traceless
subspaces, so this conclusion also holds before the transform.

For nonidentity Pauli strings $P,Q$, the tensor coefficient supported
on their respective row and column legs is exactly the sandwich
coefficient $K_{PQ}$, up to the harmless sign of a transposed $Y$.
Neither a Hamiltonian nor an anticommutator term can cancel a
coefficient involving both row and column legs.  Therefore
\[
 \operatorname{wt}(P)+\operatorname{wt}(Q)>2
 \quad\Longrightarrow\quad K_{PQ}=0.
\]
In particular every row indexed by a string of weight at least two
vanishes, directly and without a PSD argument.  The anticommutator
now has support at most two, so the same support bound excludes every
Hamiltonian term of higher weight.

For distinct physical sites, the cross-site matrices and their normalized
mixed coefficients are precisely \eqref{dyn:cross-matrices}--
\eqref{dyn:ll-inequality}.  The general generator theorem makes these
coefficients real and nonpositive at every real pair.  Restricting to
the trine grid, the discrete moment calculation
\eqref{dyn:trine-moments} eliminates every cross-site block and pair
Hamiltonian exactly as in Theorem~\ref{dyn:trine}.  That calculation
uses Hermiticity and trace annihilation, but never PSD of $K$.
Thus the generator is a sum of local Hermiticity-preserving,
trace-annihilating generators.

At one site, the nine products of the two Pauli derivative triples are
independent.  Reality of their mixed coefficient forces $K_{xy}$ to
be real and $K_{xz},K_{yz}$ to be purely imaginary, giving the displayed
local shape.  Its dissipator is then real after the Cayley transform.
For a Hamiltonian $h_xX+h_yY+h_zZ$, the remaining imaginary first-order
coefficients are
\[
 h_y s^2-2h_xs-h_y,\qquad h_y t^2+2h_xt-h_y.
\]
The complex generator criterion requires each quadratic to be
nonnegative on the real line.  Comparing leading and constant terms
forces $h_y=0$, and then the linear term forces $h_x=0$.
The real mixed coefficient is \eqref{dyn:local-q}, whose normalization
\eqref{dyn:q-circle} is nonpositive exactly under \eqref{dyn:hp-circle}.
This proves necessity without any positivity assumption.

Conversely, the displayed local algebraic generators are
Hermiticity-preserving and trace-annihilating.  Their Cayley transforms
are real binary-box endomorphisms with nonpositive mixed coefficient.
Theorem~\ref{gen:main-complex} and Lemma~\ref{dyn:choi} prove full Choi
stability of their exponentials and their tensor product.  Small-time
preservation extends to all times by the semigroup identity and
Lemma~\ref{pre:grace}, with invertibility excluding zero.
Finally, complete positivity imposes $K_j\succeq0$ in the canonical
representation, and Lemma~\ref{dyn:lmi-reduction} gives exactly the
second matrix condition in \eqref{dyn:two-lmis}.
\end{proof}

\subsection{Cancellation at a fixed positive time}

The mechanism used above has a finite-time version.  It applies to different
local channels and does not assume that they belong to a semigroup.

\begin{theorem}[Static tensor cancellation]\label{dyn:static-cancellation}
Let $n\ge2$.  For each $j$, let $\Psi_j$ be an invertible CPTP qubit map
with Bloch action
\[
 (r_\perp,r_z)\longmapsto(D_jr_\perp+\tau_j,\lambda_jr_z),
 \qquad D_j\in GL_2(\R),\quad\lambda_j>0.
\]
Then $\bigotimes_j\Psi_j$ ordinarily preserves all physical Lee--Yang
densities if and only if every $\Psi_j$ has a Lee--Yang Choi tensor.
If one factor fails this condition, a positive definite separable input,
classical on a second selected site and maximally mixed on the other sites,
detects the failure, whatever the parameters of the second factor.
\end{theorem}

\begin{proof}
The outputs of the two computational basis states give
$|\tau|^2+\lambda^2\le1$, so $|\tau|<1$.
The static one-qubit criterion in Theorem~\ref{loc:classification} says that
the Choi condition for the active factor $\Psi$ is equivalent to
\begin{equation}\label{dyn:static-criterion}
 \|D^{\mathsf T}n\|^2
 \le(x-\tau\cdot n)^2+\lambda^2(1-x^2)
 \qquad(n\in S^1,\ -1\le x\le1).
\end{equation}
Changing $x$ to $-x$ gives the alternative sign convention for its
transverse displacement.
Suppose the condition fails.  By strictness and continuity choose
$x\in(-1,1)$ and $n\in S^1$ such that, with $v=D^{\mathsf T}n$,
$V=|v|$ and $y=x-\tau\cdot n$,
\begin{equation}\label{dyn:static-failure}
 V^2>y^2+\lambda^2(1-x^2).
\end{equation}
In particular $V>0$.  Set $u=yv/V^2$, so $|u|<1$ and $u\cdot v=y$.
The matrix $W=I_2-uu^{\mathsf T}$ is positive definite and
$v^{\mathsf T}Wv=V^2-y^2>\lambda^2(1-x^2)$.
Choose $0<\theta<1$ sufficiently close to one and put
\[
 d_0=\theta\frac{Wv}{\sqrt{v^{\mathsf T}Wv}}.
\]
Then $W-d_0d_0^{\mathsf T}\succ0$, as follows by congruence with
$W^{-1/2}$, and
\begin{equation}\label{dyn:static-separation}
 d_0\cdot v>\lambda\sqrt{1-x^2}.
\end{equation}
Choose $\beta>0$ with $\beta(1-|u|)>|d_0|$ and define
\[
 A_R=R(I+u\cdot\sigma_\perp)+Z,
 \qquad B_0=\beta I+(\beta u+d_0)\cdot\sigma_\perp.
\]
The second matrix is positive definite.  The first is positive definite
once $R^2(1-|u|^2)>1$.

We claim that, for all sufficiently large finite $R$,
$F_{A_R}$ is zero-free on the closed bidisk and
\begin{equation}\label{dyn:static-input-domination}
 |F_{B_0}(s,w)|<|F_{A_R}(s,w)|\qquad(|s|,|w|\le1).
\end{equation}
On the distinguished boundary use \eqref{dyn:torus-evaluation} with
$X=\cos\gamma$, $m\in S^1$, and $\ell=X+u\cdot m$.
The squared-modulus difference divided by four is
\[
 G_R(X,m)=(R^2-\beta^2)\ell^2
 -2\beta\ell(d_0\cdot m)+1-X^2-(d_0\cdot m)^2.
\]
On $\ell=0$, its last two terms are at least
$\delta:=\lambda_{\min}(I_2-uu^{\mathsf T}-d_0d_0^{\mathsf T})>0$.
If $|\ell|\le\epsilon$, then
$|X^2-(u\cdot m)^2|\le2\epsilon$.
Choose $\epsilon>0$ so
$2\epsilon+2\beta\epsilon|d_0|<\delta/2$.
For $R\ge\beta$, this gives $G_R>0$ on that set.
On $|\ell|\ge\epsilon$, using $|\ell|\le2$ gives
\[
 G_R\ge(R^2-\beta^2)\epsilon^2-4\beta|d_0|-|d_0|^2>0
\]
for sufficiently large $R$.  This proves strict torus domination.

The diagonal entries of $A_R$ are $R+1,R-1$ and its off-diagonal
magnitude is $R|u|$.  The constant coefficient in the representation
$F_{A_R}=g(w)+s h(w)$ is nonzero on the closed disk because
$R+1>R|u|$.  On the circle,
\[
 |g(w)|^2-|h(w)|^2
 =4R+4\Re((A_R)_{01}w)>0.
\]
The maximum principle gives closed-bidisk zero-freeness of $F_{A_R}$.
Apply it separately in both coordinates to $F_{B_0}/F_{A_R}$ to obtain
\eqref{dyn:static-input-domination}.
Consequently the density proportional to
$A_R\otimes|0\rangle\langle0|+B_0\otimes|1\rangle\langle1|$
is positive definite, separable, and Lee--Yang.

At the output torus point in \eqref{dyn:torus-evaluation} with
$X=-x$, $m=n$ and $Y=\sqrt{1-x^2}$, the affine action gives
\[
 \tfrac12e^{-i\gamma}F_{\Psi(A_R)}
 =R[-x+\tau\cdot n+u\cdot v]-i\lambda Y=-i\lambda Y,
\]
and
\[
 \tfrac12e^{-i\gamma}F_{\Psi(B_0)}
 =\beta[-x+\tau\cdot n+u\cdot v]+d_0\cdot v=d_0\cdot v.
\]
The first value is nonzero, while \eqref{dyn:static-separation} makes
the second strictly larger in modulus.  The strict inequality persists
when the active coordinates move slightly into the disk.
Lemma~\ref{dyn:cancellation} then gives a zero of the two-site output
after any allowed partner channel.  Padding by maximally mixed states
proves necessity for every $n\ge2$.
Conversely, if all factors have Lee--Yang Choi tensors,
Lemma~\ref{dyn:choi} gives complete preservation at each site and therefore
ordinary preservation of their tensor product.
\end{proof}

The theorem assumes the displayed local affine form.  It is not a
factorization theorem for arbitrary isolated channels.  In particular,
the continuous-time locality theorem does not by itself imply such a
static conclusion.

\subsection{The one-qubit exception and dynamical consequences}

\begin{proposition}[The ordinary one-qubit cone]\label{dyn:one-qubit}
Let a one-qubit CPTP semigroup have Bloch equation
$\dot r=Ar+b$, with real $A,b$.  It ordinarily preserves Lee--Yang states
if and only if
\begin{equation}\label{dyn:hemisphere-generator}
 b_z\ge\sqrt{A_{zx}^2+A_{zy}^2}.
\end{equation}
\end{proposition}

\begin{proof}
By Lemma~\ref{dyn:hemisphere}, the additional condition is invariance of
the northern Bloch hemisphere.  At $r_z=0$, its minimum outward-coordinate
velocity is $b_z-\sqrt{A_{zx}^2+A_{zy}^2}$, since the transverse vector
ranges over the closed unit disk.  A negative value gives immediate
failure.  Conversely, complete positivity keeps $|r(t)|\le1$.
Under \eqref{dyn:hemisphere-generator},
$\dot r_z\ge A_{zz}r_z$ at every time.  Multiplication by
$e^{-A_{zz}t}$ and integration give
$r_z(t)\ge e^{A_{zz}t}r_z(0)\ge0$.
\end{proof}

\begin{example}[Amplitude damping]\label{dyn:amplitude-damping}
The amplitude-damping channel toward $|0\rangle$ has Kraus operators
$K_0=\diag(1,\sqrt\eta)$ and
$K_1=\sqrt{1-\eta}|0\rangle\langle1|$, where $0\le\eta\le1$.
The semigroup $\eta=e^{-\gamma t}$ satisfies
$\dot r_z=\gamma(1-r_z)$ and therefore ordinarily preserves one-qubit
Lee--Yang states.  It is not completely preserving for $0<\eta<1$.
Indeed, its action on one half of a Bell state gives, up to a positive
scalar, the polynomial
\[
 (1+\sqrt\eta\,z_1z_2)(1+\sqrt\eta\,w_1w_2)
 +(1-\eta)z_2w_2.
\]
At $z_1=z_2=w_1=w_2=ir$ its value is
$1-(2\sqrt\eta+1-\eta)r^2+\eta r^4$.
It is positive at $r=0$ and equals $2\eta-2\sqrt\eta<0$ at $r=1$.
An interior root $0<r<1$ proves the claim.
\end{example}

\begin{remark}\label{dyn:rates}
For diagonal $K=\diag(a,b,c)$ the complete criterion is
$c\ge\max(a,b)$, recovering the Pauli-semigroup portion of the known
sufficient channel family \cite[Theorem~3]{WBG26}. The full cone also permits
nonunital noise: $\sqrt\gamma(X+iZ)/2$ is an allowed jump and damps
toward an equatorial pure state. Within the completely preserving cone,
only rotations around the fixed $Z$
axis are allowed as coherent motion.  More generally, ordinary
one-qubit dynamics with Hamiltonian
$(\omega_xX+\omega_yY+\omega_zZ)/2$, downward and upward jump rates
$\gamma_\downarrow,\gamma_\uparrow$, and arbitrary $Z$ dephasing
satisfies \eqref{dyn:hemisphere-generator} exactly when
$\gamma_\downarrow-\gamma_\uparrow\ge
\sqrt{\omega_x^2+\omega_y^2}$.
\end{remark}

\begin{corollary}\label{dyn:no-correlated-preparation}
An ordinarily preserving CPTP semigroup on $n\ge2$ qubits cannot take a
product input to a nonproduct state, either at a finite time or in a limit
as $t\to\infty$.  In particular, a semigroup converging from every input
to a fixed nonproduct density matrix cannot ordinarily preserve the full
Lee--Yang class.
\end{corollary}

\begin{proof}
Theorem~\ref{dyn:main} keeps every product input a tensor product.
The set of product densities is closed, being the continuous image of
a compact product of one-site state spaces.
\end{proof}

This conclusion concerns universal preservation and preparation from
product inputs.  It does not exclude a correlated stationary state of a
nonergodic product semigroup, nor preservation along one specially chosen
trajectory.  The identity semigroup illustrates the first distinction.

\subsection{Explicit boundaries of restricted tests}

The following examples distinguish preservation of selected inputs or
classical populations from preservation of the whole physical class.
Their zeros provide direct certificates independent of the classification.

\begin{lemma}[Pair defects of real hyperbola contacts]\label{dyn:pair-defect}
Let $\mathcal L$ be a sum of local GKLS generators of the real form
\eqref{dyn:real-local-shape}, on $n\ge2$ sites, and let $q_j$ be its
local mixed coefficient \eqref{dyn:local-q}.  Define
\[
 \Delta_j=\sup_{s,w\in\R,\ s+w\ne0}
             \frac{q_j(s,w)}{(s+w)^2}.
\]
All the real double-contact inequalities from the two-site pure
hyperbola family \eqref{dyn:hyperbola-contact} hold if and only if
every $\Delta_j$ is finite and
\begin{equation}\label{dyn:pair-defect-condition}
 \Delta_i+\Delta_j\le0\qquad(i\ne j).
\end{equation}
In particular, these are necessary conditions for ordinary preservation.
They are not asserted to be sufficient for ordinary preservation.
\end{lemma}

\begin{proof}
Fix a pair of sites and number them one and two.
The double-root leading coefficient of the product of the two hyperbola
factors in the proof of Proposition~\ref{dyn:local-reality} has sign $uv$.
Thus its contact inequality is
\[
 \frac{c_0^2}{u^2v^2}q_1(a_0+u,-a_0+v)
 +q_2(b_0+c_0/u,-b_0+c_0/v)\le0.
\]
Writing $S=u+v$ and $R=c_0S/(uv)$, when $S\ne0$ this is the sum of the
two ratios defining $\Delta_1,\Delta_2$.  Every pair of real points with
nonzero sums $S,R$ is attained: choose a nonzero real $P$ with the sign
of $S/R$, taking $0<P<S^2/4$ when it is positive; let $u,v$ be the roots
of $X^2-SX+P$, put $c_0=PR/S>0$, and choose $a_0,b_0$ to fit the first
coordinate in each pair.  Thus the contact inequalities with nonzero
sums say that every value of the first ratio plus every value of the
second is nonpositive.  Fixing a finite value of either ratio bounds
the other ratio above.  Each ratio also has at least one finite value,
so both suprema are real and their sum is nonpositive.  The converse
follows directly from the definition of the suprema.

When $S=0$, also $R=0$.  The paired Pauli derivative vectors obey
$\widetilde v(-s)=-\overline{v(s)}$ in \eqref{dyn:pauli-vectors}, so positivity of the local noise matrix
gives
\[
 q_j(s,-s)=-v(s)^{\mathsf T}K_j\overline{v(s)}\le0.
\]
These contacts therefore impose no additional condition.  For $n>2$,
pad the input by pure local factors that are nonzero at the selected
contact.  Local terms on the other sites then vanish to second order
at that contact, leaving exactly the displayed two-site inequality.
This proves the claim for every pair and the necessity statement.
\end{proof}

\begin{proposition}[Real hyperbola contacts are insufficient]\label{dyn:compensation}
For $a>0$ and $c\ge0$, the two-qubit generator
\[
 \mathcal L(\rho)=a(X_1\rho X_1-\rho)+c(Z_2\rho Z_2-\rho)
\]
does not ordinarily preserve Lee--Yang densities.  Nevertheless, when
$c\ge a$, it satisfies every real double-contact inequality obtained
from the pure input family used in \eqref{dyn:hyperbola-contact}.
\end{proposition}

\begin{proof}
Here $q_1(s,w)=4asw$ and
$q_2(s,w)=-c(1+s^2)(1+w^2)$.  The identities
$4sw\le(s+w)^2$ and
$(1+s^2)(1+w^2)-(s+w)^2=(sw-1)^2$ give
$\Delta_1=a$, $\Delta_2=-c$, with equality attained in each supremum.
Lemma~\ref{dyn:pair-defect} proves the assertion about all real
hyperbola contacts.

For actual failure, use the positive input
\[
 \rho=\frac27\left[
 \begin{pmatrix}2&0\\0&1\end{pmatrix}\otimes|0\rangle\langle0|
 +\frac14\begin{pmatrix}1&1\\1&1\end{pmatrix}\otimes|1\rangle\langle1|
 \right].
\]
It has trace one.  Up to its positive scalar its polynomial is
$2+sw+\tfrac14(1+s)(1+w)zv$, which is zero-free on $\D^4$ because
$|2+sw|>1$ and $|\tfrac14(1+s)(1+w)zv|<1$ there.
The second-site dephasing vanishes on this input and its entire
first-site trajectory.  At $t=(\log2)/a$, the first diagonal block is
$\diag(13/8,11/8)$, whereas the second block is fixed.  At $s=w=19i/20$
the output polynomial becomes
\[
 \frac{1229}{3200}+\frac{39+760i}{1600}zv.
\]
Choose $zv=-1229/[2(39+760i)]$.  Its modulus is $1229/1522<1$, since
$39^2+760^2=761^2$.  Taking both $z,v$ to be the same square root puts
all variables in the open disk and makes the output zero.
\end{proof}

\begin{example}[Random interchange and incoherent exclusion]\label{dyn:exchange-examples}
Let $S$ interchange two qubits.  For $0<p<1$, the channel
$\rho\mapsto p\rho+(1-p)S\rho S$ fails ordinary preservation.
The Lee--Yang input $|+0\rangle\langle+0|$ gives an output polynomial
proportional to
$p(1+z_1)(1+w_1)+(1-p)(1+z_2)(1+w_2)$.
It vanishes at
\[
 z_1=w_1=-1+\sqrt{1-p}\,e^{i\pi/4},\qquad
 z_2=w_2=-1+\sqrt p\,e^{-i\pi/4}.
\]
Every coordinate is in the open disk: for $0<x<1$,
$|-1+xe^{\pm i\pi/4}|^2=1+x^2-\sqrt2x<1$.
In particular the random-interchange semigroup fails at each positive
time, since its mixture parameter is $(1+e^{-2\gamma t})/2$.

The two jumps $\sqrt\gamma|01\rangle\langle10|$ and
$\sqrt\gamma|10\rangle\langle01|$, with $\gamma>0$, also fail at all
sufficiently small positive times.  On the same input, let
$u=e^{-\gamma t/2}$.  The $00$--$10$ coherence decays by $u$, and the two
excited populations exchange at rate $2\gamma$, giving
\[
 2F_t=1+u(z_1+w_1)+\tfrac{1+u^4}{2}z_1w_1
                         +\tfrac{1-u^4}{2}z_2w_2.
\]
At $z_1=w_1=-r$, $z_2=w_2=ir$, this is $1-2ur+u^4r^2$.
Its value at $r=0$ is positive.  At $r=1$ it is $u^4-2u+1$, negative
for $u<1$ sufficiently close to one, because this function vanishes at
one and has derivative two there.  The intermediate value theorem
supplies an interior zero.  Thus this natural quantum extension of
population exchange does not preserve the entire physical class.
\end{example}

\begin{example}[Equal local damping does not restore preservation]\label{dyn:equal-damping}
Apply the same amplitude-damping channel to both halves of a Bell state.
Its output polynomial, up to the positive factor $1/2$, is
\[
 1+(1-\eta)^2+\eta(1-\eta)(z_1w_1+z_2w_2)
 +\eta^2z_1z_2w_1w_2+\eta(z_1z_2+w_1w_2).
\]
At all four variables $ir$, the polynomial is
$1+(1-\eta)^2-2\eta(2-\eta)r^2+\eta^2r^4$.
It is positive at zero and equals $2(1-\eta)(1-2\eta)<0$ at one for
$1/2<\eta<1$.  Hence equal local damping fails ordinary joint
preservation at every sufficiently small positive time.  For instance,
$\eta=3/4$ has the exact interior root
$r^2=(5-2\sqrt2)/3$.  Vacuum padding extends the example to any larger
system.
\end{example}


\section{Physical dynamics at arbitrary finite spins}
\label{dynspin:subsection}

The physical classification becomes more restrictive when a site has spin
at least one.  We keep the physical basis fixed and use the coherent
normalization of Section~\ref{pre:section}.  For positive integers
$\kappa_i=2j_i$, put
\[
 \mathcal H=\bigotimes_{i=1}^n\C^{\kappa_i+1},\qquad
 b_a=\prod_i\binom{\kappa_i}{a_i},\qquad
 f_\psi(z)=\sum_{a\leq\kappa}\sqrt{b_a}\psi_a z^a,
\]
and define
\begin{equation}\label{dynspin:kernel}
 F_\rho(z,w)=\sum_{a,b\leq\kappa}\sqrt{b_ab_b}\rho_{ab}z^aw^b.
\end{equation}
A density matrix is coherently Lee--Yang when this polynomial is nonzero
on $\D^{2n}$, with independent row and column variables.  Ordinary
preservation refers to this class on $\mathcal H$.  Complete preservation
means preservation after adjoining any finite collection of inert finite
spins, with the same encoding.  A pure input means a rank-one density
matrix with disk-stable vector polynomial; the vector need not be a
spin-coherent state.  Adjoints are always taken in the physical
orthonormal basis, equivalently in the polynomial inner product
$\langle z^a,z^b\rangle=\delta_{ab}/b_a$.

\begin{theorem}[Finite-spin physical generators]\label{dyn:arbitrary-spin}
Suppose that $\kappa_i\geq2$ for at least one site, and let
$\Phi_t=e^{t\mathcal L}$ be a continuous CPTP semigroup on $\mathcal H$.
The following conditions are equivalent.
\begin{enumerate}
\item $\Phi_t$ ordinarily preserves coherently Lee--Yang densities for
      every $t\geq0$.
\item For some $\varepsilon>0$, $\Phi_t$ preserves every pure coherently
      Lee--Yang input for $0\leq t<\varepsilon$.
\item $\Phi_t$ completely preserves coherently Lee--Yang densities for
      every $t\geq0$.
\item The generator is a sum of local generators,
\begin{equation}\label{dynspin:locality}
 \mathcal L=\sum_i\mathcal L_i\otimes\Id_{\widehat i}.
\end{equation}
At a site with $j_i\geq1$, its local term is
\begin{equation}\label{dynspin:isotropic}
 \mathcal L_i(\rho)=-i[h_iS_i^z,\rho]
 -\frac{\gamma_i}{2}\sum_{\alpha=x,y,z}
       [S_i^\alpha,[S_i^\alpha,\rho]],
 \qquad h_i\in\R,\quad\gamma_i\geq0.
\end{equation}
At a qubit site, its local term is
\[
 \mathcal L_i(\rho)=-i[h_iZ,\rho]
 +\sum_{\alpha,\beta=x,y,z}(K_i)_{\alpha\beta}
 \left(\sigma_\alpha\rho\sigma_\beta
       -\tfrac12\{\sigma_\beta\sigma_\alpha,\rho\}\right),
 \qquad h_i\in\R,
\]
where $B_i=B_i^{\mathsf T}\in\R^{2\times2}$, $d_i\in\R^2$, and
\begin{equation}\label{dynspin:qubit-cone}
 K_i=\begin{pmatrix}B_i&i d_i\\-i d_i^{\mathsf T}&c_i\end{pmatrix}
 \succeq0,\qquad
 \begin{pmatrix}
 c_iI_2-\operatorname{adj}B_i&d_i\\
 d_i^{\mathsf T}&\Tr B_i
 \end{pmatrix}\succeq0.
\end{equation}
\end{enumerate}
In particular, $\Phi_t=\bigotimes_i e^{t\mathcal L_i}$.
\end{theorem}

The dissipator in \eqref{dynspin:isotropic} is the familiar isotropic
spin diffusion, with Casimir multipole rates
\cite{DenisMartin22,FilippovKuzhamuratova19}.  The assertion here is its
necessity from physical zero-freeness, without an assumption of covariance,
linearity of the noise, locality, or unitality.  The axial Hamiltonian
commutes with the isotropic dissipator; a nonzero axial drift need not
make the full semigroup rotation covariant.

We first establish the reduction from arbitrary GKLS operators using
only physical pure inputs.  The spin realization is
\begin{equation}\label{dynspin:spin-action}
 S_i^x=\tfrac12[(1-z_i^2)\partial_i+\kappa_i z_i],\quad
 S_i^y=\tfrac1{2i}[(1+z_i^2)\partial_i-\kappa_i z_i],\quad
 S_i^z=\kappa_i/2-z_i\partial_i.
\end{equation}
The capacity-dependent Cayley transform is
\begin{equation}\label{dynspin:cayley}
 (Cf)(s)=2^{-|\kappa|/2}\prod_i(s_i+i)^{\kappa_i}
 f\left(\frac{s_1-i}{s_1+i},\ldots,
         \frac{s_n-i}{s_n+i}\right).
\end{equation}
It is the symmetric-power realization of \eqref{dyn:cayley}, and is
unitary in the weighted inner product.  Applied separately to both
coefficient legs of a pure density, it gives
\begin{equation}\label{dynspin:reflection}
 p(z)\mathscr Rp(w),\qquad
 \mathscr Rp(w)=(-1)^{|\kappa|}\overline{p(-\bar w)},\qquad p=Cf_\psi.
\end{equation}
Indeed this follows in capacity one from
$\overline C=\diag(-1,1)C$, and then by symmetric powers and tensor
products.  This two-leg transformation is not unitary conjugation of
the density matrix.  Every stable nonzero $p$ in the capacity box is
obtained from a physical pure input after normalization.  Its degree
may be strictly below the capacity; the inverse of
\eqref{dynspin:cayley} still belongs to the prescribed box.

\begin{lemma}[Physical order reduction]\label{dynspin:order}
Under condition (2) of Theorem~\ref{dyn:arbitrary-spin}, every jump in any
finite GKLS representation has, after removal of its scalar part, the form
\[
 V_\nu\in\operatorname{span}_{\C}\{S_i^x,S_i^y,S_i^z:1\leq i\leq n\}.
\]
The Hamiltonian has spin degree at most two: it is a sum of a scalar,
linear spin terms, on-site symmetric quadratics, and two-site bilinears.
At qubit sites the on-site quadratics reduce to a scalar and linear terms.
\end{lemma}

\begin{proof}
Write a finite-dimensional GKLS representation
\[
 \mathcal L(\rho)=-i[H,\rho]
 +\sum_\nu\left(V_\nu\rho V_\nu^*
                  -\tfrac12\{V_\nu^*V_\nu,\rho\}\right).
\]
In the vector-polynomial coordinates, let
\[
 Q_\nu=CV_\nu C^{-1}
       =\sum_{\alpha\leq\kappa}P_{\nu,\alpha}(z)\partial^\alpha
\]
be its canonical representation.  Let $R$ be the largest order among
all these operators.  Suppose $R\geq2$, and fix an index
$|\alpha|=R$.  For arbitrary $a\in\R^n$ take the stable pure input
$p_a(z)=\prod_i(z_i-a_i)^{\kappa_i}$.  In local coordinates
$u=z-a$, $v=w+a$, its doubled polynomial is exactly $u^\kappa v^\kappa$,
up to a positive normalization constant.  Put $M=2|\kappa|$.

Restrict the transformed output to $(u,v)=\tau\xi$ for any
$\xi\in\R_{>0}^{2n}$.  The initial polynomial is a nonzero multiple
of $\tau^M$, and the output remains upper-half-plane stable near time
zero.  Lemma~\ref{gen:cluster-moments} makes every coefficient of its
first variation below degree $M-2$ vanish.  Varying $\xi$ over the
positive orthant shows, by polynomial identity, that each multivariate
homogeneous part below that degree vanishes as well.

Consider the coefficient of $u^{\kappa-\alpha}v^{\kappa-\alpha}$.
Hamiltonian and anticommutator terms act on one leg only, leaving the
opposite full root factor untouched, and do not contribute.  The
sandwich term from $V_\nu$ is
$(Q_\nu p_a)(z)\mathscr R(Q_\nu p_a)(w)$.  To obtain
$u^{\kappa-\alpha}$ from a canonical term with derivative index $\beta$,
even allowing Taylor terms of its coefficient polynomial at $a$, one
must have $\beta\geq\alpha$.  Since $|\beta|\leq R=|\alpha|$, this
requires $\beta=\alpha$ and the constant Taylor coefficient.  Reflection
therefore gives the total coefficient
\begin{equation}\label{dynspin:squares}
 (-1)^R(\kappa)_\alpha^2
       \sum_\nu|P_{\nu,\alpha}(a)|^2,
 \qquad
 (\kappa)_\alpha=\prod_i\frac{\kappa_i!}{(\kappa_i-\alpha_i)!}.
\end{equation}
Its degree is $M-2R<M-2$, so it is zero.  Every square vanishes for
every real $a$.  Thus every coefficient of maximal order vanishes
identically, contradicting the definition of $R$.

All jumps have order at most one.  Proposition~\ref{gen:box-decomposition}
identifies this space as a scalar plus the local $\mathfrak{sl}_2$
actions.  Formula~\eqref{dynspin:spin-action} and the Cayley rotation
identify it with the stated spin span.  Absorb scalar jump parts into
the Hamiltonian.  The dissipative superoperator now has differential
order at most two on the doubled box.

If $CHC^{-1}$ had maximal canonical order $R\geq3$, use the same input
and choose $|\alpha|=R$.  The coefficient of
$u^{\kappa-\alpha}v^\kappa$ in the first variation can come only from
the row Hamiltonian and its constant maximal-order coefficient at $a$.
The column Hamiltonian retains $u^\kappa$, and the dissipator cannot
lower total degree by $R$.  Its degree $M-R<M-2$ forces that Hamiltonian
coefficient to vanish for every $a$, again a contradiction.
Proposition~\ref{gen:box-decomposition} identifies the resulting
order-two space with products of at most two local spin operators.
Hermiticity gives real coefficients in the Hermitian basis consisting
of scalars, local spins, local symmetric quadrupoles, and products on
distinct sites.  The antisymmetric local products are spin commutators,
and the scalar symmetric component is the Casimir.  At capacity one
the remaining quadrupoles vanish.  This proves the assertion.
\end{proof}

\begin{lemma}[Locality at finite spins]\label{dynspin:local}
Under the same hypothesis, the generator is a sum of local GKLS
generators as in \eqref{dynspin:locality}.
\end{lemma}

\begin{proof}
Put $T_i^\alpha=2S_i^\alpha$.  Lemma~\ref{dynspin:order} gives a PSD
Gram matrix $K$ in the basis of these local spin operators.  Its
$(i,j)$ block is denoted $K^{ij}$.  The transformed left and right
actions have respective derivative vectors
\begin{equation}\label{dynspin:vectors}
 \mathbf v(s)=(2s,1-s^2,i(1+s^2)),\qquad
 \widetilde{\mathbf v}(t)=(2t,t^2-1,i(1+t^2)).
\end{equation}
These vectors are independent of the capacity.  Only their zeroth-order
coefficients depend on $\kappa_i$.

Fix distinct sites $i,j$.  The physical pure product with transformed
vector $\prod_k(z_k-a_k)$ allows independent row--row contacts at
$z_i=s,z_j=t$ and row--column contacts at $z_i=s,w_j=t$, by choosing
$a_i=s$ and respectively $a_j=t$ or $a_j=-t$.  Keep all other legs
off their roots and restrict to a line with all slopes positive.
The initial root has multiplicity two.  Only the chosen mixed
principal coefficient survives at the contact; after division by the
initial leading coefficient its multiplier is the reciprocal of a
positive product of slopes.  Lemma~\ref{gen:cluster-moments} makes it
real and nonpositive for every real $s,t$.  On-site quadrupoles and
other terms leave at least one selected root untouched.

Write the pair Hamiltonian as
$\sum_{\alpha,\beta}h_{\alpha\beta}T_i^\alpha T_j^\beta$, with $h$
real.  The row--column coefficient is
$\mathbf v(s)^{\mathsf T}K^{ij}\widetilde{\mathbf v}(t)$.
The row--row coefficient is
$\mathbf v(s)^{\mathsf T}(-\Re K^{ij}-ih)\mathbf v(t)$: the conjugate Gram blocks
combine in the anticommutator, since the two sites commute.
Set
\[
 u(s)=\frac{(2s,1-s^2)}{1+s^2},\qquad R_0=\diag(1,-1).
\]
Its range is dense in the unit circle.  The normalized derivative
vectors are $(u(s),i)$ and $(R_0u(t),i)$.

Reality of the row--column coefficient, by comparison of its constant,
linear, and bilinear circle components, gives
\[
 K^{ij}=\begin{pmatrix}A&i b\\i d^{\mathsf T}&e\end{pmatrix},
 \qquad A,b,d,e\text{ real}.
\]
Its sign condition is
\begin{equation}\label{dynspin:cross-lr}
 u^{\mathsf T}AR_0v-b\cdot u-d\cdot R_0v-e\leq0
 \qquad (u,v\in S^1).
\end{equation}
For completeness, the comparison follows by averaging the imaginary
part first alone, then against $u$, $R_0v$, and $u(R_0v)^{\mathsf T}$;
the circle moments are $\mathbb E u=0$ and
$\mathbb E uu^{\mathsf T}=I_2/2$.
Reality of the row--row coefficient similarly gives $h_{\perp\perp}=0$
and $h_{zz}=0$.  With $k=h_{\perp z}$ and
$\ell^{\mathsf T}=h_{z\perp}$, its remaining sign condition is
\begin{equation}\label{dynspin:cross-ll}
 e-u^{\mathsf T}Av+k\cdot u+\ell\cdot v\leq0.
\end{equation}
Independent circle averages of \eqref{dynspin:cross-lr} and
\eqref{dynspin:cross-ll} are $-e$ and $e$, so $e=0$.
Each continuous nonpositive function then has zero average and vanishes
everywhere.  The same circle moments give $A=b=d=k=\ell=0$.
Thus all cross Gram blocks and pair Hamiltonians vanish.  The diagonal
Gram blocks are PSD and define local GKLS generators, proving locality.
\end{proof}

\begin{proof}[Proof of Theorem~\ref{dyn:arbitrary-spin}]
We prove (2)$\Rightarrow$(4) first.  The preceding lemmas give locality.
Each local semigroup preserves its pure Lee--Yang inputs near zero:
pad such an input with arbitrary pure Lee--Yang states at the other
sites.  Its evolved global polynomial is a product of the local output
polynomials, each nonzero since its density has trace one.  Stability
of that product implies stability of each factor.

Fix a site of capacity $\kappa\geq2$, retaining the normalization
$T=2S$.  Let $q(s,t)$ denote its row--column principal coefficient.
Use the pure vector polynomial $p(z)=(z-s)(z+t)$, with $s+t\ne0$.
At row coordinate $s$ and column coordinate $t$ both factors in
\eqref{dynspin:reflection} have simple roots.  A positive-slope line
has a double root, and only $q(s,t)$ contributes at the contact.
The common reflection sign cancels on dividing by the initial leading
coefficient.  Thus $q(s,t)$ is real and nonpositive; continuity adds
the omitted line $s+t=0$.

For a general Hermitian Gram matrix write
\[
 K=\begin{pmatrix}
 a&u+i\eta&r+i\mu\\
 u-i\eta&b&v+i\nu\\
 r-i\mu&v-i\nu&c
 \end{pmatrix}.
\]
Multiplication of \eqref{dynspin:vectors} gives
\[
 \Im q(s,t)=2(s+t)[-\eta(1-st)+r(1+st)+v(t-s)].
\]
Hence $\eta=r=v=0$, and
\begin{equation}\label{dynspin:local-shape}
 K=\begin{pmatrix}B&i d\\-i d^{\mathsf T}&c\end{pmatrix},
 \qquad B=B^{\mathsf T}\in\R^{2\times2},\quad d\in\R^2.
\end{equation}
With
\[
 r(s)=\frac{(2s,1-s^2)}{1+s^2},\qquad
 v(t)=\frac{(2t,t^2-1)}{1+t^2},
\]
the mixed coefficient satisfies
\begin{equation}\label{dynspin:circle}
 \frac{q(s,t)}{(1+s^2)(1+t^2)}
 =r(s)^{\mathsf T}Bv(t)+d\cdot(v(t)-r(s))-c\leq0.
\end{equation}
Both circle variables range densely and independently.  Taking them
equal gives $B\preceq cI_2$.

There is a second, independent contact available at this capacity.
Take $p(z)=(z-s)^2$ and a column coordinate away from its reflected
root.  Its pure-row principal coefficient must be real and nonpositive.
This follows either by real specialization of the column, or by a
positive-slope line through the contact; the nonzero complementary
factor cancels from the root-cluster coefficient.
Write the quadratic Hamiltonian as $H_2=T^{\mathsf T}QT$, with $Q$
real symmetric.  Put $u=r(s)$.  The pure-row principal coefficient,
divided by $(1+s^2)^2$, is
\begin{equation}\label{dynspin:row-quadratic}
 \tfrac12(c-u^{\mathsf T}Bu)+2Q_{\perp z}\cdot u
       -i(u^{\mathsf T}Q_{\perp\perp}u-Q_{zz}).
\end{equation}
Indeed its dissipative part is
$-\mathbf v(s)^{\mathsf T}(\Re K)\mathbf v(s)/2$, and its Hamiltonian part is
$-i\mathbf v(s)^{\mathsf T}Q\mathbf v(s)$, using the three-component vector in
\eqref{dynspin:vectors}.  Reality gives
$Q_{\perp\perp}=Q_{zz}I_2$.  Remove the corresponding Casimir scalar
and set $k=Q_{\perp z}$.  Nonpositivity at $u$ and $-u$ now gives
$c-u^{\mathsf T}Bu\leq0$.  Together with $B\preceq cI_2$, this yields
$B=cI_2$; the remaining inequality $2k\cdot u\leq0$ gives $k=0$.
The quadratic Hamiltonian is consequently scalar.

Equation~\eqref{dynspin:circle} reduces to
\[
 c(r\cdot v-1)+d\cdot(v-r)\leq0\qquad(r,v\in S^1).
\]
If $d\ne0$, take $r\perp d$ and
$v=(\cos\theta)r+(\sin\theta)d/|d|$ for small positive $\theta$.
The left side is $c(\cos\theta-1)+|d|\sin\theta>0$, a contradiction.
Thus $d=0$, $K=cI_3$, and positivity gives $c\geq0$.

The transformed isotropic dissipator has the explicit form
\begin{align}\label{dynspin:cayley-diffusion}
 \widetilde{\mathcal D}={}&-2c\kappa(\kappa zw+1)
 +2c\kappa z(zw-1)\partial_z
 +2c\kappa w(zw-1)\partial_w\notag\\
 &-2c(zw-1)^2\partial_z\partial_w.
\end{align}
To verify it directly, the left actions of $(T_x,T_y,T_z)$ are
\[
 2z\partial_z-\kappa,\quad
 (1-z^2)\partial_z+\kappa z,\quad
 i[(1+z^2)\partial_z-\kappa z],
\]
and the right actions are the same formulas in $w$ except that the
middle action has its sign reversed.  Substitute them into
$c\sum_\alpha(L_\alpha R_\alpha
 -(L_\alpha^2+R_\alpha^2)/2)$.
In particular this is a real operator and its pure-row and pure-column
second-order coefficients vanish.
For the remaining Hamiltonian $H=\widehat h_xT_x+\widehat h_yT_y+
\widehat h_zT_z$, use $p(z)=z-s$ and keep the reflected column off its
root.  The imaginary simple-root velocity comes only from the row
Hamiltonian and is
\[
 \widehat h_y s^2-2\widehat h_xs-\widehat h_y\geq0\qquad(s\in\R),
\]
by Lemma~\ref{gen:cluster-moments}.  Hence
$\widehat h_x=\widehat h_y=0$.  Returning to $S=T/2$ gives
\eqref{dynspin:isotropic}, with $h=2\widehat h_z$ and $\gamma=4c$.

It remains to recover the full qubit condition in a mixed-spin system.
Fix a qubit and a high-spin partner whose local form has just been
proved.  Their mixed principal coefficients will be denoted by
$q_{\rm q}(s,t)$ and $q_{\rm h}(u,v)=-2c(uv-1)^2$.
For arbitrary real $s,t$ with $s+t\ne0$, choose
\begin{equation}\label{dynspin:mask-parameters}
 \begin{gathered}
 x=(s+t)/2,\qquad a=(s-t)/2,\qquad c_0>|x|,\\
 b=\sqrt{(c_0/x)^2-1},\qquad
 u_*=b+c_0/x,\qquad v_*=-b+c_0/x.
 \end{gathered}
\end{equation}
Then $u_*v_*=1$.  The polynomial
\begin{equation}\label{dynspin:mixed-test}
 p(z,u)=(z-a)(u-b)-c_0
\end{equation}
is upper-half-plane stable, since for $z\in\HH$ its root in $u$ is
$b+c_0/(z-a)$, which lies in the lower half-plane.
It is a valid physical vector in capacities $(1,\kappa)$, including
when its degree is below $\kappa$.  Apart from the common parity sign,
its reflected factor is $(w+a)(v+b)-c_0$.

At $(z,u,w,v)=(s,u_*,t,v_*)$ both factors vanish.  Along the line with
all four slopes one, the initial quadratic coefficient, apart from
the same common sign, is $(c_0/x+x)^2>0$.  Locality makes its first
variation at the contact equal, with that sign, to
\begin{equation}\label{dynspin:mask-isolation}
 \frac{c_0^2}{x^2}q_{\rm q}(s,t)
       +x^2q_{\rm h}(u_*,v_*)
 =\frac{c_0^2}{x^2}q_{\rm q}(s,t).
\end{equation}
All lower-order terms vanish because one input factor remains zero;
the high-spin pure-row and pure-column coefficients are zero by
\eqref{dynspin:cayley-diffusion}.  Other sites can be padded by pure
stable factors kept nonzero at the contact.  The root-cluster condition
thus makes $q_{\rm q}(s,t)$ real and nonpositive.  Continuity again adds
$s+t=0$.  Consequently \eqref{dynspin:local-shape} and
\eqref{dynspin:circle} hold at every qubit as well.  For this Gram shape
the dissipator is real after Cayley transformation, as in the proof of
Proposition~\ref{dyn:local-complete}.  The same simple-root test excludes
its transverse Hamiltonian.  Lemma~\ref{dyn:lmi-reduction} gives precisely
\eqref{dynspin:qubit-cone}, including $\Tr B=0$.  This proves
(2)$\Rightarrow$(4), without treating the evolving partner as inert.

Conversely assume (4).  Formula~\eqref{dynspin:cayley-diffusion} is a
real endomorphism of the doubled capacity box whose only second-order
coefficient is $-2c(zw-1)^2\leq0$.  The axial Hamiltonian contributes
only a real first-order endomorphism.  At qubit sites, the proof of
Proposition~\ref{dyn:local-complete} gives a real box endomorphism with
its only mixed coefficient nonpositive by
Lemma~\ref{dyn:lmi-reduction}.  The sum therefore satisfies
Theorem~\ref{gen:main-real} on the full doubled capacity box.  Its
exponential preserves all complex stable polynomials after arbitrary
inert-variable extensions.  Undoing the coefficient Cayley transforms
gives the asserted physical complete preservation, in particular for
arbitrary inert finite spins.  No convexity of the stable class is used.
This proves (4)$\Rightarrow$(3).  Finally (3)$\Rightarrow$(1) and
(1)$\Rightarrow$(2) are immediate.
\end{proof}

The all-qubit ordinary classification is Theorem~\ref{dyn:main}; it is
not asserted there or here to follow from pure inputs alone.  On a single
qubit the ordinary cone is instead Proposition~\ref{dyn:one-qubit}, and
ordinary preservation need not be complete.  Theorem~\ref{dyn:arbitrary-spin}
also makes no classification claim for isolated noncovariant high-spin
channels: its semigroup hypothesis is essential to the argument.

\subsection{An application to rotation convolution semigroups}
\label{dynspin:brownian-subsection}

Gaussian processes on compact groups can already be recognized through
a faithful finite-dimensional representation and its tensor products;
see \cite[Theorem~2.8]{Voit03}, which restates the earlier compact-group
characterization.  For the real defining representation of $SO(3)$,
the spin-one conjugation action is its tensor square and contains the
defining representation as its antisymmetric part.  The following
application combines that recognition mechanism with the physical
classification above.  We give the short character proof in this case
using Hunt's formula \cite{Hunt56} and its twirling realization
\cite[Theorem~5.1]{Aniello10}.

\begin{corollary}[Lee--Yang rotation dynamics]\label{dyn:brownian-recognition}
Let $(\mu_t)_{t\geq0}$ be a weakly continuous probability convolution
semigroup on $SO(3)$ with $\mu_0=\delta_e$.  Let $U$ be its spin-one
representation and set
\[
 \Phi_t(\rho)=\int_{SO(3)}U(g)\rho U(g)^*\,\mu_t(dg).
\]
Then $\Phi_t$ preserves every pure coherently Lee--Yang input in a
neighborhood of zero if and only if the group generator is
\begin{equation}\label{dynspin:group-generator}
 \mathcal A=hX_z+\frac\gamma2(X_x^2+X_y^2+X_z^2),
 \qquad h\in\R,\quad\gamma\geq0.
\end{equation}
Here $X_\alpha$ are the left-invariant fields normalized by
$U(\exp(tE_\alpha))=e^{-itS_\alpha}$ and
$X_\alpha f(g)=\left.\frac{d}{dt}f(g\exp(tE_\alpha))\right|_{t=0}$.
In particular the L\'evy measure is zero, without a prior assumption of
continuous sample paths or rotation invariance of $\mu_t$.
\end{corollary}

\begin{proof}
The twirling maps form a continuous CPTP semigroup.  Under the stated
preservation hypothesis Theorem~\ref{dyn:arbitrary-spin} gives
\eqref{dynspin:isotropic}.  The spin-one operator space decomposes as
the angular-momentum representations $0\oplus1\oplus2$.  On its
$\ell$ component the double-commutator Casimir is $\ell(\ell+1)I$;
the axial commutator has trace zero.  Thus the generator restrictions
satisfy
\[
 \Tr\mathcal L_1=-3\gamma,\qquad \Tr\mathcal L_2=-15\gamma.
\]
Let $\chi_1,\chi_2$ be the characters of the three- and five-dimensional
representations.  For a rotation of angle $\theta$, put
\begin{equation}\label{dynspin:character-witness}
 F(g)=2-\chi_1(g)+\tfrac15\chi_2(g)
     =\tfrac45(1-\cos\theta)^2.
\end{equation}
This smooth function is strictly positive off the identity and has
zero jet through order two at the identity.  Since
$\int\chi_\ell\,d\mu_t=\Tr(e^{t\mathcal L_\ell})$, the displayed
traces give $\mathcal AF(e)=3\gamma-15\gamma/5=0$.

Write Hunt's generator with drift, diffusion matrix $a\succeq0$, and
L\'evy measure $\nu$, taking its second-order part to be
$\tfrac12\sum_{\alpha,\beta}a_{\alpha\beta}X_\alpha X_\beta$.
At the identity all local terms and the linear jump compensation
vanish on $F$.  Therefore
\[
 0=\mathcal AF(e)=\int_{SO(3)\setminus\{e\}}F(g)\,\nu(dg).
\]
The integral is finite: near the identity $F=O(\dist(g,e)^4)$, while
a L\'evy measure has the requisite local second-moment integrability;
away from the identity compactness applies.  Positivity of $F$ off
$e$ implies $\nu=0$, including for an initially infinite-activity
L\'evy measure.

On the dipole space the symmetric part of the remaining generator is
$\tfrac12(a-(\Tr a)I_3)$.  Comparing with $-\gamma I_3$, first taking
traces and then substituting back, gives $a=\gamma I_3$.  Its
antisymmetric part identifies the drift as $hX_z$.  This proves
\eqref{dynspin:group-generator}.  Conversely, that group generator
induces \eqref{dynspin:isotropic} under twirling, so
Theorem~\ref{dyn:arbitrary-spin} gives complete preservation.
\end{proof}

The group and the tested spin matter.  Conjugation cannot detect jumps
in the center $\{\pm I\}$ of $SU(2)$, so the zero-L\'evy-measure
conclusion is stated on $SO(3)$.  Moreover, Poisson jumps of rate
$\lambda$ by angle $\pi$ about the $z$ axis induce on a qubit
\[
 \Phi_t(\rho)=\tfrac12(1+e^{-2\lambda t})\rho
              +\tfrac12(1-e^{-2\lambda t})Z\rho Z.
\]
This is an allowed dephasing semigroup, but its rotation process has
nonzero jumps.  Thus spin one cannot be replaced uniformly by spin
one-half in Corollary~\ref{dyn:brownian-recognition}.


\section{Static channels and stable encodings}
\label{stat:section}

This section classifies invertible Hermiticity-preserving, trace-preserving
maps with Lee--Yang full Choi tensors, and gives structural results for
singular channels. The
variables associated with different physical qubits are always distinct.
In particular, permutation of physical wires is allowed, whereas an arbitrary
change of the computational basis is not.

\subsection{Coefficient tensors and complete preservation}

For a matrix \(M:(\C^2)^{\otimes m}\to(\C^2)^{\otimes n}\), put
\begin{equation}
 G_M(z,w)=\sum_{a\in\{0,1\}^n}\sum_{b\in\{0,1\}^m}
 M_{ab}z^aw^b.
 \label{stat:matrix-polynomial}
\end{equation}
We call \(M\) \emph{tensor-LY} if this polynomial is nonzero whenever all
its variables belong to the open unit disk \(\D\).  Zero is excluded from
this definition.  For a complex-linear map
\(\Phi:M_{2^m}(\C)\to M_{2^n}(\C)\), our Choi convention is
\begin{equation}
 J(\Phi)=\sum_{c,d}\Phi(|c\rangle\langle d|)
                 \otimes |c\rangle\langle d|.
 \label{stat:choi}
\end{equation}
Thus full Choi--LY means stability in all \(2(m+n)\) variables of
\(G_{J(\Phi)}\).  Hermiticity-preserving and trace-preserving maps are
abbreviated HP and TP, respectively.

We use the disk tensor-contraction theorem in the following form: contracting
two binary indices of a tensor-LY polynomial with unit radii produces a
tensor-LY polynomial or zero.  Tensor products, permutation of indices, and
coefficientwise conjugation preserve tensor-LY.  These facts follow from
the classical Asano contraction theorem; the tensor and channel formulation
is described in \cite[Section~2]{WBG26}.  The possible zero result at unit
radius will be excluded explicitly where necessary.

\begin{lemma}[The Choi bridge]
\label{stat:bridge}
Let \(\Phi:M_{2^m}(\C)\to M_{2^n}(\C)\) be complex linear.
If \(J(\Phi)\) is tensor-LY, then \(\Phi\), and each of its extensions
by identity maps, sends tensor-LY matrices to tensor-LY matrices or zero.
Its Hilbert--Schmidt adjoint \(\Phi^\dagger\) is also full Choi--LY.
For a CPTP map, full Choi--LY is equivalent to preservation of tensor-LY
density matrices with every finite qubit reference system left inert.
\end{lemma}

\begin{proof}
The superoperator and Choi arrays satisfy
\[
 S(\Phi)_{(a,b),(c,d)}=\Phi(|c\rangle\langle d|)_{ab},
 \qquad
 J(\Phi)_{(a,c),(b,d)}=S(\Phi)_{(a,b),(c,d)}.
\]
They therefore differ only by a permutation of tensor indices.  Applying
\(\Phi\) is contraction of this array with the input coefficient array.
Input transposes occurring in the matrix convention also only permute
indices.  The contraction theorem proves the first assertion, including
inert indices.  The superoperator of \(\Phi^\dagger\) is obtained by
conjugating entries and interchanging the input and output index pairs,
which proves the adjoint assertion.

For a CPTP map the output of a density matrix has trace one, so cannot
be the zero alternative.  Conversely, take the normalized maximally
entangled vector between the input and an equally large reference:
\[
 |\Omega_m\rangle=2^{-m/2}\sum_c|c\rangle|c\rangle.
\]
Its amplitude polynomial is
\(2^{-m/2}\prod_i(1+x_i u_i)\), and its density-matrix polynomial is
the product of this polynomial and its coefficientwise conjugate in
independent variables.  Both are stable.  Applying \(\Phi\) to one
side produces \(2^{-m}J(\Phi)\), proving the converse.
\end{proof}

Ordinary multiplication of tensor-LY matrices is likewise a tensor
contraction, with the same zero alternative.  In particular, multiplication
on either side by an invertible tensor-LY matrix whose inverse is
tensor-LY preserves the class of tensor-LY matrices.

\subsection{An elementary matrix edge and hyperbolic pencils}

Write
\[
 X=\begin{pmatrix}0&1\\1&0\end{pmatrix},\qquad
 Y=\begin{pmatrix}0&-i\\i&0\end{pmatrix},\qquad
 Z=\begin{pmatrix}1&0\\0&-1\end{pmatrix}.
\]
Subscripts denote physical sites, with identity factors suppressed.  The
real vector spaces needed below are
\begin{equation}
 \mathfrak g=\operatorname{span}_{\R}\{X,Y,iZ\},\qquad
 \mathcal E=\R I+\mathfrak g,
 \qquad
 \mathfrak e_n=\C I+\bigoplus_{j=1}^n\mathfrak g_j.
 \label{stat:edge-spaces}
\end{equation}

\begin{lemma}[Two-sided tangency]
\label{stat:edge}
Suppose \(M(t)\) is a matrix-valued curve defined for real \(t\) near
zero, differentiable at zero, with \(M(0)=I\) and \(M'(0)=B\).
If \(M(t)\) is tensor-LY for both signs of all sufficiently small
\(t\), then \(B\in\mathfrak e_n\).
Conversely, \(B\in\mathfrak e_n\) implies that \(e^{tB}\) is
tensor-LY for every real \(t\).
\end{lemma}

\begin{proof}
Fix a site \(j\) and \(|\zeta|=1\).  Specializing \(z_j=\zeta\)
is justified by first using \(r\zeta\), \(r<1\), and then applying
Hurwitz's theorem.  The specialization is stable or zero.  At \(t=0\)
its value with all remaining variables zero is one, so it is not the
zero polynomial for sufficiently small \(|t|\).
Fix interior values \(u=(z_k,w_k)_{k\ne j}\), and put
\(D(u)=\prod_{k\ne j}(1+z_kw_k)\).  The resulting affine polynomial
in \(w_j\) is
\[
 D(u)(1+\zeta w_j)+tG_B(\zeta,w_j,u)+o(t).
\]
Its root \(r(t)\) is differentiable at zero, with
\(r(0)=-\zeta^{-1}\), and stability implies \(|r(t)|\ge1\) for
both signs of small \(t\).  Differentiating the root equation gives
\[
 \frac{r'(0)}{r(0)}
 =\frac{G_B(\zeta,-\zeta^{-1},u)}{D(u)}.
\]
The derivative of \(|r(t)|^2\) at its two-sided minimum is zero.
Consequently the holomorphic function on the remaining polydisk
\begin{equation}
 u\longmapsto\frac{G_B(\zeta,-\zeta^{-1},u)}{D(u)}
 \label{stat:contact-quotient}
\end{equation}
has purely imaginary values, and is therefore a purely imaginary
constant.

Expand \(B\) in the Pauli tensor basis, separating site \(j\):
\(B=\sum_{\sigma,\tau}b_{\sigma,\tau}\sigma_j\otimes\tau\).
At \((z_j,w_j)=(\zeta,-\zeta^{-1})\), the one-site polynomials are
\[
 G_I=0,\qquad G_X=\zeta-\zeta^{-1},\qquad
 G_Y=i(\zeta+\zeta^{-1}),\qquad G_Z=2.
\]
Since the Pauli polynomials on the other sites are linearly independent,
constancy of \eqref{stat:contact-quotient} gives, for every nonidentity
word \(\tau\),
\[
 b_{X,\tau}(\zeta-\zeta^{-1})
 +ib_{Y,\tau}(\zeta+\zeta^{-1})+2b_{Z,\tau}=0.
\]
Taking \(\zeta=1,-1,i\) forces all three coefficients to vanish.
As \(j\) was arbitrary, every Pauli coefficient of weight at least
two is zero.  Thus
\(B=cI+\sum_j(x_jX_j+y_jY_j+d_jZ_j)\).
The remaining scalar contact expression is
\[
 x_j(\zeta-\zeta^{-1})+iy_j(\zeta+\zeta^{-1})+2d_j.
\]
Its real part is zero.  The same three values of \(\zeta\) imply
\(x_j,y_j\in\R\) and \(d_j\in i\R\), as asserted.  The scalar
\(c\) is unrestricted.

For the converse, if
\(N=aI+xX+yY+izZ\in\mathcal E\), direct multiplication gives
\begin{equation}
 NZN^*=\delta Z,
 \qquad \delta=a^2+z^2-x^2-y^2=\det N.
 \label{stat:unit-identity}
\end{equation}
When \(\delta>0\), write
\(G_N(s,w)=A(s)+C(s)w\).  Equation~\eqref{stat:unit-identity}
implies
\[
 |A(s)|^2-|C(s)|^2=\delta(1-|s|^2)>0\qquad(|s|<1),
\]
so \(N\) is tensor-LY.  Its inverse has the same property.
For \(b=xX+yY+izZ\), one has \(b^2=(x^2+y^2-z^2)I\).
The exponential \(e^{tb}\) belongs to \(\mathcal E\) for real
\(t\) and has determinant one.  Therefore it is tensor-LY.  Finally,
for \(B=cI+\sum_jb_j\),
\(e^{tB}=e^{tc}\bigotimes_j e^{tb_j}\), which proves the assertion.
\end{proof}

The space \(\mathcal E\) is closed under multiplication.  Its members
sufficiently close to \(I\), and their inverses, are tensor-LY by
\eqref{stat:unit-identity}.  We call these matrices local stability
units.  A traceless element \(b=xX+yY+izZ\in\mathfrak g\) is
\emph{hyperbolic} if \(x^2+y^2-z^2>0\).  A diagonal phase
conjugation followed by conjugation with \(e^{rY}\), where
\(\tanh(2r)=z/\sqrt{x^2+y^2}\), carries it to
\(\sqrt{x^2+y^2-z^2}\,X\).  The conjugating matrices and their
inverses are local stability units.

\begin{lemma}[Hyperbolic affine pencil]
\label{stat:hyperbolic}
Let \(B=cI+\sum_jb_j\in\mathfrak e_n\), and suppose one
\(b_k\) is hyperbolic.  If \(I+tB\) is tensor-LY for some nonzero
real \(t\), then \(b_j=0\) for all \(j\ne k\) and \(c\in\R\).
\end{lemma}

\begin{proof}
Conjugating by a local stability unit, assume \(b_k=aX_k\),
\(a>0\).  For a one-site matrix \(b\) define
\(q_b(z,w)=G_b(z,w)/(1+zw)\).  The function
\[
 q_X(z,z)=\frac{2z}{1+z^2}
\]
attains every nonreal complex number for \(|z|<1\).  Indeed, the
inverse equation has two reciprocal roots; neither is on the unit
circle for a nonreal target, because the displayed quotient is real
there wherever finite.  Exactly one inverse root is therefore in
the disk.

If \(b_j\ne0\) for another site, \(q_{b_j}\) is a nonconstant
holomorphic function, since \(b_j\) is nonzero and traceless.
Freeze all other variable pairs at zero, and divide the polynomial
of \(I+tB\) by \(\prod_i(1+z_iw_i)\).  The result has the form
\(C+taq_X(z_k,w_k)+tq_{b_j}(z_j,w_j)\).  The open image of
\(q_{b_j}\) allows its variables to be chosen so that the target
needed from \(q_X\) is nonreal.  The preceding inverse construction
then gives an interior zero, a contradiction.  Once the other
components vanish, an imaginary part of \(c\) would make
\(-(1+tc)/(ta)\) nonreal and give the same contradiction.
\end{proof}

\subsection{Routing and factorization of visible outputs}

Let \(\Phi:M_{2^m}(\C)\to M_{2^n}(\C)\) be HP, TP and full
Choi--LY, and put \(\phi=\Phi^\dagger\).  An output site \(i\)
is called \emph{transversely visible} if at least one of
\(\phi(X_i),\phi(Y_i)\) is nonscalar.  Denote the set of such
sites by \(A\).

\begin{lemma}[Local routing]
\label{stat:routing}
For each \(i\in A\) there is a unique input site \(\pi(i)\) such
that, writing \(k=\pi(i)\),
\begin{equation}
 \begin{split}
 \phi(X_i)&=b_{ix}I+D_{i,11}X_k+D_{i,21}Y_k,\\
 \phi(Y_i)&=b_{iy}I+D_{i,12}X_k+D_{i,22}Y_k,\\
 \phi(Z_i)&=\lambda_i Z_k,
 \end{split}
 \label{stat:local-block}
\end{equation}
where all coefficients are real and \(D_i\ne0\).  No injectivity
of \(\Phi\) or \(D_i\) is required.
\end{lemma}

\begin{proof}
For a local \(U\in\mathfrak g_i\), the identity
\(U^2=\delta I\), \(\delta\in\R\), gives
\(e^{tU}=c_\delta(t)I+s_\delta(t)U\), with real analytic scalar
coefficients satisfying \(c_\delta(0)=1\), \(s_\delta(0)=0\),
and \(s_\delta'(0)=1\).  These exponentials are stability units.
Since \(\phi\) is unital, their images, divided by
\(c_\delta(t)\), give a two-sided stable pencil
\(I+r\phi(U)\) near zero.  A zero image is excluded near \(I\).
Lemma~\ref{stat:edge} puts \(\phi(U)\) in \(\mathfrak e_m\).
HP then implies
\[
 \phi(X_i),\phi(Y_i)\in\R I+
       \sum_k\operatorname{span}_{\R}\{X_k,Y_k\},
 \qquad
 \phi(Z_i)\in\R I+\sum_k\R Z_k.
\]
Every nonscalar transverse image has a hyperbolic component, so
Lemma~\ref{stat:hyperbolic} confines it to one input site.  If two
transverse directions used different sites, their sum would have
two nonzero hyperbolic components.  Thus the whole transverse plane
uses the same site whenever its nonscalar image is nonzero.

Choose a transverse \(U_i\) with nonzero traceless image on this
site, and apply the argument to \(U_i+i\varepsilon Z_i\) for a
small nonzero real \(\varepsilon\).  The routed component remains
hyperbolic.  Lemma~\ref{stat:hyperbolic} eliminates all other
longitudinal components, and its real-scalar conclusion eliminates
the scalar coefficient of \(\phi(Z_i)\).  This proves
\eqref{stat:local-block}.
\end{proof}

\begin{theorem}[Visible-output decomposition]
\label{stat:visible}
Let \(\Phi:M_{2^m}(\C)\to M_{2^n}(\C)\) be HP, TP and full
Choi--LY.  With the preceding notation, put \(K=\pi(A)\) and
\(G_k=\pi^{-1}(k)\).  There are HP, TP, full Choi--LY maps
\(\Xi_k:M_2(\C)\to M_{2^{|G_k|}}(\C)\) such that, after
ordering the physical output groups,
\begin{equation}
 \Tr_{A^c}\circ\Phi
 =\left(\bigotimes_{k\in K}\Xi_k\right)\circ\Tr_{K^c}.
 \label{stat:visible-form}
\end{equation}
The routing function is unique.  If \(\Phi\) is CP, every
\(\Xi_k\) is CP.  The trace on the right is the ordinary,
unnormalized partial trace.
\end{theorem}

\begin{proof}
For each \(i\in A\), choose a small real open set
\(\mathcal O_i\subset\mathfrak g_i\) about a transverse direction
whose image is nonscalar.  Shrink it so that every \(U_i\in
\mathcal O_i\), and the traceless part of \(\phi(U_i)\), is
hyperbolic.  By \eqref{stat:local-block},
\(\phi(U_i)\in\mathcal E_{\pi(i)}\).  These open sets have
complex span equal to the entire local traceless matrix space,
even when \(D_i\) has rank one.

Fix choices \(U_i\in\mathcal O_i\).  For \(S\subseteq A\) set
\[
 F_S(t_S)=\phi\left(\prod_{i\in S}(I+t_iU_i)\right).
\]
We prove by induction that, for all sufficiently small real parameters,
\begin{equation}
 F_S(t_S)=\prod_{k\in\pi(S)} F_{k,S}(t_{S\cap G_k}),
 \qquad F_{k,S}\in\mathcal E_k,\qquad F_{k,S}(0)=I,
 \label{stat:group-induction}
\end{equation}
where each factor is multiaffine in its own parameters.  Distinct
factors act on distinct sites.  They and their inverses are stability
units near zero.  The empty and singleton cases follow from unitality
and Lemma~\ref{stat:routing}.

Suppose \eqref{stat:group-induction} holds, and add \(i\notin S\),
with route \(k\).  Put
\[
 \Delta(t_S)=\phi\left(U_i\prod_{j\in S}(I+t_jU_j)\right),
 \qquad B(t_S)=F_S(t_S)^{-1}\Delta(t_S).
\]
The new input product is tensor-LY.  Its image is
\(F_S+t_i\Delta\), which is nonzero near the identity.  Left
multiplication by the stable inverse gives the two-sided stable
pencil \(I+t_iB(t_S)\).  Lemma~\ref{stat:edge} applies to this
pencil.  At \(t_S=0\) its coefficient is \(\phi(U_i)\), whose
component on site \(k\) is hyperbolic.  Continuity and
Lemma~\ref{stat:hyperbolic} therefore imply
\begin{equation}
 B(t_S)\in\mathcal E_k
 \label{stat:localized-increment}
\end{equation}
on a sufficiently small parameter neighborhood.

Fix a parameter \(t_j\) belonging to a different group \(G_\ell\),
and fix all the other parameters.  Write
\(F_{\ell,S}=C_0+t_jC_1\).  The matrices \(C_0,C_1\) are
linearly independent for sufficiently small fixed values of the
other parameters: at zero they are \(I,\phi(U_j)\), with a
nonzero hyperbolic traceless part in the second matrix.
Remove the fixed other-site factors from \(\Delta=F_SB\).
An affine function of \(t_j\) is thereby written as
\[
 (C_0+t_jC_1)\otimes R(t_j),\qquad R(t_j)=F_{k,S}B(t_S).
\]
Choose linear functionals \(\alpha,\beta\) with
\[
 \alpha(C_0)=1,\quad\alpha(C_1)=0,\qquad
 \beta(C_0)=0,\quad\beta(C_1)=1.
\]
Applying them shows that both \(R(t_j)\) and \(t_jR(t_j)\) are
affine.  Hence \(R\), and therefore \(B\), is constant in
\(t_j\).  Repeating this for every parameter outside \(G_k\)
gives the update
\[
 F_{S\cup\{i\}}=
 \prod_{\ell\ne k}F_{\ell,S}
 \left(F_{k,S}+t_iF_{k,S}B(t_{S\cap G_k})\right).
\]
The last factor lies in the real algebra \(\mathcal E_k\), which
is closed under multiplication.  It is multiaffine, as is seen by
setting every other group's parameters to zero in the original
definition of \(F_{S\cup\{i\}}\).  This completes the induction;
only finitely many neighborhood shrinkings are needed.

Setting the other groups' parameters to zero identifies each
\(F_{k,S}\) as the image of its own group product.  Extracting
coefficients in \eqref{stat:group-induction} consequently gives
\[
 \phi\left(\prod_{i\in S}U_i\right)=
 \prod_{k\in\pi(S)}
 \phi\left(\prod_{i\in S\cap G_k}U_i\right).
\]
This is a complex multilinear identity.  Its validity on a product
of real open sets extends first to their real spans and then to their
complex spans.  Adjoining identity factors spans the entire matrix
algebra on \(A\).  Thus the restriction of \(\phi\) to that
algebra is the tensor product of maps on the groups, each with image
on its routed input site.

For an observable \(O\) on \(G_k\), extract the one-site image of
\(\phi(O\otimes I)\) by normalized partial trace over the other
input sites.  This gives a unital HP map \(\psi_k\), CP if
\(\phi\) is CP.  Its adjoint \(\Xi_k\) is TP and has the same
positivity properties.  Equivalently, \(\Xi_k\) is obtained by
preparing every other input maximally mixed, applying \(\Phi\),
and tracing the other outputs.  These operations have tensor-LY
Choi arrays.  Their contraction is nonzero because \(\Xi_k\) is
TP, so \(\Xi_k\) is full Choi--LY.  Taking adjoints proves
\eqref{stat:visible-form}.  The adjoint of the partial trace in
that formula tensors with identity, which also verifies its
normalization.
\end{proof}

\begin{theorem}[Invertible static classification]
\label{stat:invertible}
An invertible complex-linear HP, TP map
\(\Phi:M_{2^n}(\C)\to M_{2^n}(\C)\) is full Choi--LY if and
only if
\begin{equation}
 \Phi=(\Phi_1\otimes\cdots\otimes\Phi_n)\circ\mathcal P,
 \label{stat:invertible-form}
\end{equation}
where \(\mathcal P\) permutes the physical input wires and the
\(\Phi_i\) are invertible one-qubit HP, TP, full Choi--LY maps.
The map \(\Phi\) is CP if and only if all its local factors are CP.
\end{theorem}

\begin{proof}
Apply the preceding results to \(\phi=\Phi^\dagger\).  The induced
map on the real space of transverse fields modulo scalars is injective.
Indeed, if \(\phi(U)=bI\), unitality gives
\(\phi(U-bI)=0\); injectivity forces \(U=bI\), and a traceless
transverse field is then zero.  In particular every output is visible,
each of its two-dimensional transverse planes maps injectively to its
routed two-dimensional input plane, and two such planes cannot route
to the same input plane.  With \(n\) planes on each side, the routes
form a permutation and all groups have one site.

Theorem~\ref{stat:visible} now gives a tensor product of local factors.
Each factor is invertible: a nonzero kernel vector of one factor,
tensored with suitable nonzero vectors for the other factors, would
give a kernel vector of the full map.  The same theorem proves
inheritance of CP.  Conversely, tensor products and physical wire
permutations preserve full Choi--LY and all the stated properties.
\end{proof}

\begin{remark}
The proof of Theorem~\ref{stat:invertible} does not require a
semigroup representation, complete positivity, or injectivity of an
unproved local restriction.  For completeness, its adjoint local
longitudinal coefficient in \eqref{stat:local-block} is positive:
stability of \(I+s\lambda_iZ\), obtained from \(I+sZ\) for small
\(s>0\), gives \(|1+s\lambda_i|\ge|1-s\lambda_i|\), hence
\(\lambda_i\ge0\); injectivity excludes zero.
\end{remark}

\begin{example}[Ordinary static preservation is weaker]
\label{stat:ordinary-counterexample}
Let \(n\ge2\), \(d=2^n\),
\(\tau=\diag(3/4,1/4)\), and \(\sigma=\tau^{\otimes n}\).
For
\[
 0<a<\frac1{4^n+1},\qquad
 \Phi_a(M)=aM+(1-a)\Tr(M)\sigma,
\]
the map \(\Phi_a\) is invertible and CPTP, and sends every density
matrix to a matrix with a polynomial nonzero on the closed unit
polydisk.  Nevertheless it is neither a product of local channels
up to a physical permutation nor full Choi--LY.

Indeed, it fixes \(\sigma\) and multiplies the trace-zero subspace
by \(a\), proving invertibility.  On the closed polydisk,
\[
 |G_\sigma|\ge2^{-n}=d^{-1},\qquad |G_\rho|\le d
\]
for every density matrix \(\rho\).  The second inequality follows
by bounding the two monomial-vector norms by \(\sqrt d\) and
\(\|\rho\|_{\mathrm{op}}\) by one.  Hence
\[
 |G_{\Phi_a(\rho)}|\ge(1-a)d^{-1}-ad>0.
\]
On input \(|0^n\rangle\langle0^n|\), the occupation indicators at
any two distinct output sites have covariance \(a(1-a)/16>0\).
The output is not product, whereas a product of local channels,
with or without a wire permutation, sends product inputs to products.

One may also exhibit the Choi zero directly.  Its polynomial is
\[
 a\prod_i(1+z_ix_i)(1+w_iy_i)
 +(1-a)G_\sigma(z,w)\prod_i(1+x_iy_i).
\]
Set all output variables except \(z_1,w_1\) to zero, and set
\(z_1=w_1=t\), \(x_1=y_1=v\).  Define
\[
 p(t)=(3/4)^{n-1}(3/4+t^2/4),\quad
 C(t)=\frac{1-a}{a}p(t),\quad C_0=C(1).
\]
Choose \(v_0=e^{i\theta}\) with
\(\cos\theta=-1/(1+C_0)\), \(\pi/2<\theta<\pi\).  Then
\((1+v_0)^2/(1+v_0^2)=-C_0\).
For small \(\varepsilon>0\), put
\[
 t=1-\varepsilon,\qquad v=(1-\varepsilon)v_0,
 \qquad q=-\frac{(1+tv)^2}{C(t)(1+v^2)}-1.
\]
As \(\varepsilon\downarrow0\), \(q\to0\), so all of
\(|t|,|v|,|q|\) are smaller than one for sufficiently small
positive \(\varepsilon\).  Set \(x_2=y_2=\sqrt q\) and all
remaining reference variables to zero.  The Choi polynomial becomes
\(a(1+tv)^2+(1-a)p(t)(1+v^2)(1+q)=0\), at an interior point.
Lemma~\ref{stat:bridge} turns this into an ancillary physical witness.
The same example is the late-time channel
\(e^{t(\mathcal R_\sigma-\Id)}=e^{-t}\Id+(1-e^{-t})\mathcal R_\sigma\),
where \(\mathcal R_\sigma(M)=\Tr(M)\sigma\).  Its ordinary
strict-preservation property for \(t>\log(4^n+1)\) therefore does
not imply preservation at times approaching zero.
\end{example}

\subsection{Saturated columns and reflected polynomial factors}

For a polynomial \(p\), let \(\|p\|_2\) be the Euclidean norm of
its coefficient vector.  If \(p\) is multiaffine in a specified
finite set \(S\), define its full reflection by
\begin{equation}
 p^*(z)=z^S\overline{p(1/\bar z)}.
 \label{stat:reflection}
\end{equation}
The set \(S\) is part of the definition, and may include variables
absent from \(p\).  In contrast, \(\overline p(w)=
\overline{p(\bar w)}\) below denotes coefficientwise conjugation
without reflection.

\begin{lemma}[Saturation and disjoint factors]
\label{stat:saturation}
Let \(F(z,x)=\sum_{b\in\{0,1\}^m}A_b(z)x^b\) be a stable
multiaffine polynomial.  Suppose
\(\|A_{e_i}\|_2=\|A_0\|_2\) for every input site \(i\).
Then, over the rational function field in \(z\),
\begin{equation}
 F(z,x)=A_0(z)\prod_{i=1}^m
       \left(1+x_i\frac{A_{e_i}(z)}{A_0(z)}\right).
 \label{stat:saturation-product}
\end{equation}
Writing the quotients in coprime form \(q_i/p_i\), there is a
polynomial factorization
\begin{equation}
 F(z,x)=C(z)\prod_{i=1}^m(p_i(z)+x_iq_i(z))
 \label{stat:primitive-product}
\end{equation}
in which the output-variable supports of the displayed factors are
pairwise disjoint.  All these factors are stable, and
\(\|p_i\|_2=\|q_i\|_2\).
\end{lemma}

\begin{proof}
Setting all input variables except \(x_i\) to zero gives a stable
affine polynomial \(A_0+x_iA_{e_i}\), whence
\(|A_{e_i}(z)|\le|A_0(z)|\) in the output polydisk.
The inequality extends continuously to the output torus.  Parseval
and norm equality imply equality of these moduli everywhere on that
torus.  At a torus point where \(A_0\ne0\), boundary Hurwitz makes
the input polynomial stable; its nonzero constant coefficient
excludes the zero alternative.  Write this input polynomial as
\(P_i(x_{\ne i})+x_iQ_i(x_{\ne i})\).  The holomorphic ratio
\(Q_i/P_i\) is contractive on the remaining input polydisk, and
has modulus one at its origin.  It is therefore constant by the
maximum modulus principle.  Applying this to every input gives
\eqref{stat:saturation-product} at these torus points.
The zero set of the nonzero polynomial \(A_0\) has torus measure
zero.  Clearing denominators and using uniqueness of Fourier
coefficients proves the rational identity everywhere.

Each \(p_i+x_iq_i\) is primitive over \(\C[z]\) and irreducible
over the rational function field, where it is linear in its own
input variable.  Gauss's lemma applied to the primitive part of
\(F\), considered as a polynomial in \(x\), gives
\eqref{stat:primitive-product} with \(C\) polynomial.  Degrees in
any one output variable add under multiplication in this integral
domain.  Since every such degree of \(F\) is at most one, each
output variable belongs to at most one displayed factor.  Stability
of \(F\) implies stability of every factor and, by setting
\(x_i=0\), of \(p_i\).  Finally coefficient norms multiply on
disjoint supports.  Comparing \(A_0\) with \(A_{e_i}\) gives the
claimed local norm equality.
\end{proof}

\begin{lemma}[Reduced reflected pairs]
\label{stat:reflected-pair}
Let \(p+xq\) be stable and multiaffine, where \(p,q\) are
coprime, have equal coefficient norms, and their union of output
supports is \(S\).  Then
\begin{equation}
 q=e^{i\theta}p^*
 \label{stat:pair-reflection}
\end{equation}
for some real \(\theta\).  Conversely, for any stable multiaffine
\(p\) on a specified set \(S\), the polynomial
\(p+e^{i\theta}xp^*\) is stable.
\end{lemma}

\begin{proof}
Stable affine slices and equality of norms imply \(|p|=|q|\) on
the output torus.  Multiplication by \(z^S\) gives
\(pp^*=qq^*\).  Coprimality implies \(p\mid q^*\); write
\(q^*=ph\), so \(p^*=qh\).
The constant coefficient of the stable polynomial \(p\) is nonzero.
Thus \(p^*\) has degree one in every variable of \(S\).
If \(h\) depends on \(z_j\), degree additivity makes \(q\)
independent of \(z_j\), and the support-union property makes
\(p\) depend on \(z_j\).  Independence of \(q\) implies
\(z_j\mid q^*=ph\).  Since \(p\) has nonzero constant term,
\(z_j\mid h\), and then \(z_j\mid p^*=qh\).  The latter
says \(p\) is independent of \(z_j\), a contradiction.  Hence
\(h\) is constant.  Reflecting twice shows \(|h|=1\), proving
\eqref{stat:pair-reflection}.

For the converse, \(p^*/p\) is holomorphic on the polydisk and
has modulus at most one.  To justify this assertion even in the
presence of boundary zeros, use \(p_r(z)=p(rz)\), \(r<1\).
The polynomial \(p_r\) has no closed-polydisk zero, and
\(p_r^*/p_r\) has unit modulus on the torus.  Successive maximum
modulus principles make it contractive inside.  Letting \(r\uparrow1\)
gives the asserted bound.  Hence
\(1+e^{i\theta}x p^*/p\ne0\) for \(x\in\D\).
\end{proof}

Reflection descriptions of rational inner functions are classical;
see \cite{Knese24}.  The elementary proof above specifies
the full support set and the coprime multiaffine reduction used here.

\subsection{All stable isometries and their coding obstruction}

\begin{theorem}[Stable isometries]
\label{stat:isometry}
Let \(V:(\C^2)^{\otimes m}\to(\C^2)^{\otimes n}\) be an
isometry.  Its coefficient polynomial is stable if and only if,
after a physical output permutation,
\begin{equation}
 V=|\psi\rangle\otimes V_1\otimes\cdots\otimes V_m,
 \label{stat:isometry-form}
\end{equation}
where \(\psi\) is a normalized stable pure state on an optional
prepared block \(S_0\).  The remaining nonempty, disjoint blocks
\(S_i\) carry one-input isometries whose polynomials are
\begin{equation}
 G_{V_i}(z,x_i)=
 \frac{p_i(z)+e^{i\theta_i}x_i p_i^*(z)}{\|p_i\|_2},
 \qquad \gcd(p_i,p_i^*)=1,
 \qquad [z^{S_i}]p_i(z)^2=0.
 \label{stat:isometry-block}
\end{equation}
Each \(p_i\) is stable and multiaffine, and reflection uses the
whole block \(S_i\).  The sets \(S_0,S_1,\ldots,S_m\) partition
the physical outputs.  In particular, an isometric channel
\(M\mapsto VMV^*\) is full Choi--LY exactly for these isometries.
\end{theorem}

\begin{proof}
Write \(G_V=\sum_bA_b(z)x^b\).  Isometry gives
\(\langle A_b,A_c\rangle=\delta_{bc}\), so
Lemma~\ref{stat:saturation} applies.  In its factorization,
disjointness makes both norms and inner products multiplicative.
Comparing columns \(0\) and \(e_i\) yields
\(\|p_i\|_2=\|q_i\|_2>0\) and
\(\langle p_i,q_i\rangle=0\).  A factor with no output variable
could not satisfy the latter equality, so every logical block is
nonempty.  Normalize the common prepared factor and each pair.
There is no residual normalization, because the all-zero column
has norm one.  Variables absent from the polynomial are fixed-zero
prepared outputs.

Lemma~\ref{stat:reflected-pair} gives the reflected form.
Furthermore,
\[
 \langle p_i,p_i^*\rangle
 =\overline{\sum_{T\subseteq S_i}(p_i)_T(p_i)_{S_i\setminus T}}
 =\overline{[z^{S_i}]p_i(z)^2},
\]
so orthogonality is exactly the final condition in
\eqref{stat:isometry-block}.  Conversely that condition gives
orthogonal columns, reflection preserves their norms, and
Lemma~\ref{stat:reflected-pair} gives stability.  Tensor products
prove the converse for \(V\).  The isometric channel's Choi
polynomial is \(G_V(z,x)\overline{G_V}(w,y)\) in independent
variables, proving the final equivalence.
\end{proof}

\begin{lemma}[Strict longitudinal bias]
\label{stat:strict-bias}
If \(p\) is stable and multiaffine on a nonempty specified set
\(S\), and \(\gcd(p,p^*)=1\), then for every \(j\in S\),
\begin{equation}
 \|p|_{z_j=0}\|_2^2>\|\partial_jp\|_2^2.
 \label{stat:strict-slice}
\end{equation}
\end{lemma}

\begin{proof}
Write \(p=a+z_jb\).  Stable affine slices imply \(|b|\le|a|\)
on the remaining polydisk, hence \(\|b\|_2\le\|a\|_2\).
If \(b=0\), strictness follows from \(a\ne0\).
Otherwise suppose equality holds.  Torus saturation gives
\(aa^*=bb^*\), reflecting both polynomials over
\(R=S\setminus\{j\}\), while the full reflection is
\(p^*=b^*+z_ja^*\).
Let \(g=\gcd(a,b)\) and
\(h=a/g+z_jb/g\).  This polynomial is primitive over \(\C[z_R]\)
and has positive degree in \(z_j\).  Over \(\C(z_R)[z_j]\),
\[
 p^*=\frac{a^*}{b}(a+z_jb).
\]
Thus \(h\) divides both \(p\) and \(p^*\) over the fraction
field.  Gauss's lemma gives the same divisibility in the polynomial
ring, contradicting coprimality.  This proof also covers empty
\(R\) and variables absent from either slice.
\end{proof}

\begin{corollary}[Single-output erasure]
\label{stat:no-code}
For a reduced logical block in Theorem~\ref{stat:isometry}, every
physical output \(j\in S_i\) satisfies
\begin{equation}
 V_i^*Z_jV_i=\delta_j Z_i,
 \qquad
 \delta_j=
 \frac{\|p_i|_{z_j=0}\|_2^2-\|\partial_jp_i\|_2^2}
      {\|p_i\|_2^2}>0.
 \label{stat:bias-compression}
\end{equation}
Erasure of any such output is not exactly correctable on the
logical space.  Every nontrivial stable isometric quantum code
therefore has distance one.
\end{corollary}

\begin{proof}
Reflection complements occupied subsets, so the two diagonal
expectations of \(Z_j\) have opposite signs.  The cross term is
zero because complementary subsets cancel:
\[
 \langle p_i,Z_jp_i^*\rangle
 =\overline{\sum_{T\subseteq S_i}
       (-1)^{\mathbf{1}_{j\in T}}(p_i)_T(p_i)_{S_i\setminus T}}=0.
\]
Lemma~\ref{stat:strict-bias} now proves
\eqref{stat:bias-compression}.  For erasure of output \(j\),
the two Kraus operators may be taken as
\(E_a=\langle a|_j\otimes I\), \(a=0,1\).
The Knill--Laflamme criterion \cite{KnillLaflamme97} requires
\(V^*E_a^*E_bV\) to be scalar for all \(a,b\).
Subtracting the diagonal cases requires \(V^*Z_jV\) to be scalar,
contrary to \eqref{stat:bias-compression}.  Equivalently, this
weight-one phase observable is not detectable by the code.
\end{proof}

\begin{remark}
Outputs in the prepared block are exempt from the erasure conclusion.
Nor is there a uniform positive lower bound on the individual bias:
for \(p(z_1,z_2)=1+r z_1\), \(0\le r<1\), reflected over
\(\{1,2\}\), one has \(p^*=z_2(r+z_1)\), coprime orthogonal
columns, and \(\delta_1=(1-r^2)/(1+r^2)\to0\).  The conclusion
concerns exact correction only.
\end{remark}

\begin{example}[Coherent information invisible in one-site marginals]
\label{stat:repetition}
The repetition isometry \(V|0\rangle=|00\rangle\),
\(V|1\rangle=|11\rangle\) has polynomial \(1+xz_1z_2\).
The square channel \(\Phi(M)=V\Tr_2(M)V^*\) is full Choi--LY,
with polynomial
\[
 (1+x_1z_1z_2)(1+y_1w_1w_2)(1+x_2y_2).
\]
All its single-output transverse adjoints vanish, although
\(\Phi^\dagger(X_1X_2)=X_1\otimes I_2\).
Thus Theorem~\ref{stat:visible} cannot identify its omitted
outputs as a classical or input-independent sector.
\end{example}

\begin{proposition}[No stable Stinespring representation in general]
\label{stat:no-dilation}
There is a two-qubit CPTP full Choi--LY map with no Stinespring
isometry whose coefficient tensor is LY, for any finite qubit
environment.
\end{proposition}

\begin{proof}
Let \(\Delta\) be complete computational-basis dephasing on two
qubits, let \(P\) swap the qubits, and set
\(\Phi=\tfrac12\Delta+\tfrac12\operatorname{Ad}_P\circ\Delta\).
This is CPTP.  With \(r_j=z_jw_j\) and \(s_j=x_jy_j\), its
Choi polynomial is
\begin{equation}
 \begin{split}
 K(r,s)&=\tfrac12(1+r_1s_1)(1+r_2s_2)
       +\tfrac12(1+r_1s_2)(1+r_2s_1)\\
 &=1+\tfrac12(r_1+r_2)(s_1+s_2)+r_1r_2s_1s_2.
 \end{split}
 \label{stat:symmetrized-kernel}
\end{equation}
For fixed \(s_1,s_2\in\D\), this is the symmetric multiaffine
polarization in \(r_1,r_2\) of
\((1+rs_1)(1+rs_2)\).  The Grace--Walsh--Szeg\H{o} coincidence
theorem on the convex circular domain \(\D\) implies its
nonvanishing for \(r_1,r_2\in\D\); see \cite{BB10II}.
Thus the full Choi tensor is stable.

However,
\(\Phi^\dagger(Z_1)=\Phi^\dagger(Z_2)=(Z_1+Z_2)/2\).
If a stable Stinespring isometry existed, apply
Theorem~\ref{stat:isometry} to its physical and environmental
output qubits together.  Every one-site observable on a physical
output would pull back either to a scalar or to an operator on
one logical input, because its output belongs to one disjoint
encoding block or to the prepared block.  Tracing the environment
merely inserts identities at its sites.  The displayed two-input
longitudinal observable is incompatible with this conclusion.
\end{proof}

\subsection{All Schur channels}

\begin{lemma}[A scalar equal-diagonal criterion]
\label{stat:two-by-two}
For a nonzero PSD matrix
\(H=\begin{psmallmatrix}a&b\\\bar b&d\end{psmallmatrix}\),
its tensor polynomial is stable if and only if \(a>0\) and
\(a\ge d\).  For arbitrary complex \(\beta,\gamma\), the
polynomial \(1+\beta x+\gamma y+xy\) is stable if and only if
\(\gamma=\bar\beta\) and \(|\beta|\le1\).
\end{lemma}

\begin{proof}
A stable polynomial must have nonzero constant coefficient, so
\(a>0\).  The necessary boundary root-modulus inequality is
\[
 |a+\bar b\zeta|^2-|b+d\zeta|^2
 =(a-d)(a+d+2\Re(\bar b\zeta))\ge0,
 \qquad |\zeta|=1.
\]
Its average gives \(a\ge d\).  Conversely, PSD gives
\(|b|^2\le ad\), and \(a\ge d\) gives \(|b|\le a\).
If \(|b|<a\), the affine constant factor \(a+\bar b z\)
is nonzero on the closed disk, and the displayed inequality
and the maximum modulus principle bound the other coefficient
divided by this one by one.  This proves stability.  If
\(|b|=a\), PSD and \(a\ge d\) force \(d=a\), and the
polynomial factors as
\(a(1+(\bar b/a)z)(1+(b/a)w)\), also stable.

For the second assertion, boundary root-modulus inequalities give
\(|1+\beta\zeta|\ge|\gamma+\zeta|\), and the analogous
inequality with \(\beta,\gamma\) exchanged.  Squaring and
averaging gives \(|\beta|=|\gamma|\).  The remaining nonnegative
linear trigonometric polynomial has mean zero, so vanishes
identically, yielding \(\beta=\bar\gamma\).  Setting one
variable to zero gives \(|\beta|\le1\).  The first assertion
applied to the resulting equal-diagonal PSD matrix proves the
converse.
\end{proof}

\begin{theorem}[Stable correlation matrices and Schur channels]
\label{stat:schur}
Let \(C\succeq0\) be a \(2^m\)-by-\(2^m\) matrix with
\(C_{aa}=1\).  Its tensor polynomial is stable if and only if
\begin{equation}
 C=\bigotimes_{i=1}^m
 \begin{pmatrix}1&\eta_i\\\bar\eta_i&1\end{pmatrix},
 \qquad |\eta_i|\le1.
 \label{stat:correlation-product}
\end{equation}
Equivalently, a CPTP map fixing every computational-basis state is
full Choi--LY if and only if it is a tensor product of one-qubit
phase/dephasing channels.  Parameters \(\eta_i=0\) are allowed.
\end{theorem}

\begin{proof}
Separate one physical site and write
\(C=\begin{psmallmatrix}A&B\\B^*&D\end{psmallmatrix}\).
For arbitrary \(z\in\D^{m-1}\), compress the other sites against
the product vector \(v(z)=\sum_a z^a|a\rangle\).
The two-by-two compression is PSD.  Its tensor polynomial is an
interior specialization of \(G_C\), hence stable, and its
constant coefficient is strictly positive.  By
Lemma~\ref{stat:two-by-two},
\[
 \langle v(z),(A-D)v(z)\rangle\ge0.
\]
At any fixed positive interior coordinate moduli, the angular
average of this expression is
\(\sum_a(A_{aa}-D_{aa})|z|^{2a}=0\).
Continuity and nonnegativity force pointwise equality on every
such torus.  Varying the moduli proves the polynomial identity
\(\langle v(z),(A-D)v(z)\rangle=0\).  Independence of the
monomials \(\bar z^a z^b\) gives \(A=D\) entrywise.

Writing \(u\) for the remaining independent row and column
variables, the full polynomial has the form
\[
 a(u)(1+xy)+b(u)x+c(u)y.
\]
The specialization \(a(u)\) is nonzero throughout the polydisk.
Lemma~\ref{stat:two-by-two} gives
\(c(u)/a(u)=\overline{b(u)/a(u)}\) and
\(|b(u)/a(u)|\le1\) for every complex \(u\) there.
Both ratios are holomorphic on this connected domain, so they
are constant.  This factors off one matrix in
\eqref{stat:correlation-product}.  The remaining factor \(A\)
is PSD, has diagonal one, and is tensor-LY.  Induction proves
the assertion.  The converse follows from the one-site criterion
and multiplication on disjoint variables.

For a Schur multiplier \(\mathcal S_C(M)=C\circ M\), CP and
TP are equivalent to \(C\succeq0\) and diagonal one.  Its full
Choi polynomial is \(G_C(r,s)\), where \(r_i=z_iu_i\) and
\(s_i=w_iv_i\).  Each disk value is attainable by these
products, so full Choi stability is exactly tensor stability
of \(C\).  Finally, any CPTP map fixing all computational-basis
states has diagonal Kraus operators: each Kraus image of a basis
vector must lie in the one-dimensional range of the corresponding
pure output.  It is therefore a Schur map, proving the equivalent
formulation.
\end{proof}

\subsection{Channels with pure computational-basis outputs}

\begin{theorem}[Pure-basis-output classification]
\label{stat:pure-basis}
Let \(\Phi:M_{2^m}(\C)\to M_{2^n}(\C)\) be CPTP, and suppose
every computational-basis output \(\Phi(|a\rangle\langle a|)\)
is pure.  The outputs need not be orthogonal or distinct.  Then
\(\Phi\) is full Choi--LY if and only if it has the following
form.

Partition the inputs as \(E\sqcup K=[m]\), and the outputs into
an optional prepared block \(S_0\) and nonempty disjoint blocks
\(S_i\), \(i\in K\).  Prepare a normalized stable pure state
\(\psi\) on \(S_0\).  On each \(S_i\) take a reduced stable
two-column matrix \(V_i\) with normalized, linearly independent
columns \(v_{i0},v_{i1}\).  Reduced means that their amplitude
polynomials are coprime.  Equivalently,
\begin{equation}
 G_{V_i}(z,x_i)=
 \frac{p_i(z)+e^{i\theta_i}x_i p_i^*(z)}{\|p_i\|_2},
 \qquad \gcd(p_i,p_i^*)=1,
 \label{stat:pure-pair}
\end{equation}
where \(p_i\) is stable and multiaffine and reflection uses the
full block.  Put \(\mu_i=\langle v_{i0},v_{i1}\rangle\), and
choose \(|\eta_i|\le1\) subject to
\begin{equation}
 \mu_i\eta_i=0.
 \label{stat:overlap-noise}
\end{equation}
With
\[
 C_i=\begin{pmatrix}1&\eta_i\\\bar\eta_i&1\end{pmatrix},
 \qquad
 \mathcal B_i(M)=V_i(C_i\circ M)V_i^*,
\]
the channel, after output permutation, is exactly
\begin{equation}
 \Phi(\rho)=|\psi\rangle\langle\psi|\otimes
       \left(\bigotimes_{i\in K}\mathcal B_i\right)
                         (\Tr_E\rho).
 \label{stat:pure-basis-form}
\end{equation}
Empty \(K\) is allowed.  If \(\mu_i=0\), the block is a stable
isometric encoding preceded by arbitrary local phase/dephasing.
If \(\mu_i\ne0\), it measures the input in the computational
basis and prepares the corresponding nonorthogonal reflected
pure state.
\end{theorem}

\begin{proof}
Assume first that \(\Phi\) is full Choi--LY.  Choose normalized
output vectors \(\psi_a\), with amplitude polynomials \(p_a\).
Complete input dephasing is full Choi--LY; its Choi polynomial is
a product of factors \(1+x_iy_iu_iv_i\).  After composing with
this map and grouping each input row/column pair into its product,
we obtain the stable polynomial
\begin{equation}
 F(z,w,s)=\sum_a A_a(z,w)s^a,
 \qquad A_a(z,w)=p_a(z)\overline p_a(w).
 \label{stat:pure-density-kernel}
\end{equation}
Grouping is an equivalence of disk stability because disk-variable
products cover the disk.  Each \(A_a\) has coefficient norm one,
being the polynomial of a normalized rank-one projection.
Lemma~\ref{stat:saturation} gives, for every nonempty subset
\(T\subseteq[m]\),
\begin{equation}
 A_T A_0^{|T|-1}=\prod_{i\in T}A_{e_i}.
 \label{stat:pure-column-identity}
\end{equation}

Substitute \eqref{stat:pure-density-kernel} into this identity.
The rational function
\[
 R_T(z)=\frac{p_T(z)p_0(z)^{|T|-1}}
                   {\prod_{i\in T}p_{e_i}(z)}
\]
satisfies \(R_T(z)\overline{R_T(\bar w)}=1\) in independent
variables.  Fixing a generic \(w\) shows that it is a constant
\(\xi_T\) of modulus one.  We may therefore align the output
column phases so that the polynomial
\begin{equation}
 G(z,x)=p_0(z)\prod_i
              \left(1+x_i\frac{p_{e_i}(z)}{p_0(z)}\right)
 \label{stat:phase-aligned-amplitude}
\end{equation}
has exactly those phase-aligned columns.  More explicitly, its
coefficient at \(x^T\) is \(\xi_T^{-1}p_T\); hence it is a
multiaffine polynomial, not merely a rational expression.
The stable constant slice \(A_0=p_0\overline p_0\) makes
\(p_0\) zero-free in the output polydisk.  The singleton modulus
domination for \eqref{stat:pure-density-kernel}, specialized to
\(w=\bar z\), gives \(|p_{e_i}/p_0|\le1\).
Thus \eqref{stat:phase-aligned-amplitude} is stable.  It defines
a matrix \(V\) with normalized columns; no orthogonality is
being asserted.

Apply Lemma~\ref{stat:saturation} to \(G\), whose column norms
are all one, and then Lemma~\ref{stat:reflected-pair} to its
primitive factors.  We obtain a stable prepared factor and
disjoint reduced reflected pairs with equal local norms.  A
dependent primitive pair \(q_i=c_ip_i\) must have both
polynomials constant, so has empty output support; conversely
every nonempty primitive pair has independent columns.  Assign
the empty factors to \(E\) and the others to \(K\).  Their
equal norms imply \(|c_i|=1\).  Replacing each erased input
variable by \(c_i^{-1}x_i\) is a disk-preserving phase change
of columns, and makes the two erased-input columns identical.
Normalize the prepared factor and all retained pairs.  The total
normalization is one because the all-zero full column has norm
one.  All aligned output vectors now have the form
\begin{equation}
 v_a=\psi\otimes\bigotimes_{i\in K}v_{i,a_i},
 \label{stat:aligned-vectors}
\end{equation}
independent of \(a_E\).  Common prepared factors and variables
absent from \(G\) belong to \(S_0\).  This proves all the
polynomial and support assertions about the pairs in the theorem.

Choose Kraus operators \(K_\ell\) of \(\Phi\).  Since the
positive sum of their output projections on input \(a\) has
rank one, every \(K_\ell|a\rangle\) is proportional to \(v_a\).
Writing \(K_\ell|a\rangle=c_{\ell a}v_a\) gives
\begin{equation}
 \Phi(|a\rangle\langle b|)=C_{ab}|v_a\rangle\langle v_b|,
 \qquad C_{ab}=\sum_\ell c_{\ell a}\bar c_{\ell b},
 \label{stat:kraus-correlation}
\end{equation}
where \(C\succeq0\) and \(C_{aa}=1\).  Original trace
preservation is exactly
\begin{equation}
 C_{ab}\langle v_b,v_a\rangle=\delta_{ab}.
 \label{stat:original-trace}
\end{equation}
In particular, if \(a_K=b_K\), the aligned vectors in
\eqref{stat:aligned-vectors} are identical, so
\begin{equation}
 C_{ab}=\delta_{ab}\quad\text{whenever }a_K=b_K.
 \label{stat:erased-trace}
\end{equation}

For each retained block, its independent normalized columns have
positive definite Gram matrix
\[
 H_i=V_i^*V_i=\begin{pmatrix}1&\mu_i\\\bar\mu_i&1\end{pmatrix},
 \qquad |\mu_i|<1.
\]
Both \(V_i^*\) and
\[
 H_i^{-1}=\frac1{1-|\mu_i|^2}
       \begin{pmatrix}1&-\mu_i\\-\bar\mu_i&1\end{pmatrix}
\]
are tensor-LY, the latter by Lemma~\ref{stat:two-by-two}.
Consequently \(L_i=H_i^{-1}V_i^*\) is tensor-LY and satisfies
\(L_iV_i=I\).  Contraction cannot give zero because of this
identity.  The rectangular matrix
\(L=\langle\psi|\otimes\bigotimes_{i\in K}L_i\) is stable
and obeys \(Lv_a=|a_K\rangle\).

The CP map \(\Psi(M)=L\Phi(M)L^*\) is full Choi--LY or zero,
and satisfies
\[
 \Psi(|a\rangle\langle b|)=C_{ab}|a_K\rangle\langle b_K|.
\]
Its trace is \(C_{ab}\delta_{a_K,b_K}=\delta_{ab}\) by
\eqref{stat:erased-trace}.  Thus \(\Psi\) is TP and nonzero.
The compression itself need not be TP on arbitrary output
matrices; this argument proves TP only for its composition with
\(\Phi\), which is all that is required.

The Choi polynomial of \(\Psi\) is
\[
 \sum_{a,b}C_{ab}z^{a_K}w^{b_K}u^av^b.
\]
For \(i\in K\), replace \(z_iu_i,w_iv_i\) by \(r_i,s_i\);
for \(i\in E\), put \(r_i=u_i,s_i=v_i\).
The variables \(r_i,s_i\) range independently over the disk,
so stability of this Choi polynomial is precisely tensor
stability of the entire correlation matrix \(C\).
Theorem~\ref{stat:schur} gives
\(C=\bigotimes_i\begin{psmallmatrix}1&\eta_i\\\bar\eta_i&1\end{psmallmatrix}\),
with \(|\eta_i|\le1\).  Taking strings differing only at an
erased input in \eqref{stat:erased-trace} forces \(\eta_i=0\)
there.  Taking strings differing at a retained input in
\eqref{stat:original-trace} gives
\(\eta_i\bar\mu_i=0\), equivalently
\eqref{stat:overlap-noise}.  Combining these tensor-product
coefficients with \eqref{stat:aligned-vectors} proves
\eqref{stat:pure-basis-form}; the erased coefficient is exactly
\(\delta_{a_E,b_E}\), hence ordinary partial trace.

Conversely, start from the displayed data.  Each \(C_i\) is a
PSD correlation matrix, so its Schur map is CP.  Following it
by \(M\mapsto V_iMV_i^*\) preserves CP.  Normalization of
the columns gives the correct diagonal traces, while
\eqref{stat:overlap-noise} makes the off-diagonal traces zero.
Thus each \(\mathcal B_i\) is TP and has pure basis outputs.
Its two component Choi tensors are stable by
Lemma~\ref{stat:reflected-pair} and
Theorem~\ref{stat:schur}; their composition is nonzero because
it is TP.  Stable pure preparation, partial trace, tensor
products and physical permutations prove the converse.
The descriptions of the two cases follow directly from
\(\mu_i\eta_i=0\).
\end{proof}

Theorems~\ref{stat:visible}, \ref{stat:schur}, and
\ref{stat:pure-basis} concern different singular strata.
In particular, none asserts a classification of general channels
with mixed computational-basis outputs.  Proposition~\ref{stat:no-dilation}
prevents extending the isometry theorem to that class by an
assumption of a stable Stinespring representation.

\subsection{Pauli-diagonal maps without invertibility}

\begin{lemma}[Factorization on slit domains]
\label{stat:slit-factorization}
Let \(\Omega_1,\ldots,\Omega_n\) be nonempty connected open
subsets of \(\C\) with \(\C\setminus\Omega_j\subseteq\R\).
A multiaffine polynomial nonvanishing on \(\prod_j\Omega_j\)
is a nonzero scalar times a product of one-variable affine
polynomials, with constant factors permitted.
\end{lemma}

\begin{proof}
Write \(p(r_1,r')=a(r')+r_1b(r')\).  If \(b=0\), omit that
coordinate.  Otherwise, on a nonempty connected open neighborhood
where \(b\ne0\), the holomorphic function \(-a/b\) avoids
\(\Omega_1\) and is therefore real-valued.  It is constant,
say \(c\notin\Omega_1\).  The polynomial identity principle
gives \(a+cb=0\) everywhere, so \(p=(r_1-c)b\).
The remaining factor is nonvanishing on the remaining domain.
Induction proves the assertion.
\end{proof}

\begin{theorem}[All-rank Pauli-diagonal rigidity]
\label{stat:pauli}
Suppose a complex-linear HP, TP map on \(n\) qubits is diagonal
in the Pauli-string basis:
\(\Phi(\sigma_\alpha)=\lambda_\alpha\sigma_\alpha\), where
\(\sigma_\alpha=\bigotimes_j\sigma_{\alpha_j}\) and
\(\sigma_0=I\).  If \(\Phi\) is full Choi--LY, then
\begin{equation}
 \lambda_\alpha=\prod_j\lambda_{\alpha_j^{(j)}},
 \qquad \Phi=\bigotimes_j\Phi_j,
 \label{stat:pauli-product}
\end{equation}
where \(\alpha_j^{(j)}\) denotes the string with the stated
Pauli at site \(j\) and identities elsewhere.  Each local
factor is HP, TP and full Choi--LY.  If \(\Phi\) is CP, all
factors are CP, and the converse holds.  Zero eigenvalues are
allowed.
\end{theorem}

\begin{proof}
HP makes the \(\lambda_\alpha\) real, and TP gives
\(\lambda_0=1\).  Pauli orthogonality yields
\[
 J(\Phi)=2^{-n}\sum_\alpha\lambda_\alpha
                 \sigma_\alpha\otimes\sigma_\alpha^{\mathrm T}.
\]
The four one-site polynomials are
\(f_I=1+zw\), \(f_X=z+w\), \(f_Y=i(z-w)\),
and \(f_Z=1-zw\); the transpose contributes a minus sign
only for \(Y\).

At every site select a desired axis, independently of the others.
For axis \(X\), set the output variables to \(z=w=1\).
Only the identity and \(X\) terms remain, and their normalized
ratio is
\[
 r_X=\frac{x+y}{1+xy}.
\]
Its image on the bidisk is
\(\Omega_\perp=\C\setminus(({-\infty},-1]\cup[1,\infty))\).
For a real \(|r|\ge1\), the inverse equation would require
\(y=(r-x)/(1-rx)\), whereas
\[
 |r-x|^2-|1-rx|^2=(r^2-1)(1-|x|^2)\ge0.
\]
The endpoint equations also require a boundary variable.
Conversely, setting \(x=y=u\) gives an inverse quadratic with
reciprocal roots.  A unit-circle root would place \(r\) on
the omitted rays.  For every nonzero \(r\in\Omega_\perp\),
exactly one inverse root is therefore in the disk; zero is
attained at \(u=0\).

For axis \(Y\), set \(z=i,w=-i\).  The output factor is
\(f_Y=-2\), and the input transpose changes the sign again.
The surviving ratio is
\(r_Y=i(x-y)/(1+xy)\), with the same image
\(\Omega_\perp\), by replacing \((x,y)\) with \((ix,-iy)\).
For axis \(Z\), set \(z=s,w=-s,x=t,y=t\), all interior.
The output specialization kills \(X\) and the input
specialization kills \(Y\).  The ratio is
\[
 r_Z=\frac{1+s^2}{1-s^2}\frac{1-t^2}{1+t^2}.
\]
Each factor ranges over the right half-plane, so their product
ranges exactly over \(\Omega_Z=\C\setminus(-\infty,0]\).
Indeed products have arguments strictly between \(-\pi\) and
\(\pi\), and the principal square root realizes every point
of this slit plane as such a product.

For axes \(a_1,\ldots,a_n\), divide out the nonzero identity
factors after these substitutions.  The resulting polynomial is
\begin{equation}
 p_a(r)=\sum_{S\subseteq[n]}\lambda_{a_S}\prod_{j\in S}r_j,
 \qquad p_a(0)=1,
 \label{stat:axis-polynomial}
\end{equation}
where \(a_S\) has the chosen axes on \(S\).
Boundary specialization at the \(X,Y\) sites gives stable or
zero by Hurwitz.  The zero alternative is impossible: the
ratio maps have the nonempty open product image just computed,
and the nonzero polynomial \eqref{stat:axis-polynomial} cannot
vanish identically there.  The \(Z\) substitutions are interior.
Thus \(p_a\) is nonvanishing on the corresponding product of
slit domains.

Lemma~\ref{stat:slit-factorization} and \(p_a(0)=1\) give
\(p_a(r)=\prod_j(1+\lambda_{a_j^{(j)}}r_j)\).
The normalization at zero is algebraic and does not require
zero to belong to \(\Omega_Z\).  Comparing coefficients, and
varying the axes, proves \eqref{stat:pauli-product} for every
Pauli string.  Local factors inherit the Choi and positivity
properties by the marginal construction in the proof of
Theorem~\ref{stat:visible}.  Tensor products prove the converse.
\end{proof}

For CPTP local Pauli factors the exact one-site condition is
\[
 \lambda_z\ge\max\{|\lambda_x|,|\lambda_y|\},
\]
or equivalently
\(\min(p_0,p_z)\ge\max(p_x,p_y)\) in the Pauli-error
probabilities.  This is the Pauli specialization of the local
classification in Theorem~\ref{loc:classification}; its
sufficiency is also \cite[Theorem~3]{WBG26}.
Because the Pauli transform between error probabilities and
eigenvalues is an invertible tensor-product transform,
\eqref{stat:pauli-product} says that the entire error
distribution is a product distribution.

\subsection{Classical channels and the purity boundary}

Consider the channel that measures \(m\) input bits in the
computational basis and prepares a diagonal distribution
\(P(b\mid a)\) on \(n\) output bits.  Its full Choi
polynomial reduces to
\begin{equation}
 K(r,s)=\sum_{a,b}P(b\mid a)r^bs^a=\sum_aA_a(r)s^a,
 \qquad K(1,s)=\prod_i(1+s_i).
 \label{stat:classical-kernel}
\end{equation}
The coefficients are nonnegative.  As above, reducing each
row/column pair to its product is an equivalence of disk
stability.  Define the collision purity of column \(a\) by
\(\mathcal C(a)=\sum_bP(b\mid a)^2=\|A_a\|_2^2\).

\begin{theorem}[The constant-purity boundary]
\label{stat:classical-purity}
Every stable stochastic kernel \eqref{stat:classical-kernel}
satisfies \(\mathcal C(a)\le\mathcal C(0)\) for every input.
Equality at every singleton input \(e_i\) holds if and only
if, for a partition \(S_0,S_1,\ldots,S_m\) of the output
variables, the kernel is
\begin{equation}
 K(r,s)=C(r_{S_0})\prod_{i=1}^m
                 \bigl(p_i(r_{S_i})+s_i p_i^*(r_{S_i})\bigr).
 \label{stat:classical-purity-form}
\end{equation}
Here \(C\) and the \(p_i\) are stable, multiaffine, have
nonnegative coefficients, satisfy \(C(1)=p_i(1)=1\), and
\(\gcd(p_i,p_i^*)=1\); reflection is over the full indicated
block.  Empty blocks \(S_i\) are allowed and contribute
\(1+s_i\).  Every column of such a kernel has the same
collision purity.
\end{theorem}

\begin{proof}
Fix the output variables in the polydisk.  For any input subset
\(T\), set its variables equal to \(t\) and the other input
variables to zero.  The resulting univariate polynomial is
nonvanishing for \(|t|<1\), has constant coefficient \(A_0\),
and top coefficient \(A_T\) in degree \(|T|\).  If that top
coefficient is nonzero, the product of root moduli is at least
one, whence \(|A_T|\le|A_0|\).  If it is zero the inequality
is immediate.  Continuity and torus Parseval prove the purity
inequality.

Under singleton equality, Lemma~\ref{stat:saturation} gives
disjoint primitive factors \(C\prod_i(p_i+s_iq_i)\).
The stable constant slice makes the constant coefficients of
\(C,p_i\) nonzero.  Fix their phases to make these constants
positive.  Each coefficient of each factor can be isolated
as a coefficient of \(K\) by selecting the constant monomial
in all the other factors; disjoint supports therefore make
all coefficients of \(C,p_i,q_i\) nonnegative.
Lemma~\ref{stat:reflected-pair} gives
\(q_i=e^{i\theta_i}p_i^*\), and nonnegativity forces the
phase to be one.  Stochasticity at output variables one gives
\(p_i(1)=q_i(1)>0\).  Normalizing these common values and
the prepared factor gives \eqref{stat:classical-purity-form},
with no scalar remainder because every column has mass one.
Variables absent from \(K\) are fixed-zero prepared outputs.

Conversely, Lemma~\ref{stat:reflected-pair} makes each factor
stable.  The displayed normalization makes the nonnegative
coefficient kernel stochastic.  Reflection preserves
coefficient norms and disjoint supports make them multiply,
so all columns have equal collision purity.
\end{proof}

\begin{corollary}[Deterministic classical channels]
\label{stat:deterministic}
A deterministic stochastic kernel is full Choi--LY if and
only if every output bit is identically zero or copies one
input bit.  An input may be copied to any disjoint block of
outputs or discarded.
\end{corollary}

\begin{proof}
Write \(P(b\mid a)=\mathbf{1}_{b=f(a)}\).  All column
purities equal one.  The rational product identity in
Lemma~\ref{stat:saturation} therefore gives
\[
 r^{f(a)}=r^{f(0)+\sum_{i:a_i=1}(f(e_i)-f(0))}.
\]
Stable evaluation at the full origin requires \(f(0)=0\).
Each output coordinate is consequently a sum of its
singleton Boolean values.  Since it must still be a bit
when all inputs are one, at most one singleton value can
be one.  This proves necessity.  Conversely the kernel of
such a map is
\(\prod_i(1+s_i\prod_{j\in S_i}r_j)\), with disjoint
copied-output blocks, and is manifestly stable.
\end{proof}

\begin{theorem}[One classical output]
\label{stat:classical-one-output}
A stochastic kernel from \(m\) input bits to one output bit
is full Choi--LY if and only if
\begin{equation}
 P(1\mid a)=\alpha+\sum_i\lambda_i a_i,
 \qquad \alpha\ge0,\quad\lambda_i\ge0,\quad
                   2\alpha+\sum_i\lambda_i\le1.
 \label{stat:randomized-copy}
\end{equation}
Thus it chooses to copy input \(i\) with probability
\(\lambda_i\), to output zero with probability
\(1-2\alpha-\sum_i\lambda_i\), or to output a fresh fair
bit with probability \(2\alpha\).
\end{theorem}

\begin{proof}
Write \(K(r,s)=A(s)+rB(s)\).  Stochasticity gives
\(A+B=\prod_i(1+s_i)\).  Put \(H=(A-B)/(A+B)\), so
\[
 K=\frac{\prod_i(1+s_i)}2\bigl((1+r)+(1-r)H(s)\bigr).
\]
The quotient \((1+r)/(1-r)\) ranges over the open right
half-plane; thus stability is exactly \(\Re H(s)\ge0\)
on the input polydisk.  Set
\(u_i=(1-s_i)/(1+s_i)\).  Then \(H\) becomes a real
multiaffine polynomial in the right-half-plane variables
\(u_i\), because every
\(s^a/\prod_i(1+s_i)\) is a product of factors
\((1\pm u_i)/2\).

A real polynomial with nonnegative real part on a product
of right half-planes has the form
\(c+\sum_i\lambda_i u_i\), with \(c,\lambda_i\ge0\).
To see this, restrict to \(u_i=v_i t\), \(v_i>0\),
\(\Re t>0\).  A real one-variable polynomial of degree
at least two has a leading term with negative real part
along some interior right-half-plane ray at large radius.
Every such restriction therefore has degree at most one.
Varying \(v\) in its open positive orthant forces all
higher homogeneous parts to vanish.  Positive real
evaluations near zero, or with one variable tending to
infinity, show that the remaining coefficients are
nonnegative.

Evaluating the transformed polynomial identity at
\(u_i=(-1)^{a_i}\) gives
\(1-2P(1\mid a)=c+\sum_i\lambda_i(-1)^{a_i}\).
This is polynomial evaluation, including at vertices
corresponding to poles of the original change of variables.
Consequently \eqref{stat:randomized-copy} holds with
\(\alpha=(1-c-\sum_i\lambda_i)/2\).  The original
probability at \(a=0\) makes \(\alpha\ge0\), and
\(c\ge0\) gives the last inequality.  Conversely these
conditions define the stated stochastic mixture and a
positive-real polynomial \(H\), proving stability.
\end{proof}

Every single-output marginal of a multibit classical
full Choi--LY channel therefore has the form
\eqref{stat:randomized-copy}.  Marginalization sets the
other output variables to one; boundary Hurwitz applies,
and stochastic normalization excludes zero.  This does
not determine the joint output correlations.

\begin{proposition}[Nonconvexity among reset channels]
\label{stat:reset-nonconvex}
For \(n\ge3\), the distribution
\(\nu_\theta=(1-\theta)\delta_{0^n}
+\theta\operatorname{Unif}(\{0,1\}^n)\),
\(0\le\theta\le1\), has a stable polynomial precisely when
\begin{equation}
 0\le\theta\le\frac1{1+\cos^n(\pi/n)}
 \quad\text{or}\quad\theta=1.
 \label{stat:reset-threshold}
\end{equation}
For \(n=1,2\), all \(\theta\in[0,1]\) are allowed.
Consequently the set of full Choi--LY channels is not
convex, even among classical reset channels.
\end{proposition}

\begin{proof}
The distribution polynomial is
\(f_\theta(r)=1-\theta+2^{-n}\theta\prod_j(1+r_j)\).
The endpoints \(\theta=0,1\) are stable.  For
\(0<\theta<1\), an interior zero means that a product
of \(n\) points in the open disk centered at one with
radius one equals \(-2^n(1-\theta)/\theta\).
Write these nonzero points as
\(\rho_j e^{i\alpha_j}\), where
\(|\alpha_j|<\pi/2\) and
\(0<\rho_j<2\cos\alpha_j\).
A negative real product requires
\(\sum_j\alpha_j\) to be an odd multiple of \(\pi\),
which is impossible for \(n=1,2\).
For \(n\ge3\), concavity of \(\log\cos\) implies
\[
 \prod_j\rho_j
 <2^n\prod_j\cos\alpha_j
 \le[2\cos(\pi/n)]^n.
\]
An odd multiple of larger absolute value can only reduce
the bound.  Every smaller positive product magnitude is
attained by setting all arguments to \(\pi/n\) and
choosing equal moduli.  Therefore an interior zero occurs
exactly when
\(2^n(1-\theta)/\theta<[2\cos(\pi/n)]^n\), the
complement of \eqref{stat:reset-threshold} away from the
separately treated endpoint.  Equality requires boundary
points and is stable for the open-disk definition.
Resetting to this distribution multiplies its polynomial
by the stable input trace factor \(\prod_i(1+s_i)\),
so the same criterion applies to the channel.
\end{proof}


\section{The complete one-qubit static classification}\label{loc:section}

We give exact parameter criteria for a single qubit, including every
singular stratum.  The stability classification first requires only
Hermiticity preservation and trace preservation.  Complete positivity is
then imposed by explicit matrix or scalar inequalities.

Write a complex-linear Hermiticity-preserving, trace-preserving map in
Bloch coordinates as
\begin{equation}\label{loc:bloch}
 \Phi(I)=I+t\cdot\sigma,\qquad
 \Phi(\sigma_j)=\sum_iT_{ij}\sigma_i,
 \qquad t\in\R^3,\quad T\in\R^{3\times3},
\end{equation}
where $\sigma=(X,Y,Z)$.  On trace-one Hermitian matrices its action is
$r\mapsto t+Tr$.  Invertibility means complex-linear invertibility on
$M_2(\C)$, equivalently $\det T\ne0$.  Full Choi stability always
refers to all four variables of $F_{J(\Phi)}$.

\begin{theorem}[All one-qubit HP/TP maps]\label{loc:classification}
The map \eqref{loc:bloch} has a full Lee--Yang Choi kernel if and only if
it belongs to exactly one of the following two classes.
\begin{enumerate}
\item There is a nonzero real $2\times2$ matrix $D$ such that
\begin{equation}\label{loc:transverse-form}
 T=\begin{pmatrix}D&0\\0&\lambda\end{pmatrix},\qquad
 \lambda>0,\qquad t=(\tau,0),\qquad \tau\in\R^2,\quad\|\tau\|<1,
\end{equation}
and, for every $n\in\R^2$ with $\|n\|=1$,
\begin{equation}\label{loc:angular}
 \|D^{\mathsf T}n\|^2\leq m_\lambda(\tau\cdot n),\qquad
 m_\lambda(r)=\min_{-1\leq x\leq1}
                  \bigl((x+r)^2+\lambda^2(1-x^2)\bigr).
\end{equation}
\item The transverse linear action vanishes and
\begin{equation}\label{loc:erasure-form}
 T=\diag(0,0,\lambda),\qquad
 \lambda\geq0,\quad t_z\geq0,\quad\|t_\perp\|\leq1.
\end{equation}
\end{enumerate}
For $|r|<1$, the minimum in \eqref{loc:angular} is
\begin{equation}\label{loc:minimum}
 m_\lambda(r)=
 \begin{cases}
 \lambda^2\bigl(1-r^2/(1-\lambda^2)\bigr),
   &0<\lambda<1,\ |r|\leq1-\lambda^2,\\
 (1-|r|)^2,&0<\lambda<1,\ |r|>1-\lambda^2,\\
 (1-|r|)^2,&\lambda\geq1.
 \end{cases}
\end{equation}
The complex-linear rank is $2+\rank D$ in the first class and
$1+\boldsymbol{1}_{\{\lambda>0\}}$ in the second.  In particular the
first class is invertible exactly when $D$ is invertible.
\end{theorem}

\subsection{Boundary kernels and the proof of the classification}

\begin{lemma}[Trace-one scalar kernels]\label{loc:scalar-kernel}
A Hermitian trace-one matrix with transverse Bloch norm $R$ and
longitudinal coordinate $h$ has a stable coefficient polynomial if and
only if $R\leq1$ and $h\geq0$.  For $0\leq R\leq1$, the function
\begin{equation}\label{loc:positive-real}
 K_R(z,w)=\frac{1+zw+R(z+w)}{1-zw}
\end{equation}
has strictly positive real part on $\D^2$, and the infimum of that real
part is zero.
\end{lemma}

\begin{proof}
A rotation about $Z$ makes the transverse component $(R,0)$ and only
rotates the disk variables.  The coefficient polynomial becomes
\[
 f(z,w)=\frac{1+h+R(z+w)+(1-h)zw}{2}
       =\frac{1-zw}{2}(K_R(z,w)+h).
\]
For $a=|z|<1$, $b=|w|<1$ and $0\leq R\leq1$, multiplication by
$|1-zw|^2$ shows that the real part of the numerator of $K_R$ is at
least
\begin{align*}
 1-a^2b^2-R\bigl(a(1-b^2)+b(1-a^2)\bigr)
 &\geq(1-ab)(1-a)(1-b)>0.
\end{align*}
This proves sufficiency.  Setting $z=w=is$ and letting $s\uparrow1$
shows that the infimum of $\Re K_R$ is zero.

For necessity, stability gives the boundary root-modulus inequality
obtained by treating $f$ as affine in $w$.  At $z=e^{i\theta}$ its
squared difference between the constant and linear coefficients is
$h(1+R\cos\theta)$.  Its average is $h$, so $h\geq0$.  If $h>0$,
the same inequality for all $\theta$ gives $R\leq1$.  If $h=0$ and
$R>1$, setting $w=0$ gives the interior zero $z=-1/R$.
\end{proof}

\begin{lemma}[The local block form]\label{loc:block}
Every map in Theorem~\ref{loc:classification} whose full Choi kernel is
stable satisfies
\begin{equation}\label{loc:block-intermediate}
 T=D\oplus\lambda,\qquad\lambda\geq0.
\end{equation}
If $D\ne0$, then $t_z=0$ and $\lambda>0$.
\end{lemma}

\begin{proof}
Let $\phi=\Phi^\dagger$.  It is unital and Hermiticity preserving.
Lemma~\ref{stat:bridge} makes it a preserver of complex matrix-LY
polynomials, with a possible zero output.  Near the identity that zero
alternative is excluded by continuity.  Apply
Lemma~\ref{stat:edge} to the two-sided stable curves
$\phi(e^{sX})$, $\phi(e^{sY})$, and $\phi(e^{siZ})$.
Hermiticity preservation, and its anti-Hermitian counterpart, give
\begin{equation}\label{loc:adjoint-block}
 \begin{aligned}
 \phi(X)&=t_xI+D_{11}X+D_{12}Y,\\
 \phi(Y)&=t_yI+D_{21}X+D_{22}Y,\\
 \phi(Z)&=t_zI+\lambda Z.
 \end{aligned}
\end{equation}
The stable matrix $I+sZ$, $s>0$ small, maps to
$(1+st_z)I+s\lambda Z$.  Its root-modulus inequality is
\[
 |1+st_z+s\lambda|^2-|1+st_z-s\lambda|^2
       =4s\lambda(1+st_z)\geq0,
\]
so $\lambda\geq0$.

If $D\ne0$, choose a real transverse field $A$ such that the transverse
component of $\phi(A)$ is nonzero.  For small nonzero real
$\varepsilon$, both $A+i\varepsilon Z$ and the traceless part of its
image are hyperbolic.  The input pencil
$I+s(A+i\varepsilon Z)$ is stable for both signs of sufficiently
small $s$: its exponentials are stability units by
Lemma~\ref{stat:edge}, and the identity
$(A+i\varepsilon Z)^2\in\R I$ rewrites these exponentials as scalar
multiples of the pencil near the identity.  Its image is a stable pencil
whose scalar part is $t_A+i\varepsilon t_z$.
Lemma~\ref{stat:hyperbolic} forces this scalar to be real, hence $t_z=0$.

Finally specialize the two input Choi variables to
$u=e^{i\theta}$, $v=-e^{-i\theta}$.  The input matrix is
$Z+i n\cdot\sigma_\perp$, with $n$ ranging over the real unit circle.
The output matrix is
$\lambda Z+i(Dn)\cdot\sigma_\perp$.  Hurwitz's theorem makes its
polynomial stable or zero.  If $\lambda=0$, choose $n$ with $Dn\ne0$.
The polynomial is then nonzero but its constant coefficient is zero,
contradicting stability.  This proves $\lambda>0$.
\end{proof}

\begin{proof}[Proof of Theorem~\ref{loc:classification}]
By Lemma~\ref{loc:block}, it remains to treat the two possible transverse
blocks.  Suppose first that $D=0$, and rotate about $Z$ so that
$t_\perp=(R,0)$.  Direct expansion of the Choi kernel gives
\begin{equation}\label{loc:erasure-kernel}
 F_{J(\Phi)}=
 \frac{(1+uv)(1-zw)}2
 \left(K_R(z,w)+t_z+\lambda\frac{1-uv}{1+uv}\right).
\end{equation}
The last quotient ranges over the right half-plane.  Thus
$R\leq1$, $t_z\geq0$, and $\lambda\geq0$ imply strict nonvanishing
by Lemma~\ref{loc:scalar-kernel}.
Conversely, specialization $u=v=0$ gives a trace-one Hermitian kernel
with transverse norm $R$ and longitudinal coordinate $t_z+\lambda$,
so $R\leq1$.  If $\lambda=0$, that same lemma gives $t_z\geq0$.
If $\lambda>0$ and $t_z<0$, choose $z,w$ with
$\Re(K_R+t_z)<0$.  A value of $(1-uv)/(1+uv)$ in the right
half-plane then cancels the parenthesis in
\eqref{loc:erasure-kernel}, a contradiction.  This proves the second
class completely, without positivity.

Now suppose $D\ne0$.  Lemma~\ref{loc:block} gives
$\lambda>0$ and $t_z=0$.  Set
\begin{equation}\label{loc:coefficients}
 \begin{gathered}
 A_0=(1+\lambda)/2,\qquad B_0=(1-\lambda)/2,
 \qquad \zeta=(t_x-it_y)/2,\\
 a=(D_{11}+D_{22}+i(D_{12}-D_{21}))/2,\qquad
 b=(D_{11}-D_{22}+i(D_{12}+D_{21}))/2.
 \end{gathered}
\end{equation}
The full polynomial has the reciprocal decomposition
\begin{equation}\label{loc:reciprocal}
 F_{J(\Phi)}=p+vp^\sharp,\qquad
 p=g+uh,\quad
 g=A_0+\zeta w+\bar\zeta z+B_0zw,\quad
 h=\bar b w+\bar a z,
\end{equation}
where
$p^\sharp(z,w,u)=zwu\,\overline{p(1/\bar z,1/\bar w,1/\bar u)}$.
This is an algebraic identity with no positivity assumption.

If the full polynomial is stable, its specializations $p$ and $g$ are
stable, since their constant coefficient is $A_0>0$.
Lemma~\ref{loc:scalar-kernel} applied to $g$ gives
$R=\|t_\perp\|\leq1$.  The root test in $u$ and continuity give
$|h|\leq|g|$ on the output two-torus.  Write
$z=e^{i(\theta+\varphi)}$, $w=e^{i(\theta-\varphi)}$ and
$n=(\cos\varphi,-\sin\varphi)$.  Direct expansion yields
\begin{equation}\label{loc:torus-norms}
 |h|^2=\|D^{\mathsf T}n\|^2,\qquad
 |g|^2=(\cos\theta+t_\perp\cdot n)^2+
                         \lambda^2\sin^2\theta.
\end{equation}
Minimization in $x=\cos\theta$ proves \eqref{loc:angular}.
If $R=1$, rotate so that $t_\perp=e_1$.  The endpoint $x=-1$
gives $\|D^{\mathsf T}n\|\leq1-n_1$ for every unit $n$.
At $n=e_1$ the first row of $D$ vanishes.  Setting
$n=(\cos\alpha,\sin\alpha)$ then gives
$|\sin\alpha|\|D^{\mathsf T}e_2\|\leq1-\cos\alpha$.
Division by $|\sin\alpha|$ and passage to $\alpha=0$ kills the
second row.  Since $D\ne0$, necessarily $R<1$.

Conversely, assume \eqref{loc:transverse-form} and
\eqref{loc:angular}.  The polynomial $g$ has no zero on the closed
bidisk.  Indeed its constant coefficient in $w$, namely
$A_0+\bar\zeta z$, has no zero on the closed disk, and at $|z|=1$
the squared root-modulus difference is
\[
 |A_0+\bar\zeta z|^2-|\zeta+B_0z|^2
       =\lambda(1+2\Re(\bar\zeta z))>0.
\]
Here $2|\zeta|=R<1$.  The maximum modulus principle in $z$ proves
the asserted closed-bidisk nonvanishing.  Equations
\eqref{loc:angular} and \eqref{loc:torus-norms}, followed by the
maximum modulus principle in each variable, imply $|h/g|\leq1$ on
the bidisk.  Hence $p=g+uh$ is stable.

For any stable multiaffine $p$, one has $|p^\sharp/p|\leq1$ in the
open polydisk.  To include boundary zeros, apply the maximum modulus
principle to $p_r(z,w,u)=p(rz,rw,ru)$, $r<1$, which has no zeros on
the closed polydisk.  On its distinguished boundary
$|p_r^\sharp|=|p_r|$, so $|p_r^\sharp/p_r|\leq1$ inside.
Let $r\uparrow1$.  This proves the assertion, and hence
$p+vp^\sharp\ne0$ when $|v|<1$.  Sufficiency follows.

The formula \eqref{loc:minimum} follows by minimizing the quadratic
$(1-\lambda^2)x^2+2rx+r^2+\lambda^2$: for $\lambda<1$ its critical
point is $-r/(1-\lambda^2)$, and otherwise its minimum on the interval
is at an endpoint.  Finally the Pauli transfer matrix is
$\begin{psmallmatrix}1&0\\t&T\end{psmallmatrix}$, whose rank is
$1+\rank T$.  This gives all rank assertions.
\end{proof}

\subsection{Complete positivity and finite inequalities}

\begin{theorem}[The CPTP subfamilies]\label{loc:cptp}
In the first class of Theorem~\ref{loc:classification}, complete
positivity is equivalent to
\begin{equation}\label{loc:choi-psd}
 J=\begin{pmatrix}
 A_0&0&\zeta&a\\
 0&B_0&\bar b&\zeta\\
 \bar\zeta&b&B_0&0\\
 \bar a&\bar\zeta&0&A_0
 \end{pmatrix}\succeq0,
\end{equation}
with coefficients \eqref{loc:coefficients}.  For a positive
trace-preserving map of this form with $0<\lambda<1$, the stability
criterion \eqref{loc:angular} is equivalent to the finite matrix
inequality
\begin{equation}\label{loc:lmi}
 \frac{DD^{\mathsf T}}{\lambda^2}
       +\frac{\tau\tau^{\mathsf T}}{1-\lambda^2}\preceq I_2.
\end{equation}
At $\lambda=1$, positivity forces $\tau=0$, and the stability
criterion is $\|D\|_{\mathrm{op}}\leq1$.

For the singular rank-three subclass, input and output rotations about
$Z$ put $D=\diag(d,0)$ with $d>0$ and $\tau=(p,q)$.
The complete CPTP and full-Choi-LY conditions are the two scalar
inequalities
\begin{equation}\label{loc:rank-three}
 q^2+\bigl(d+\sqrt{p^2+\lambda^2}\bigr)^2\leq1,
 \qquad
 \frac{d^2}{\lambda^2}
       +\frac{p^2}{1-\lambda^2-q^2}\leq1,
 \qquad d,\lambda>0.
\end{equation}
The first inequality makes the denominator in the second strictly
positive.

In the transverse-erasure class, the complete conditions are
\begin{equation}\label{loc:erasure-cptp}
 t_z,\lambda\geq0,\qquad
 \|t_\perp\|^2+(t_z+\lambda)^2\leq1.
\end{equation}
\end{theorem}

\begin{proof}
The matrix in \eqref{loc:choi-psd} is the Choi matrix in the order
$(00,01,10,11)$, with each pair ordered as (output,input).
Its partial trace over the output equals the input identity.  Choi's
criterion gives the first assertion.

For the finite inequality, let $k=1-\lambda^2>0$ and fix a transverse
unit vector $n$.  Reversing $n$ if necessary, set
$r=\tau\cdot n\geq0$ and $s^2=\|D^{\mathsf T}n\|^2$.
The desired scalar form of \eqref{loc:lmi} is
\begin{equation}\label{loc:lmi-scalar}
 Q=\lambda^2-s^2-\lambda^2r^2/k\geq0.
\end{equation}
If $r\leq k$, this is exactly the first branch of
\eqref{loc:minimum}.  If $r>k$, put $a=k/r\in(0,1)$.
Positivity means that the affine image of the Bloch unit ball lies in
the unit ball.  Its support in the unit output direction
$(an,\sqrt{1-a^2})$ therefore satisfies
\[
 ar+\sqrt{a^2s^2+\lambda^2(1-a^2)}\leq1.
\]
Squaring the nonnegative sides gives
\[
 0\leq k-2ar+a^2(r^2-s^2+\lambda^2)=\frac{k^2}{r^2}Q.
\]
Thus positivity itself implies \eqref{loc:lmi-scalar} in the clipped
region.  Conversely \eqref{loc:lmi-scalar} bounds the unconstrained
minimum over all real $x$, hence implies \eqref{loc:angular} on
$[-1,1]$.  Varying $n$ proves \eqref{loc:lmi}.  Positivity at the
input Bloch points $\pm e_z$ gives
$\|\tau\|^2+\lambda^2\leq1$, proving both $\lambda\leq1$ and the
stated endpoint assertion.

Suppose now $D$ has rank one.  Real rotations in the transverse input
and output planes put it in the form $\diag(d,0)$, $d>0$; these are
implemented by rotations about $Z$ and preserve both properties under
consideration.  In these coordinates the Choi matrix satisfies
$2J=I+H$, where
\[
 H=pX\otimes I+qY\otimes I+dX\otimes X+\lambda Z\otimes Z.
\]
Pauli multiplication gives
\[
 H^2=(p^2+q^2+d^2+\lambda^2)I
          +2pd\,I\otimes X-2d\lambda\,Y\otimes Y.
\]
The last two Pauli matrices anticommute.  The eigenvalues of $H^2$ are
$q^2+(\sqrt{p^2+\lambda^2}\pm d)^2$, each twice.  Also
$\Tr H=\Tr H^3=0$, by Pauli orthogonality.  Newton's identities make
the characteristic polynomial of $H$ even, so its spectrum is
symmetric about zero.  Its least eigenvalue is consequently
$-\sqrt{q^2+(\sqrt{p^2+\lambda^2}+d)^2}$.
This proves the first inequality in \eqref{loc:rank-three} as the
exact CP condition.  Since $d,\lambda>0$, it also implies
$\lambda<1$ and $1-\lambda^2-q^2>0$.
Taking the Schur complement of the lower-right entry in
\eqref{loc:lmi} now gives precisely the second inequality in
\eqref{loc:rank-three}.

Finally, a transverse-erasure map measures in the $Z$ basis and prepares
\begin{equation}\label{loc:prepared-states}
 \rho_\pm=\tfrac12\bigl(I+t_xX+t_yY+(t_z\pm\lambda)Z\bigr).
\end{equation}
Its Choi matrix is block diagonal on the input and is positive exactly
when both prepared matrices are positive, equivalently
$\|t_\perp\|^2+(t_z\pm\lambda)^2\leq1$.
Since $t_z,\lambda\geq0$, the two inequalities reduce to
\eqref{loc:erasure-cptp}.
\end{proof}

\begin{remark}\label{loc:singular-distinction}
Full Choi stability does not require both states in
\eqref{loc:prepared-states} to be individually Lee--Yang: the quantity
$t_z-\lambda$ may be negative.  The choice
$t_\perp=0$, $t_z=\lambda=1/3$ is a rank-two CPTP example with
positive longitudinal drift.  A rank-three example retaining transverse
information is $t=0$, $D=\diag(1/5,0)$, $\lambda=2/5$.
At rank one, the maps are precisely resets
$M\mapsto\Tr(M)\rho$; their CPTP full-Choi-LY condition is
$t_z\geq0$ and $\|t\|\leq1$.
These singular examples explain why the zero longitudinal drift conclusion for an
invertible map cannot be extended without a rank hypothesis.
\end{remark}


\section{Independently covariant channels between finite spins}\label{spin:section}

Let $U_k$ denote the spin-$k/2$ representation on
$\mathcal H_k=\C^{k+1}$, now allowing $k=0$.
For input capacities $\kappa=(\kappa_1,\ldots,\kappa_n)$ and output
capacities $\ell=(\ell_1,\ldots,\ell_n)$, a complex-linear map
$\Phi:\operatorname{End}(\bigotimes_i\mathcal H_{\kappa_i})
\to\operatorname{End}(\bigotimes_i\mathcal H_{\ell_i})$
is independently covariant if
\begin{equation}\label{spin:covariance}
 \Phi\bigl(U_\kappa(g)MU_\kappa(g)^*\bigr)
 =U_\ell(g)\Phi(M)U_\ell(g)^*,\qquad
 U_k(g)=\bigotimes_iU_{k_i}(g_i),
\end{equation}
for every $g=(g_1,\ldots,g_n)\in SU(2)^n$.
The input and output sites in this definition are matched in advance.

For a capacity vector $k$, write $b^k_a=\prod_i\binom{k_i}{a_i}$.
Use the Choi convention
$J(\Phi)=\sum_{a,b}\Phi(|a\rangle\langle b|)\otimes
|a\rangle\langle b|$.  Its coherent polynomial is
\begin{equation}\label{spin:choi-weights}
 F_{J(\Phi)}(z,w,u,v)=
 \sum_{c,e,a,b}\sqrt{b^{\ell}_cb^{\ell}_e
                          b^{\kappa}_ab^{\kappa}_b}\,
 [\Phi(|a\rangle\langle b|)]_{ce}\,z^cw^eu^av^b.
\end{equation}
Full coherent Choi stability means stability in all four groups of
variables.  For one matched pair of sites put
\begin{equation}\label{spin:invariants}
 K_o=1+zw,\qquad K_i=1+uv,\qquad
 A=(1+zu)(1+wv),\qquad B=K_oK_i.
\end{equation}

\begin{theorem}[Covariant root classification]\label{spin:covariant}
Suppose $\Phi$ is complex linear, trace preserving, and independently
covariant as in \eqref{spin:covariance}.  Its full coherent Choi kernel is
Lee--Yang if and only if $\Phi=\bigotimes_i\Phi_i$, where each local
factor is specified by an integer $0\leq d_i\leq\min(\kappa_i,\ell_i)$
and an unordered list $s_{i1},\ldots,s_{id_i}\geq0$ as follows.  With
$q_i(r)=\prod_{a=1}^{d_i}(r+s_{ia})$, its local kernel is
\begin{equation}\label{spin:covariant-kernel}
 F_{J(\Phi_i)}=
 \frac{K_{o,i}^{\ell_i-d_i}K_{i,i}^{\kappa_i-d_i}
            \prod_{a=1}^{d_i}(A_i+s_{ia}B_i)}
      {(\ell_i+1)\int_0^1q_i(x)\,dx}.
\end{equation}
Every such map is completely positive, although neither positivity nor
Hermiticity preservation was assumed.  Its complex-linear rank is
\begin{equation}\label{spin:covariant-rank}
 \rank\Phi=\prod_i(d_i+1)^2.
\end{equation}
It is injective exactly when $d_i=\kappa_i$ at every site, surjective
exactly when $d_i=\ell_i$ at every site, and invertible exactly when
$d_i=\kappa_i=\ell_i$ at every site.
\end{theorem}

We prove the theorem together with a factorization through an intermediate
spin representation.  The representation theory of covariant maps and
the general covariant completely positive simplex are established in
\cite{AlNuwairan14}; coherent descriptions and multipliers are also
studied in \cite{ARS24}.  We use the multiplicity-free commutant
description, and derive the additional root restriction directly from
stability.

\begin{lemma}[Range of the coherent ratio]\label{spin:ratio}
The image of $A/B$ on $\D^4$ is precisely
$\Omega=\C\setminus(-\infty,0]$.
\end{lemma}

\begin{proof}
For $a\in\D$, the disk automorphisms
\[
 f_a(x)=\frac{x-a}{1-\bar a x},\qquad
 h_a(y)=\frac{y+\bar a}{1+ay}
\]
satisfy
\[
 1+f_a(x)h_a(y)=
 \frac{(1-|a|^2)(1+xy)}{(1-\bar a x)(1+ay)}.
\]
Apply $f_a$ to $z,v$ and $h_a$ to $u,w$.  The ratio $A/B$ is unchanged.
Choosing $a=z$ makes the new $z$ zero, and reduces the ratio to
$(1+w'v')/(1+u'v')$.  It is nonzero.  If it equalled $-t$ for some
$t>0$, then $1+t=-v'(w'+tu')$, whereas
$|v'|(|w'|+t|u'|)<1+t$.  This proves inclusion in $\Omega$.
Conversely, setting $u=z$ and $v=w=-z$ gives
\[
 \frac AB=\left(\frac{1+z^2}{1-z^2}\right)^2.
\]
The successive maps $z\mapsto z^2$, $x\mapsto(1+x)/(1-x)$, and
$y\mapsto y^2$ map onto the disk, the right half-plane, and $\Omega$,
respectively.  This proves surjectivity.
\end{proof}

\begin{lemma}[Polynomial factorization on a product of slit planes]
\label{spin:slit-product}
If a nonzero polynomial $Q\in\C[r_1,\ldots,r_n]$ has no zero on
$\Omega^n$, then
\begin{equation}\label{spin:slit-factor}
 Q(r)=c\prod_iq_i(r_i),\qquad c\ne0,
\end{equation}
where each $q_i$ is monic and all its roots belong to $(-\infty,0]$.
\end{lemma}

\begin{proof}
Write $Q(x,y)=\sum_{k=0}^da_k(y)x^k$, with $a_d\not\equiv0$.
On any connected open patch of $\Omega^{n-1}$ where $a_d\ne0$, every
root of the monic polynomial $Q(x,y)/a_d(y)$ is real and nonpositive.
Its coefficients $a_k/a_d$ are therefore real-valued holomorphic
functions, and hence constant on the patch.  Polynomial identity gives
$Q(x,y)=a_d(y)p(x)$ everywhere, for a fixed monic polynomial $p$ with
all roots in $(-\infty,0]$.  Since $p$ is nonzero on $\Omega$, the
polynomial $a_d$ has no zero on $\Omega^{n-1}$.  Induction proves the
claim, skipping variables of degree zero.  This reasoning includes
repeated roots and slices on which the degree might initially drop.
\end{proof}

For $0\leq r\leq k$, let $R_{k\to r}$ be partial trace from a
Dicke-encoded $k$-qubit input to $r$ qubits, identified with
$\operatorname{End}(\mathcal H_r)$.  The reduced operator is supported
on the symmetric sector: a vector antisymmetric in a retained pair is
orthogonal to the fully symmetric input sector.  Define the adjoint map
\begin{equation}\label{spin:raw-cloner}
 C_{r\to k}(N)=
 W_k^*\bigl(W_rNW_r^*\otimes I_2^{\otimes(k-r)}\bigr)W_k.
\end{equation}
These maps are completely positive and covariant; $R_{k\to r}$ is
trace preserving.  At $r=0$ they are the ordinary trace and scalar
multiplication of $I_{k+1}$, respectively.

\begin{lemma}[Coherent compression kernels]\label{spin:compression}
For $0\leq r\leq\min(\kappa,\ell)$, the map
$T_r^{\kappa,\ell}=C_{r\to\ell}R_{\kappa\to r}$ has coherent Choi
kernel
\begin{equation}\label{spin:basis-kernel}
 A^rK_o^{\ell-r}K_i^{\kappa-r}.
\end{equation}
These $\min(\kappa,\ell)+1$ maps form a basis of all complex-linear
$SU(2)$-covariant maps between the two matrix spaces.  Their trace
effects are
\begin{equation}\label{spin:basis-effect}
 (T_r^{\kappa,\ell})^*(I_{\ell+1})
            =\frac{\ell+1}{r+1}I_{\kappa+1}.
\end{equation}
\end{lemma}

\begin{proof}
The Dicke coefficients give, for $c=a+k$ and $e=b+k$,
\begin{equation}\label{spin:compression-coefficients}
 [C_{r\to\ell}(|a\rangle\langle b|)]_{c,e}
 =\frac{\sqrt{\binom ra\binom rb}\binom{\ell-r}{k}}
             {\sqrt{\binom\ell c\binom\ell e}},
\end{equation}
and zero if the two shifts do not agree.  In particular,
\begin{equation}\label{spin:kernel-multiplication}
 F_{C_{r\to\ell}(N)}(z,w)=(1+zw)^{\ell-r}F_N(z,w).
\end{equation}
Adjunction gives the coefficients of $R_{\kappa\to r}$, with
$\ell$ replaced by $\kappa$.  If its input indices are $a+m,b+m$,
the two square-root binomial weights from the output and input of the
Choi kernel cancel exactly the two denominators in these formulas.
What remains is
\[
 \binom ra\binom rb\binom{\ell-r}{k}\binom{\kappa-r}{m},
\]
the coefficient in the expansion of \eqref{spin:basis-kernel}.

Covariance is equivalent to the Choi matrix commuting with
$U_\ell\otimes\overline{U_\kappa}$.  The Clebsch--Gordan decomposition
is multiplicity free, with spins
$|\ell-\kappa|/2,|\ell-\kappa|/2+1,\ldots,(\ell+\kappa)/2$.
Its commutant therefore has dimension $\min(\kappa,\ell)+1$.
The displayed covariant kernels are independent, since division by
$K_o^\ell K_i^\kappa$ gives the successive powers of $A/B$, whose
range is open by Lemma~\ref{spin:ratio}.  They consequently span the
entire complex commutant, without a Hermiticity assumption.

Finally $C_{r\to\ell}(I_{r+1})=I_{\ell+1}$: the projection onto the
symmetric sector of the first $r$ qubits acts identically on the fully
symmetric $\ell$-qubit sector.  Covariance and irreducibility imply that
$C_{r\to\ell}^*(I_{\ell+1})$ is a scalar multiple of $I_{r+1}$.
Taking traces identifies this scalar as $(\ell+1)/(r+1)$.
Composing with the trace-preserving $R_{\kappa\to r}$ proves
\eqref{spin:basis-effect}.
\end{proof}

\begin{proof}[Proof of Theorem~\ref{spin:covariant}]
For the independent product group, the commutant is the tensor product
of the local commutants in Lemma~\ref{spin:compression}.  There is
therefore a unique polynomial $Q$ of separate degrees at most
$\min(\kappa_i,\ell_i)$ such that
\begin{equation}\label{spin:global-invariant}
 F_{J(\Phi)}=
 \prod_iK_{o,i}^{\ell_i}K_{i,i}^{\kappa_i}\,
 Q\left(\frac{A_1}{B_1},\ldots,\frac{A_n}{B_n}\right).
\end{equation}
The separate degree bounds make this a polynomial.  The prefactor
never vanishes in the polydisk.  The variables of different sites are
independent, so Lemma~\ref{spin:ratio} shows that full coherent stability
is equivalent to nonvanishing of $Q$ on $\Omega^n$.
Lemma~\ref{spin:slit-product} gives
$Q=c\prod_iq_i$, where
$q_i(r)=\prod_{a=1}^{d_i}(r+s_{ia})$ and $s_{ia}\geq0$.
The kernel thus factors, up to one nonzero scalar, into the numerators
of \eqref{spin:covariant-kernel}.

For one factor write $q(r)=\sum_{a=0}^dq_ar^a$ and set
$\Psi_{\kappa,\ell,q}=\sum_{a=0}^dq_aT_a^{\kappa,\ell}$.
Every coefficient $q_a$ is nonnegative, so this map is completely
positive.  By \eqref{spin:basis-effect} its trace effect is
\begin{equation}\label{spin:normalization}
 \Psi_{\kappa,\ell,q}^*(I_{\ell+1})
 =c_\ell(q)I_{\kappa+1},\qquad
 c_\ell(q)=(\ell+1)\sum_{a=0}^d\frac{q_a}{a+1}
          =(\ell+1)\int_0^1q(x)\,dx>0.
\end{equation}
In the global factorization, trace preservation forces
$c\prod_ic_{\ell_i}(q_i)=1$.  Thus the initial scalar is positive
real and the map is the tensor product of the individually normalized
completely positive maps in \eqref{spin:covariant-kernel}.
Conversely these maps are covariant and trace preserving by construction,
and their kernels are stable because $q_i$ has no root in $\Omega$.
This proves the classification and automatic complete positivity.

To prove the rank assertion, first consider the square case $d=\kappa
=\ell$.  The map $\Psi_{d,d,q}$ is also the Dicke compression
\begin{equation}\label{spin:core-compression}
 \Psi_{d,d,q}(M)=W_d^*\left[
     \left(\bigotimes_{a=1}^d\mathcal T_{s_a}\right)(W_dMW_d^*)\right]W_d,
 \qquad \mathcal T_s(N)=N+s\Tr(N)I_2,
\end{equation}
where the tensor product in brackets acts on the $d$ qubit matrix
factors.  Its kernel is $\prod_a(A+s_aB)$ by diagonalization of the
binary variables, which verifies the identity with $\Psi_{d,d,q}$.
As an operator on Hilbert--Schmidt space, $\mathcal T_s$ is positive
definite: its eigenvalues are $1$ on traceless matrices and $1+2s$ on
the scalar space.  The embedding $M\mapsto W_dMW_d^*$ is a
Hilbert--Schmidt isometry.  Its compression of the positive definite
tensor product is therefore positive definite and invertible.  Define
\begin{equation}\label{spin:core}
 \Lambda_{d,q}=\frac{\Psi_{d,d,q}}{c_d(q)},\qquad
 \Lambda_{0,1}=\Id_{\C}.
\end{equation}

The multiplication identity \eqref{spin:kernel-multiplication} gives
$C_{d\to\ell}C_{a\to d}=C_{a\to\ell}$, and adjunction gives
$R_{d\to a}R_{\kappa\to d}=R_{\kappa\to a}$.  Hence
\begin{equation}\label{spin:raw-factorization}
 \Psi_{\kappa,\ell,q}
 =C_{d\to\ell}\Psi_{d,d,q}R_{\kappa\to d}.
\end{equation}
The map $C_{d\to\ell}$ is injective by
\eqref{spin:kernel-multiplication}; its adjoint $R_{\ell\to d}$ is
therefore surjective, and the same holds for $R_{\kappa\to d}$.
The middle map in \eqref{spin:raw-factorization} is invertible on a
space of dimension $(d+1)^2$.  This proves local rank $(d+1)^2$.
Ranks multiply under tensor products, which gives
\eqref{spin:covariant-rank} and all assertions about injectivity,
surjectivity, and invertibility.
\end{proof}

\begin{corollary}[Trace, invertible core, and symmetric cloning]
\label{spin:trace-core-cloner}
Each local map in Theorem~\ref{spin:covariant} has the exact factorization
\begin{equation}\label{spin:normalized-factorization}
 \Phi_i=\mathcal C_{d_i\to\ell_i}\circ\Lambda_{d_i,q_i}
                                      \circ R_{\kappa_i\to d_i},
 \qquad
 \mathcal C_{d\to\ell}=\frac{d+1}{\ell+1}C_{d\to\ell}.
\end{equation}
All three maps are trace preserving and completely positive, and the
middle map is an invertible coherent Lee--Yang map.
For equal input and output capacities the local factors are also unital.
\end{corollary}

\begin{proof}
The first two assertions follow from \eqref{spin:raw-factorization},
\eqref{spin:normalization}, and
$c_d(q)/c_\ell(q)=(d+1)/(\ell+1)$.  For equal capacities,
covariance makes $\Phi(I)$ scalar, and trace preservation makes this
scalar one.
\end{proof}

The normalized map $\mathcal C_{d\to\ell}$ is the universal symmetric
cloner of \cite[Eq.~(3.3)]{Werner98}; its dimension-ratio normalization
and concatenation law are classical.  Here stability forces the entire
factorization \eqref{spin:normalized-factorization}.  The integer $d$
specifies an intermediate spin dimension and the exact algebraic rank;
it is not an assertion of noiseless recovery or an operational quantum
communication capacity.  If a specified covariance relation permutes the
matched output sites, the theorem applies after undoing that permutation;
an invertible map then matches equal capacities.

\begin{example}[The spin-one spectral region]\label{spin:spin-one}
For $\kappa=\ell=2$, let $\lambda_1,\lambda_2$ be the multipliers of a
trace-preserving covariant map on the dipole and quadrupole subspaces.
The exact coherent Choi condition is
\begin{equation}\label{spin:spin-one-region}
 0\leq\lambda_2\leq\lambda_1,\qquad
 3\lambda_1\leq2+\lambda_2,\qquad
 3\lambda_1^2+\lambda_2^2\geq4\lambda_2.
\end{equation}
Indeed, in the degree-two case write $q(r)=r^2+ar+b$.
Its roots are real and nonpositive exactly when
$a,b\geq0$ and $a^2\geq4b$.  The maps with kernels $A^2$, $AB$,
and $B^2$ have scalar multipliers $1,3/2,3$, respectively, dipole
multipliers $1,1,0$, and quadrupole multipliers $1,0,0$.
For the middle map, partial trace annihilates the quadrupole
$\diag(1,-2,1)$, while the difference of the two extreme projectors is
sent to $S^z$; covariance determines its action on each irreducible
operator space.  Therefore
\[
 \lambda_1=\frac{1+a}{1+3a/2+3b},\qquad
 \lambda_2=\frac1{1+3a/2+3b}.
\]
Solving for $a,b$ gives \eqref{spin:spin-one-region} when
$\lambda_2>0$.  Degree one gives
$\lambda_2=0$ and $0<\lambda_1\leq2/3$, and degree zero gives the
origin, proving all endpoint cases.  In particular, the ordinary
qutrit depolarizing map with $\lambda_1=\lambda_2=t$ fails coherent
Choi stability for every $0<t<1$, although both endpoints are stable.
\end{example}

\begin{remark}\label{spin:scope}
Theorem~\ref{spin:covariant} assumes independent local covariance, not
only covariance under a common rotation.  Its proof uses the static coherent Choi kernel.  The
argument does not apply the
invertible qubit classification to a polarized higher-spin map: a
Dicke lift is generally singular on the full qubit matrix space.
No unrestricted noncovariant higher-spin static classification is
asserted here.
\end{remark}


\section{Quantitative rigidity and correlated tests}
\label{rob:section}

The finite-test rigidity of Theorem~\ref{dyn:trine} is an exact statement.
Its quantitative behaviour depends on the allowed inputs.  We prove that
products of two GHZ states detect precisely the failure of a diagonal
Hamiltonian to have additive energies.  The approximate-additivity theorem
of Kalton and Roberts then gives a dimension-independent comparison with
distance from one-site dynamics.  In contrast, even all product Lee--Yang
inputs admit no such comparison.

Throughout this section, a density matrix acts on
$\mathcal H_n=(\C^2)^{\otimes n}$, and
\begin{align*}
 F_X(z,w)&=\sum_{x,y\in\{0,1\}^n}X_{xy}z^xw^y,\\
 \mathcal Y_n&=\{\rho\geq0:\Tr\rho=1,\ 
                 F_\rho\ne0\text{ on }\D^{2n}\}.
\end{align*}
The condition is zero-freeness, not merely that $F_\rho$ is a nonzero
polynomial.  We use the unhalved trace norm $\|\cdot\|_1$ and diamond norm
$\|\cdot\|_\diamond$, and write
$\operatorname{dist}_1(X,\mathcal Y_n)=
\inf_{\rho\in\mathcal Y_n}\|X-\rho\|_1$.
Let $\operatorname{LocHP}_n$ be the real vector space of maps
\[
 \mathcal M=\sum_{i=1}^n
       \mathcal M_i\otimes\Id_{\{1,\ldots,n\}\setminus\{i\}},
\]
where every $\mathcal M_i$ is complex linear and Hermiticity preserving.
Let $\operatorname{LocGKLS}_n$ be the cone obtained by requiring each
$\mathcal M_i$ to be a one-qubit GKLS generator.  In particular, local
Hamiltonian generators belong to both classes.

\subsection{The diagonal quotient norm}

Identify a subset $S\subseteq[n]$ with the computational basis vector
whose occupied sites are exactly $S$.  For a real function $h$ on $2^{[n]}$,
put
\begin{gather*}
 H=\diag(h(S):S\subseteq[n]),\qquad
 \mathcal L_H(X)=-i[H,X],\\
 \operatorname{osc}h=\max_S h(S)-\min_S h(S).
\end{gather*}

\begin{proposition}
\label{rob:distance}
For every diagonal Hermitian $H$,
\begin{align}
 d(H)
 &:=
 \operatorname{dist}_\diamond(\mathcal L_H,\operatorname{LocHP}_n)
 \notag\\
 &=\operatorname{dist}_\diamond(\mathcal L_H,\operatorname{LocGKLS}_n)
 \notag\\
 &=\min_{b\in\R^n}\operatorname{osc}
       \left(h(S)-\sum_{i\in S}b_i\right).
 \label{rob:distance-formula}
\end{align}
If $h(S)=h([n]\setminus S)$ for all $S$, then
$d(H)=\operatorname{osc}h=\|\mathcal L_H\|_\diamond$.
\end{proposition}

\begin{proof}
Fix $\mathcal M\in\operatorname{LocHP}_n$ and set
$E_{xy}=|x\rangle\langle y|$.  The coefficient of $E_{xy}$ in
$\mathcal M(E_{xy})$ equals
\[
 \langle x|\mathcal M(E_{xy})|y\rangle
 =\sum_i\langle x_i|\mathcal M_i(E_{x_i y_i})|y_i\rangle.
\]
If $x_i=y_i$, the summand is real.  The summands for $(x_i,y_i)=(0,1)$
and $(1,0)$ are conjugates.  Thus, for a vector $b\in\R^n$ independent
of $x,y$, its imaginary part is $b\cdot(y-x)$.  Since
$\|E_{xy}\|_1=1$ and every matrix entry is bounded by the trace norm,
\[
 \|\mathcal L_H-\mathcal M\|_\diamond
 \geq \max_{x,y}
       |(h(x)-b\cdot x)-(h(y)-b\cdot y)|
 =\operatorname{osc}(h-b\cdot x).
\]
Here and below $h(x)$ is the subset notation written in occupation
coordinates.

For a Hermitian matrix $D$ one has
\begin{equation}
 \|-i[D,\cdot]\|_\diamond
   =\lambda_{\max}(D)-\lambda_{\min}(D).
 \label{rob:commutator-norm}
\end{equation}
Indeed, subtract the midpoint of the spectral interval from $D$.
The trace-norm ideal inequality, also after adjoining an identity
ancilla, gives the upper bound $2\|D\|_{\mathrm{op}}$ for the resulting
commutator.  The matrix unit between eigenvectors at the two spectral
endpoints gives the reverse bound.  Apply this identity to
$D=H-H_b$, where
$H_b=\sum_i b_i|1\rangle\langle1|_i$.
The local Hamiltonian generator $\mathcal L_{H_b}$ attains the
corresponding oscillation.  Because
$\operatorname{LocGKLS}_n\subseteq\operatorname{LocHP}_n$,
the two distance infima agree with the infimum of these oscillations.
It is a minimum: the objective is continuous and
\[
 \operatorname{osc}(h-b\cdot x)
 \geq \operatorname{osc}(b\cdot x)-\operatorname{osc}h
 =\|b\|_{\ell^1}-\operatorname{osc}h.
\]

Suppose now that $h$ is invariant under complementation.  Replacing
$x$ by $1-x$ shows that
$\operatorname{osc}(h-b\cdot x)=\operatorname{osc}(h+b\cdot x)$.
The oscillation is a convex function of its argument, so
\[
 \operatorname{osc}h
 \leq \tfrac12\operatorname{osc}(h-b\cdot x)
       +\tfrac12\operatorname{osc}(h+b\cdot x)
 =\operatorname{osc}(h-b\cdot x).
\]
Taking $b=0$ proves the last assertion.
\end{proof}

\subsection{Two correlated blocks}

For disjoint nonempty $A,B\subseteq[n]$, define
\begin{equation}
\begin{aligned}
 |\psi_{A,B}\rangle
 &=\frac{|0_A\rangle+|1_A\rangle}{\sqrt2}
  \otimes\frac{|0_B\rangle+|1_B\rangle}{\sqrt2}
  \otimes|0_{[n]\setminus(A\cup B)}\rangle,\\
 \rho_{A,B}&=|\psi_{A,B}\rangle\langle\psi_{A,B}|.
\end{aligned}
 \label{rob:ghz-probe}
\end{equation}
The pure-state coefficient polynomial is
$\frac12(1+\prod_{i\in A}z_i)(1+\prod_{i\in B}z_i)$.
It is zero-free on $\D^n$, and hence $\rho_{A,B}\in\mathcal Y_n$.
Set
\begin{align*}
 \Delta_{A,B}(H)&=h(A\cup B)-h(A)-h(B)+h(\varnothing),\\
 \Delta(H)&=\max_{A\cap B=\varnothing}|\Delta_{A,B}(H)|.
\end{align*}

\begin{theorem}[Correlated diagonal rigidity]
\label{rob:diagonal}
Let $n\geq2$ and let $H$ be a diagonal Hermitian Hamiltonian.  For every
pair of disjoint nonempty subsets $A,B$ the following limit exists:
\begin{equation}
 \lim_{t\downarrow0}
 \frac{\operatorname{dist}_1
       (e^{-itH}\rho_{A,B}e^{itH},\mathcal Y_n)}{t}
 =\frac{|\Delta_{A,B}(H)|}{2}.
 \label{rob:exact-defect}
\end{equation}
If $\delta(H)$ denotes the maximum of these limits, then
\begin{equation}
 \delta(H)=\frac{\Delta(H)}2,
 \qquad
 \delta(H)\leq d(H)\leq4K\,\delta(H),
 \label{rob:diagonal-bounds}
\end{equation}
where $K$ is any universal constant in the real Kalton--Roberts
approximate-additivity theorem.  In particular, one may take $K=45$.
\end{theorem}

We first establish the polynomial and boundary facts used in the proof.

\begin{lemma}[Block compression]
\label{rob:compression}
Let $f(u_1,\ldots,u_k,v)$ be multiaffine and zero-free on the open unit
polydisk.  If $a(v)$ and $d(v)$ are its constant and full-degree
coefficients in the $u$ variables, then $a(v)+d(v)s$ is zero-free on
the open unit polydisk in $(s,v)$.

Consequently, if
$V:\C^2\otimes\C^2\longrightarrow\mathcal H_n$ is the isometry
\[
 V|a,b\rangle=|a_A,b_B,0_{[n]\setminus(A\cup B)}\rangle,
\]
then $F_{V^*\sigma V}$ is zero-free on $\D^4$ for every
$\sigma\in\mathcal Y_n$.
\end{lemma}

\begin{proof}
Fix the remaining variables $v$ in their polydisk.  The polynomial
$g(u)=f(u,\ldots,u,v)$ has no zero in $\D$, and $a(v)=g(0)\ne0$.
If $\deg g=k$, all its roots have modulus at least one, and the
product-of-roots identity gives $|d(v)|\leq|a(v)|$.
If $\deg g<k$, then $d(v)=0$.  Thus $a(v)+d(v)s\ne0$ for $|s|<1$.
This proves the first assertion; it is the unit-disk form of the
Asano coefficient contraction \cite{AsanoJapanese70}.

For the second assertion, first set all row and column variables
outside $A\cup B$ to zero.  Apply the first assertion successively to
the row variables in $A$, the row variables in $B$, and the corresponding
two groups of column variables.  The resulting coefficient polynomial
is exactly $F_{V^*\sigma V}$.  No normalization or zero-polynomial
exception occurs in this argument.
\end{proof}

The compression in Lemma~\ref{rob:compression} is completely positive
and trace nonincreasing, but need not preserve trace.  We use its
unnormalized output.  For every Hermitian $X$,
$\|V^*XV\|_1\leq\|X\|_1$, as follows either from trace-norm duality
or from the positive and negative parts of $X$.

\begin{lemma}[A two-qubit boundary contact]
\label{rob:contact}
Write $\rho_0=|++\rangle\langle++|$, where
$|+\rangle=(|0\rangle+|1\rangle)/\sqrt2$, and put
\begin{align*}
 u&=|00\rangle+|10\rangle+|01\rangle+|11\rangle,\\
 v&=|00\rangle-|10\rangle-|01\rangle+|11\rangle,\qquad
 \ell(X)=\Im\langle v|X|u\rangle.
\end{align*}
If $t_j\downarrow0$, $X_j$ is Hermitian, $F_{X_j}$ is zero-free on
$\D^4$, and $\|X_j-\rho_0\|_1=O(t_j)$, then
$\ell(X_j)=o(t_j)$.  Moreover,
\begin{equation}
 |\ell(X)|\leq2\|X\|_1
 \quad\text{for every Hermitian }X.
 \label{rob:functional-norm}
\end{equation}
\end{lemma}

\begin{proof}
The Hermitian observable representing $\ell$ is
\[
 R=\frac{|u\rangle\langle v|-|v\rangle\langle u|}{2i}.
\]
The vectors $u,v$ are orthogonal and both have norm two.  On their span
$R$ has eigenvalues $2,-2$, and it vanishes on the orthogonal
complement.  Thus $\|R\|_{\mathrm{op}}=2$, proving
\eqref{rob:functional-norm}.

Consider
\[
 g_j(s)=F_{X_j}(s-1,s-1,1,1)
       =a_{2,j}s^2+a_{1,j}s+a_{0,j}.
\]
For $X_j=\rho_0$ this polynomial is $s^2$.  Continuity of its
coefficients gives $a_{2,j}=1+O(t_j)$ and
$a_{1,j},a_{0,j}=O(t_j)$.  The column variables have been specialized
to boundary points, so we justify their use.  For fixed $j$, the
polynomials $F_{X_j}(s-1,s-1,r,r)$, with $0<r<1$, are zero-free on
$|s-1|<1$.  Hurwitz's theorem shows that their limit as $r\uparrow1$
is zero-free there or identically zero.  For all sufficiently large
$j$, $a_{2,j}\ne0$, which excludes the latter alternative.

Let $r_{1,j},r_{2,j}$ be the two roots of $g_j$, counted with
multiplicity.  Vieta's formulas give their sum and product as $O(t_j)$;
the quadratic formula then gives $|r_{k,j}|=O(\sqrt{t_j})$.
Because $|r_{k,j}-1|\geq1$,
\[
 \Re r_{k,j}\leq\tfrac12|r_{k,j}|^2=O(t_j).
\]
Their real parts have sum $O(t_j)$, so the same real parts are bounded
below by $-O(t_j)$.  Consequently
$\Re r_{k,j}=O(t_j)$ and $\Im r_{k,j}=O(\sqrt{t_j})$.
It follows that
\[
 \Im(r_{1,j}r_{2,j})=O(t_j^{3/2}),\qquad
 \Im a_{0,j}=O(t_j^{3/2})+O(t_j^2)=o(t_j).
\]
Finally $a_{0,j}=F_{X_j}(-1,-1,1,1)=\langle v|X_j|u\rangle$.
\end{proof}

\begin{proof}[Proof of Theorem~\ref{rob:diagonal}]
Fix $A,B$, let $V$ be as in Lemma~\ref{rob:compression}, and abbreviate
$\Delta=\Delta_{A,B}(H)$ and
$\rho(t)=e^{-itH}\rho_{A,B}e^{itH}$.
The image of $V$ is invariant under $H$, and
$V^*\rho(t)V=e^{-itD}\rho_0e^{itD}$, where the logical diagonal
Hamiltonian $D$ has energies
$h(\varnothing),h(A),h(B),h(A\cup B)$.
Since $\rho_0=|u\rangle\langle u|/4$ and $\langle v|u\rangle=0$,
differentiation gives
\begin{equation}
 \ell(V^*\rho(t)V)=-\Delta t+O(t^2).
 \label{rob:contact-derivative}
\end{equation}

Let $t_j\downarrow0$ be arbitrary, and choose
$\sigma_j\in\mathcal Y_n$ with
\[
 \|\rho(t_j)-\sigma_j\|_1
 \leq\operatorname{dist}_1(\rho(t_j),\mathcal Y_n)+t_j^2.
\]
Because $\rho_{A,B}\in\mathcal Y_n$ and
$\rho(t_j)=\rho_{A,B}+O(t_j)$, both the distance on the right and
$\|\sigma_j-\rho_{A,B}\|_1$ are $O(t_j)$.
Lemma~\ref{rob:compression} shows that
$X_j=V^*\sigma_j V$ has a zero-free coefficient polynomial.
It is Hermitian and $X_j=\rho_0+O(t_j)$, so
Lemma~\ref{rob:contact} applies without rescaling its trace.
Equations~\eqref{rob:functional-norm} and
\eqref{rob:contact-derivative} imply
\[
 2\|\rho(t_j)-\sigma_j\|_1
 \geq |\ell(V^*\rho(t_j)V-X_j)|
 =|\Delta|t_j+o(t_j).
\]
As the sequence was arbitrary, this proves the lower bound
$|\Delta|/2$ for the liminf in \eqref{rob:exact-defect}.

For a matching upper bound, subtract $(\Delta/4)Z_1Z_2$ from $D$,
where $Z=\diag(1,-1)$.  The remaining logical energy is affine in the
two occupation bits: its mixed difference is zero.  Its evolution of
the logical vector $|++\rangle$, embedded by $V$, is therefore a
product of two GHZ states with unit-modulus relative phases and a
vacuum on the remaining sites.  This comparison state belongs to
$\mathcal Y_n$.  Since all four logical configurations have equal
probability and the two values of $Z_1Z_2$ occur equally often, the
inner product between the actual and comparison pure vectors is
$\cos(t\Delta/4)$.  Their trace distance is
\[
 2|\sin(t\Delta/4)|\leq t|\Delta|/2\qquad(t\geq0).
\]
This proves \eqref{rob:exact-defect}.  Empty subsets have mixed
difference zero, and the family of pairs is finite, so
$\delta(H)=\Delta(H)/2$.

We now invoke the precise inherited input.  The Kalton--Roberts
theorem \cite[Theorem~4.1]{KaltonRoberts83} supplies a universal
$K\leq45$ with the following property: if a real set function $f$ on
a set algebra satisfies $f(\varnothing)=0$ and
\[
 |f(A\cup B)-f(A)-f(B)|\leq\eta
 \quad(A\cap B=\varnothing),
\]
then there is an additive real set function $\mu$ such that
$\sup_S|f(S)-\mu(S)|\leq K\eta$.
Apply it on $2^{[n]}$ to $f(S)=h(S)-h(\varnothing)$ and
$\eta=\Delta(H)$.  On this finite algebra,
$\mu(S)=\sum_{i\in S}b_i$ for some real $b_i$, and hence
\[
 \operatorname{osc}(h-b\cdot x)\leq2K\Delta(H).
\]
If $\Delta(H)=0$, the same conclusion follows directly from finite
additivity.  Proposition~\ref{rob:distance} now yields
$d(H)\leq2K\Delta(H)=4K\delta(H)$.

Finally choose a minimizer $b$ in \eqref{rob:distance-formula}.
The flow of $H_b$ consists of independent diagonal one-qubit
unitaries.  It rotates each polynomial variable by a unit complex
number and thus preserves $\mathcal Y_n$.
The Duhamel formula and unitary trace-norm invariance imply, for every
density matrix $\rho$ and $t\geq0$,
\[
 \|e^{t\mathcal L_H}\rho-e^{t\mathcal L_{H_b}}\rho\|_1
 \leq t\|\mathcal L_H-\mathcal L_{H_b}\|_\diamond=t\,d(H).
\]
Apply this estimate to each $\rho_{A,B}$, divide by $t$, and use
\eqref{rob:exact-defect}.  Taking the maximum gives
$\delta(H)\leq d(H)$.
\end{proof}

\begin{remark}
\label{rob:scope}
The universal constant in \eqref{rob:diagonal-bounds} is inherited
from Kalton--Roberts.  The bridge to that theorem is the exact
root-contact identity \eqref{rob:exact-defect}.
The probes may contain arbitrarily large blocks, their family is
exponential in $n$, and the defect uses distance to the full set
$\mathcal Y_n$.  No efficient certification procedure or corresponding
statement for non-diagonal Hamiltonians is asserted.
\end{remark}

\subsection{Product inputs do not give uniform robustness}

Let $\mathcal P_n\subseteq\mathcal Y_n$ be the set of all product
density matrices in $\mathcal Y_n$, including mixed products, and put
\[
 \delta_{\mathrm{prod}}(\mathcal L)
 =\sup_{\rho\in\mathcal P_n}\limsup_{t\downarrow0}
   \frac{\operatorname{dist}_1(e^{t\mathcal L}\rho,\mathcal Y_n)}{t}.
\]
We give two different mechanisms by which this defect can be small.
For use with mixed states, recall the elementary estimate
\begin{equation}
 \|e^{-itH}\rho e^{itH}-\rho\|_1
 \leq2|t|\sqrt{\Tr(\rho H^2)-(\Tr(\rho H))^2}.
 \label{rob:variance-bound}
\end{equation}
To prove it, purify $\rho$ to a unit vector $\Psi$ and evolve the
purification by $H\otimes I$.  For a pure-state projection
$P=|\Psi\rangle\langle\Psi|$, the two possibly nonzero eigenvalues
of $-i[H\otimes I,P]$ are
$\pm\sqrt{\operatorname{Var}_\rho H}$.
Its trace norm is therefore twice this square root.
The variance is constant along the unitary orbit, so integration gives
the same bound for purified states.  Partial trace gives
\eqref{rob:variance-bound}.

\begin{proposition}[A nearest-neighbour obstruction]
\label{rob:path}
On a path with $m\geq1$ edges and $n=m+1$ vertices, let
\[
 H_m=\frac1{2m}\sum_{j=1}^m Z_jZ_{j+1},
 \qquad \mathcal L_m=-i[H_m,\cdot].
\]
Then
\begin{gather*}
 \|\mathcal L_m\|_\diamond=d(H_m)=1,\\
 \operatorname{dist}_1(e^{t\mathcal L_m}\rho,\mathcal Y_n)
 \leq\sqrt{\frac3m}|t|
 \quad(\rho\in\mathcal P_n,\ t\in\R).
\end{gather*}
For the equatorial pure product
$|\psi_\theta\rangle=\bigotimes_{j=1}^n
(|0\rangle+e^{i\theta_j}|1\rangle)/\sqrt2$ and its density
$\rho_\theta$, the sharper formula is
\begin{equation}
 \|e^{t\mathcal L_m}\rho_\theta-\rho_\theta\|_1
 =2\sqrt{1-\cos(t/(2m))^{2m}}
 \leq\frac{|t|}{\sqrt m}.
 \label{rob:path-exact}
\end{equation}
\end{proposition}

\begin{proof}
The diagonal energy is invariant under simultaneous complementation.
Equal computational spins have energy $1/2$, alternating spins have
energy $-1/2$, and these are the extrema.  Thus
Proposition~\ref{rob:distance} gives the two norm identities.

For a product density, the computational spins are independent signs.
Write $W_j=Z_jZ_{j+1}$ for their edge products, regarded as random
variables in this diagonal distribution.  Disjoint edges have zero
covariance, each variance is at most one, and each adjacent covariance
has absolute value at most one by Cauchy--Schwarz.  Hence
\[
 \operatorname{Var}\left(\sum_{j=1}^m W_j\right)
 \leq m+2(m-1)\leq3m.
\]
Equation~\eqref{rob:variance-bound} gives the stated mixed-state bound.
The initial state itself belongs to $\mathcal Y_n$ and is an
admissible comparison state in the distance.

For $\psi_\theta$, the computational distribution is uniform.
The root sign and the $m$ edge products determine all vertex signs
bijectively, so the edge products are independent uniform signs.
It follows that
\[
 \langle\psi_\theta|e^{-itH_m}|\psi_\theta\rangle
 =\cos(t/(2m))^m.
\]
The trace distance of unit pure states $\psi,\phi$ is
$2\sqrt{1-|\langle\psi,\phi\rangle|^2}$; this follows by restricting
the difference of their projections to their two-dimensional span.
It proves the equality in \eqref{rob:path-exact}.
Independence gives $\operatorname{Var}_{\rho_\theta}H_m=1/(4m)$,
so \eqref{rob:variance-bound} proves its last inequality.
Finally, $\rho_\theta\in\mathcal Y_n$, since its coefficient polynomial
is a product of linear factors with roots on the unit circle.
\end{proof}

\begin{lemma}[One-qubit physical states]
\label{rob:onequbit}
A one-qubit density matrix
$\tau=\left(\begin{smallmatrix}a&\overline b\\ b&d\end{smallmatrix}\right)$
belongs to $\mathcal Y_1$ if and only if $a\geq d$.
\end{lemma}

\begin{proof}
Its polynomial is
$f(z,w)=a+bz+\overline b w+dzw$.
Suppose $a\geq d$.  Positivity gives $|b|^2\leq ad\leq a^2$,
and $a>0$ because $a+d=1$.  The denominator of
\[
 q(z)=\frac{\overline b+dz}{a+bz}
\]
has no zero in $\D$.  For $|z|=1$,
\[
 |a+bz|^2-|\overline b+dz|^2
 =(a-d)(a+d+2\Re(bz))\geq0.
\]
If $|b|<a$, the maximum principle gives $|q(z)|\leq1$ in $\D$.
The remaining case $|b|=a$ forces $a=d=|b|$, and cancellation makes
$q$ a constant of modulus one.  Thus in either case
$f(z,w)=(a+bz)(1+wq(z))$ is nonzero for $|z|,|w|<1$.

Conversely, suppose $a<d$.  Then
$|b|\leq\sqrt{ad}<(a+d)/2$, so the displayed boundary difference is
strictly negative on the unit circle.  By continuity, some $z$ inside
the circle satisfies
$|a+bz|<|\overline b+dz|$.
The denominator in
$w=-(a+bz)/(\overline b+dz)$ is nonzero and $|w|<1$;
this gives an interior zero of $f$.
\end{proof}

\begin{proposition}[An exponential product-input obstruction]
\label{rob:rankone}
For $n\geq2$, put $P_n=|1^n\rangle\langle1^n|$ and
$\mathcal K_n=-i[P_n,\cdot]$.  Then
\begin{equation}
 \|\mathcal K_n\|_\diamond=1,\qquad
 d(P_n)=\frac{n-1}{n}.
 \label{rob:rankone-distance}
\end{equation}
For every $\rho\in\mathcal P_n$ and every real $t$,
\begin{equation}
 \operatorname{dist}_1(e^{t\mathcal K_n}\rho,\mathcal Y_n)
 \leq\|e^{t\mathcal K_n}\rho-\rho\|_1
 \leq2^{1-n/2}|t|.
 \label{rob:rankone-bound}
\end{equation}
For equatorial pure products, the middle expression is exactly
\[
 4\sqrt{2^{-n}(1-2^{-n})}\,|\sin(t/2)|.
\]
\end{proposition}

\begin{proof}
The commutator norm is one by \eqref{rob:commutator-norm}.
To evaluate \eqref{rob:distance-formula}, average $b$ over coordinate
permutations.  Convexity of oscillation and permutation symmetry of
$h(x)=\mathbf1_{\{x=1^n\}}$ show that this cannot increase the objective.
We may therefore take $b_i=c$ for all $i$.  The residual energies are
\[
 \{-ck:0\leq k\leq n-1\}\ \cup\ \{1-cn\}.
\]
If $c\leq1/n$, the last value minus the value at weight $n-1$ is
$1-c\geq(n-1)/n$.  If $c\geq1/n$, the difference between weights
zero and $n-1$ has absolute value $c(n-1)\geq(n-1)/n$.
For $c=1/n$, all residual energies lie in
$[-(n-1)/n,0]$ and both endpoints occur.  This proves
\eqref{rob:rankone-distance}.

The coefficient polynomial of a product density is the product of
the individual coefficient polynomials in disjoint variables.  Thus
it is zero-free precisely when all factors are zero-free.
Lemma~\ref{rob:onequbit} implies that every occupied population is at
most $1/2$.  Consequently $p=\Tr(P_n\rho)\leq2^{-n}$.
For a purification $\Psi$ of $\rho$,
\[
 \langle\Psi|(e^{-itP_n}\otimes I)|\Psi\rangle
 =1-p+pe^{-it}.
\]
The pure-state trace-distance formula and partial-trace contractivity
therefore give
\[
 \|e^{t\mathcal K_n}\rho-\rho\|_1
 \leq4\sqrt{p(1-p)}\,|\sin(t/2)|
 \leq2\sqrt p\,|t|.
\]
This proves \eqref{rob:rankone-bound}, using the initial LY state as
comparison.  For an equatorial pure product $p=2^{-n}$, and no partial
trace is needed, giving the asserted equality.
\end{proof}

\begin{corollary}
\label{rob:product-obstruction}
There is no function $\omega:[0,\infty)\to[0,\infty)$ with
$\omega(0)=0$ and $\lim_{s\downarrow0}\omega(s)=0$, independent of $n$,
such that
\[
 \operatorname{dist}_\diamond
       (\mathcal L,\operatorname{LocGKLS}_n)
 \leq\omega(\delta_{\mathrm{prod}}(\mathcal L))
\]
for every Hamiltonian generator $\mathcal L$ of diamond norm one.
This remains false if the Hamiltonian is required to be a sum of
nearest-neighbour two-site terms on a path.
\end{corollary}

\begin{proof}
Proposition~\ref{rob:path} gives a sequence with distance one and
$0\leq\delta_{\mathrm{prod}}(\mathcal L_m)\leq\sqrt{3/m}\to0$,
contradicting the required behaviour of $\omega$.
Proposition~\ref{rob:rankone} separately gives exponentially small
product defect and distance tending to one.
\end{proof}

The normalization of the path interactions is by their total strength;
no positive lower bound on each edge coefficient is assumed.  These
examples concern a global diamond norm and do not obstruct estimation
of individual interaction coefficients in other access models.
Nor do they contradict Theorem~\ref{rob:diagonal}: correlated inputs
can detect what all product inputs detect only weakly.  For instance,
the GHZ state on all $n$ sites belongs to $\mathcal Y_n$ and has
$\Tr(P_n\rho)=1/2$, so the product-population bound used above does not
extend to arbitrary physical Lee--Yang states.


\section*{Disclosure of artificial intelligence use}
Most of the mathematical content, proofs, and text in this article
were generated with substantial assistance from
OpenAI GPT-6 Pro in ChatGPT chat mode and GPT-6 Astra in Codex, both
using the ultra reasoning setting. These were the only AI tools used.
The author organized and directed the research through repeated
interactions, including questions, choices of research directions,
requests for extensions and critical review, and selection and
organization of the material. AI assistance was also used for
literature searches, proof audits, exact computational checks, and
manuscript revision. These checks do not constitute complete
proof-assistant verification, which is not claimed.
The author takes full responsibility for the accuracy, originality,
attribution, and conclusions of the article.


\enlargethispage{2\baselineskip}
\makeatletter\def\@biblabel#1{#1.}\makeatother
\providecommand{\bysame}{\leavevmode\hbox to3em{\hrulefill}\thinspace}
\providecommand{\MR}{\relax\ifhmode\unskip\space\fi MR }
% \MRhref is called by the amsart/book/proc definition of \MR.
\providecommand{\MRhref}[2]{%
  \href{http://www.ams.org/mathscinet-getitem?mr=#1}{#2}
}
\providecommand{\href}[2]{#2}
\begin{thebibliography}{10}

\bibitem{AlNuwairan14}
Muneerah Al~Nuwairan, \emph{The extreme points of {$SU(2)$}-irreducibly
  covariant channels}, Internat. J. Math. \textbf{25} (2014), 1450048.

\bibitem{Aniello10}
P.~Aniello, A.~Kossakowski, G.~Marmo, and F.~Ventriglia, \emph{Brownian motion
  on {Lie} groups and open quantum systems}, J. Phys. A: Math. Theor.
  \textbf{43} (2010), no.~26, 265301.

\bibitem{AsanoJapanese70}
Taro Asano, \emph{Rigorous theorems concerning {Heisenberg} ferromagnets},
  Bussei Kenkyu \textbf{14} (1970), no.~2, 72--96, In Japanese; title
  translated from the original.

\bibitem{ARS24}
Tommaso Aschieri, B{\l}a{\.{z}}ej Ruba, and Jan~Philip Solovej,
  \emph{{SU(2)}-equivariant quantum channels: Semiclassical analysis}, Comm.
  Math. Phys. \textbf{405} (2024), 298.

\bibitem{BB09}
Julius Borcea and Petter Br{\"a}nd{\'e}n, \emph{The {Lee--Yang} and
  {P{\'o}lya--Schur} programs. {I}. {Linear} operators preserving stability},
  Invent. Math. \textbf{177} (2009), no.~3, 541--569.

\bibitem{BB10II}
\bysame, \emph{The {Lee--Yang} and {P{\'o}lya--Schur} programs. {II}. {Theory}
  of stable polynomials and applications}, Comm. Pure Appl. Math. \textbf{62}
  (2009), no.~12, 1595--1631.

\bibitem{BB09Circular}
\bysame, \emph{{P{\'o}lya--Schur} master theorems for circular domains and
  their boundaries}, Ann. of Math. (2) \textbf{170} (2009), no.~1, 465--492.

\bibitem{DenisMartin22}
J\'{e}r\^{o}me Denis and John Martin, \emph{Extreme depolarization for any
  spin}, Phys. Rev. Research \textbf{4} (2022), no.~1, 013178.

\bibitem{FilippovKuzhamuratova19}
Sergey~N. Filippov and Ksenia~V. Kuzhamuratova, \emph{Quantum informational
  properties of the {Landau--Streater} channel}, J. Math. Phys. \textbf{60}
  (2019), no.~4, 042202.

\bibitem{GKS76}
Vittorio Gorini, Andrzej Kossakowski, and E.~C.~G. Sudarshan, \emph{Completely
  positive dynamical semigroups of {$N$}-level systems}, J. Math. Phys.
  \textbf{17} (1976), no.~5, 821--825.

\bibitem{Hunt56}
G.~A. Hunt, \emph{Semi-groups of measures on {Lie} groups}, Trans. Amer. Math.
  Soc. \textbf{81} (1956), no.~2, 264--293.

\bibitem{KaltonRoberts83}
Nigel~J. Kalton and James~W. Roberts, \emph{Uniformly exhaustive submeasures
  and nearly additive set functions}, Trans. Amer. Math. Soc. \textbf{278}
  (1983), no.~2, 803--816.

\bibitem{Knese24}
Greg Knese, \emph{Rational inner functions on the polydisk---a survey},
  arXiv:2409.14604, 2024.

\bibitem{KnillLaflamme97}
Emanuel Knill and Raymond Laflamme, \emph{Theory of quantum error-correcting
  codes}, Phys. Rev. A \textbf{55} (1997), no.~2, 900--911.

\bibitem{Lindblad76}
G{\"o}ran Lindblad, \emph{On the generators of quantum dynamical semigroups},
  Comm. Math. Phys. \textbf{48} (1976), 119--130.

\bibitem{Voit03}
Michael Voit, \emph{A {L\'{e}vy}-characterization for {Gaussian} processes on
  matrix groups}, Publ. Math. Debrecen \textbf{63} (2003), no.~1--2, 115--132.

\bibitem{WeiGenerators2026}
Dongsheng Wei, \emph{Generators of stability-preserving semigroups, spectral
  gaps, and classical ground-state computation}, 2026, Unpublished companion
  manuscript.

\bibitem{Werner98}
Reinhard~F. Werner, \emph{Optimal cloning of pure states}, Phys. Rev. A
  \textbf{58} (1998), no.~3, 1827--1832.

\bibitem{WBG26}
Benjamin Wong, Sergey Bravyi, David Gosset, and Yinchen Liu, \emph{{Lee--Yang}
  tensors and {Hamiltonian} complexity}, arXiv:2602.03605, 2026, Version 1,
  February 3, 2026.

\end{thebibliography}

\end{document}

